\chapter{Modules and commutative algebra}\label{ch:modules}

The paper's arithmetic is carried by modules: Selmer groups and Tamagawa factors are finite $\Zl$-modules measured by length and by
Fitting ideals; the Hecke algebra $\T$ is a finite free $\Z$-algebra whose completion $\T\otimes\Zl$ splits into local factors $\T_{\mf m}$; the
cohomology $H^1(X_0(N),\Zl)$ is a lattice with a perfect alternating pairing for which $\T$ is self-adjoint; and the inequality of Lemma lem:ddt417
compares a Gram determinant with a congruence ideal. This chapter develops, with complete proofs except where a reference is given, the commutative
algebra behind all of this: exact sequences and the snake lemma, $\Hom$ and $\otimes$, Nakayama's lemma and localisation, modules over principal ideal
domains and saturation, discrete valuation rings and $\Zl$, length (with the Jordan-Holder theorem) and Fitting ideals, Pontryagin duality, pairings and
Gram determinants, finite free algebras over a discrete valuation ring with their idempotents and generalised eigenspaces, congruence ideals and
Gorenstein algebras. It uses the ring theory of \Cref{ch:rings} and the finite abelian groups of \Cref{ch:groups}, and does not reprove them.

\section*{Core results at a glance}

Rings are commutative with $1$ (except in item \ref{card:modules:jh}); $\mathcal O$ is a discrete valuation ring with uniformiser $\varpi$, residue field $k$,
fraction field $K$ and valuation $v$; $M^\vee=\Hom(M,R)$ is the linear dual and $M^*$ the Pontryagin dual; ``length'' is $\len_{\Zl}$; $\T_{\mf m}$ is a factor
of $\T\otimes\Zl$.

\begin{enumerate}[label=\textbf{\thechapter.\arabic*}]
\item \textbf{Conventions.} Appendix B (B-C4, B-C10, B-C13, B-C15); rings commutative; $\mathcal O$, $\varpi$, $k$, $K$, $v$ with $v(\varpi)=1$; $M[\mf a]$; lattices;
pairings, left and right orthogonals, self-adjointness; lengths in place of the order $\#\mathcal O/\Delta$; relations are rows; the letters $I$, $J$ of
Lemma lem:crit; $\T_{\mf m}=e_{\mf m}(\T\otimes\Zl)$, never the algebraic localisation $\T_{(\mf m)}$. \Cref{conv:modules:rings,conv:modules:tm}.
\item \textbf{Modules} (definition: module, submodule, quotient, finitely generated, cyclic, homomorphism, kernel, image, cokernel, $\Hom_R$, $\End_R$;
action of an algebra, faithful action). \Cref{def:modules:module,def:modules:algebra-modules}.
\item \textbf{Basic facts.} The set $\Hom_R(M,N)$ is an $R$-module, the ring $\End_R(M)$ is an $R$-algebra, and $Rm\simeq R/\Ann_R(m)$ for $m\in M$. \emph{(proved here)} \Cref{prop:modules:module-basics}.
\item \textbf{Isomorphism theorems.} $M/\ker f\simeq f(M)$; $(A+B)/B\simeq A/(A\cap B)$; $(M/B)/(A/B)\simeq M/A$ with the correspondence of submodules; for $f$
surjective with kernel $K_0$, $M/(K_0+A)\simeq N/f(A)$. \emph{(proved here)} \Cref{thm:modules:isomorphism}.
\item \textbf{Exact sequences; splitting lemma} (definition: exactness, short exact sequences). For $0\to M'\to M\to M''\to0$: a section of $M\to M''$ exists iff a retraction of $M'\to M$ exists iff $M\simeq M'\oplus M''$ compatibly.
\emph{(proved here)} \Cref{def:modules:exact}, \Cref{lem:modules:splitting}.
\item \textbf{Two submodules.} $0\to M/(A\cap B)\to M/A\oplus M/B\to M/(A+B)\to0$ is exact; $A+B$ is direct iff $A\cap B=0$. \emph{(proved here)} \Cref{prop:modules:two-submodules}.
\item \textbf{Snake lemma.} A commutative diagram with exact rows $M'\to M\to M''\to0$ and $0\to N'\to N\to N''$ gives an exact sequence
$\ker f'\to\ker f\to\ker f''\xrightarrow{\delta}\coker f'\to\coker f\to\coker f''$, with $0\to$ at the left if $M'\to M$ is injective and $\to0$ at the right if $N\to N''$ is surjective.
\emph{(proved here)} \Cref{lem:modules:snake}.
\item \textbf{Multiplication by $n$.} For $0\to A\to B\to C\to0$: $0\to A[n]\to B[n]\to C[n]\to A/n\to B/n\to C/n\to0$ is exact; $\#\ker\phi=\#\coker\phi$ for an endomorphism of a
finite abelian group. \emph{(proved here)} \Cref{cor:modules:snake}.
\item \textbf{Five lemma.} In a commutative ladder of exact rows of five terms, if the four outer maps are isomorphisms so is the middle one (with the two
one-sided halves). \emph{(proved here)} \Cref{lem:modules:five}.
\item \textbf{Sums and products.} $\Hom$ turns direct sums in the first variable, and products in the second, into products; $M=A\oplus B$ iff $A+B=M$ and $A\cap B=0$ iff $A=eM$,
$B=(1-e)M$ for an idempotent $e\in\End M$. \emph{(proved here)} \Cref{def:modules:sums-products}, \Cref{prop:modules:sums}.
\item \textbf{Free modules.} Universal property; every module has a presentation $F_1\to F_0\to M\to0$; $R^m\simeq R^n$ implies $m=n$ ($R\ne0$); $A\in\Mat_n(R)$ is invertible iff
$\det A\in R^\times$. \emph{(proved here)} \Cref{def:modules:free}, \Cref{prop:modules:free}.
\item \textbf{Projective modules.} Equivalent: every surjection onto $P$ splits; $P$ is a summand of a free module; $\Hom(P,-)$ is exact; maps from $P$ lift along
surjections. $eR$ is projective for an idempotent $e$; $\Z/2$ over $\Z/6$ is projective, not free. \emph{(proved here)} \Cref{def:modules:projective},
\Cref{prop:modules:projective}.
\item \textbf{Hom.} $\Hom(R/\mf a,M)=M[\mf a]$; $\Hom(N,-)$ and $\Hom(-,N)$ are left exact, not right exact; $M\simeq M^{\vee\vee}$ for $M$ free of finite rank; over a
domain $M^\vee$ is torsion-free. \emph{(proved here)} \Cref{prop:modules:hom}.
\item \textbf{Tensor product.} Existence and universal property; $(R/\mf a)\otimes M=M/\mf aM$; distributivity; adjunction $\Hom(M\otimes N,P)\simeq\Hom(M,\Hom(N,P))$; base change
$S\otimes_RM$. Examples: $\Z/m\otimes\Z/n=\Z/\gcd(m,n)$, $\Zl\otimes\Z/n=\Z/\ell^{v_\ell(n)}$, $\Qlmod\otimes A=0$ for $A$ torsion. \emph{(proved here)} \Cref{thm:modules:tensor},
\Cref{ex:modules:tensor-table}.
\item \textbf{Right exactness.} $-\otimes N$ is right exact and not left exact; $N$ is flat if $-\otimes N$ is exact; projective modules are flat. \emph{(proved here)}
\Cref{prop:modules:right-exact}.
\item \textbf{Flatness.} $S^{-1}R$ is flat; over a principal ideal domain flat equals torsion-free (so $\Zl$, $\Q$ are flat over $\Z$); for $R\to S$ flat, $S\otimes-$ commutes with
kernels, images, cokernels, finite intersections, $M[\mf a]$ ($\mf a$ finitely generated) and $\Ann$ of finitely generated modules; saturation commutes with flat
extensions of principal ideal domains. \emph{(proved here)} \Cref{prop:modules:flat}.
\item \textbf{Noetherian modules.} Ascending chain condition, maximal condition, every submodule finitely generated: equivalent; closed under submodules,
quotients, extensions; finitely generated modules over Noetherian rings are finitely presented; in a nonzero Noetherian ring every proper ideal lies in a
maximal ideal; rings finite over Noetherian rings are Noetherian. Hilbert's basis theorem. \emph{(proved here, except Hilbert's theorem: stated; reference)}
\Cref{prop:modules:noetherian}.
\item \textbf{Determinant trick.} If $M$ has $n$ generators and $\phi(M)\subseteq\mf aM$, then $\phi^n+a_1\phi^{n-1}+\dots+a_n=0$ with $a_i\in\mf a^i$; if $\mf aM=M$ some $x\equiv1\bmod\mf a$
kills $M$; surjective endomorphisms of finitely generated modules are bijective. \emph{(proved here)} \Cref{lem:modules:determinant-trick}.
\item \textbf{Local rings.} $R$ is local iff its non-units form an ideal $\mf m$; $1+x$ is a unit for $x$ in the Jacobson radical; examples $\Z/p^n$, $k[x]/(x^2)$, $\T_{\mf m}$;
non-examples $\Zl\times\Zl$, $\Z_{(11)}[x]/(x^2+x-1)$. \emph{(proved here)} \Cref{def:modules:local-ring}, \Cref{prop:modules:local-ring}.
\item \textbf{Nakayama's lemma.} Over a local ring $(R,\mf m,k)$ and $M$ finitely generated: $\mf mM=M\Rightarrow M=0$ (also for ideals in the Jacobson radical); $N+\mf mM=M\Rightarrow N=M$;
$x_i$ generate $M$ iff they span $M/\mf mM$, and the minimal number of generators is $\dim_kM/\mf mM$; surjectivity and invertibility can be tested modulo $\mf m$;
finitely generated projective modules are free. \emph{(proved here)} \Cref{thm:modules:nakayama}.
\item \textbf{Localisation.} $S^{-1}R$ has the universal property for inverting $S$; $S^{-1}M=S^{-1}R\otimes M$ and localisation is exact; primes of $S^{-1}R$ are the primes
of $R$ missing $S$; $R_{\mf p}$ is local; $M=0$ iff all $M_{\mf m}=0$; localisation commutes with sums, quotients and $\Ann$ of finitely generated modules (and with $\Hom$
out of finitely presented modules: stated); $(\Z/n)_{(p)}=\Z/p^{v_p(n)}$; localising $R_1\times R_2$ at $\mf p_1\times R_2$ gives $(R_1)_{\mf p_1}$. \emph{(proved here, except the
statement on $\Hom$: stated; reference)} \Cref{thm:modules:localisation}.
\item \textbf{Submodules of free modules.} Over a principal ideal domain a submodule of $R^n$ is free of rank $\le n$; finitely generated projective modules are free.
\emph{(proved here)} \Cref{thm:modules:submodule-free}.
\item \textbf{Smith normal form.} Over a principal ideal domain, $UAV=\operatorname{diag}(d_1,\dots,d_s,0,\dots)$ with $d_1\mid\cdots\mid d_s$, $U$, $V$ invertible; $(d_1\cdots d_k)$ is the ideal of
$k\times k$ minors, so the $(d_i)$ are unique. \emph{(proved here)} \Cref{thm:modules:smith}.
\item \textbf{Structure theorem.} Over a principal ideal domain, $M\simeq R^r\oplus\bigoplus R/(d_i)$ ($d_1\mid\cdots\mid d_s$ non-units) $\simeq R^r\oplus\bigoplus R/(p_j^{a_j})$, with $r$, the $(d_i)$ and the
multiset of $(p_j^{a_j})$ unique. \emph{(proved here)} \Cref{thm:modules:structure-pid}.
\item \textbf{Torsion.} $M_{\tors}=\ker(M\to M\otimes K)$; over a principal ideal domain finitely generated torsion-free modules are free and $M\simeq M_{\tors}\oplus M/M_{\tors}$; if
$R\subseteq R'$ with $R'$ a principal ideal domain free of rank $e$ over $R$, a finitely generated $R'$-module that is torsion-free of $R$-rank $n$ is $R'$-free of rank $n/e$
(so $T_3E$ is free of rank one over $\Z_3[\zeta_3]$ in Lemma lem:IV). \emph{(proved here)} \Cref{def:modules:torsion}, \Cref{prop:modules:torsion}.
\item \textbf{Saturation.} For $L$ in a lattice $M$ over a principal ideal domain: $M/L$ torsion-free iff $L=M\cap KL$ iff $L$ is a direct summand iff (over $\mathcal O$)
$L\cap\varpi M=\varpi L$; $\operatorname{sat}(L)=M\cap KL$ is the least saturated submodule containing $L$; kernels to torsion-free modules, $M\cap W$, $M[\mf a]$ and orthogonals are
saturated. \emph{(proved here)} \Cref{def:modules:saturated}, \Cref{prop:modules:saturated}.
\item \textbf{Stacked bases.} For $L\subseteq M\simeq R^n$ of rank $d$ there is a basis $e_i$ of $M$ with $d_1e_1,\dots,d_de_d$ a basis of $L$, $d_1\mid\cdots\mid d_d$ unique; $L$ is saturated iff
all $d_i$ are units; for $d=n$, $[M:L]=\abs{\det B}$ over $\Z$ and $\len(M/L)=v(\det B)$ over $\mathcal O$. \emph{(proved here)} \Cref{thm:modules:stacked-bases}.
\item \textbf{Dual of an inclusion.} For lattices $L\subseteq M$ of the same rank, $M^\vee\into L^\vee$ and $L^\vee/M^\vee\simeq\Hom(M/L,K/R)$, which is $(M/L)^*$ over $\Z$ or $\Zl$; so
$[L^\vee:M^\vee]=[M:L]$. \emph{(proved here)} \Cref{prop:modules:dual-inclusion}.
\item \textbf{Congruence module of a decomposition.} If $M\otimes K=V_1\oplus V_2$ and $L_i=M\cap V_i$, then $M/(L_1\oplus L_2)\simeq p_1(M)/L_1\simeq p_2(M)/L_2$; if $V_1=Kf$ with $f$
primitive, the quotient is cyclic $\simeq R/(r)$, $rh=f+g_0$ with $p_1(h)=f/r$, $g_0\in L_2$, and $f\equiv g\pmod{r'M}$ for some $g\in L_2$ iff $r'\mid r$. \emph{(proved here)}
\Cref{prop:modules:congruence-decomposition}.
\item \textbf{Discrete valuation rings.} For a Noetherian local domain $\mathcal O$, not a field: $\mf m$ principal iff $\mathcal O$ is a principal ideal domain iff the nonzero ideals
are the $\mf m^k$ iff $\mathcal O$ is the ring of a discrete valuation iff $\mathcal O$ is integrally closed of dimension one; then $x=u\varpi^{v(x)}$ uniquely.
$\mathcal O$ is \emph{complete} if $\mathcal O\to\varprojlim\mathcal O/\varpi^n$ is bijective. \emph{(proved here,
except the equivalence with integral closure: stated; reference)} \Cref{thm:modules:dvr}.
\item \textbf{Modules over a discrete valuation ring.} $M\simeq\mathcal O^r\oplus\bigoplus_{i\le s}\mathcal O/\varpi^{a_i}$ uniquely; torsion-free implies free; $\dim_kM/\varpi M=r+s$, $\dim_kM[\varpi]=s$;
formulas for $M[\varpi^k]$, $M/\varpi^kM$; cyclic iff $r+s\le1$; $L$ saturated iff $L\cap\varpi M=\varpi L$. \emph{(proved here)} \Cref{prop:modules:dvr-modules}.
\item \textbf{The ring $\Zl$.} $\Zl=\varprojlim\Z/\ell^n$ is a discrete valuation ring with uniformiser $\ell$ and residue field $\Fl$; $\Zl/\ell^n=\Z/\ell^n$; integers prime to $\ell$ are units,
so a $\Zl$-module killed by such an integer is $0$; $\Ql=\Zl[1/\ell]$ and $\Q\cap\Zl=\Z_{(\ell)}$; $\Zl$ is flat over $\Z$ and $A\otimes\Zl=\Zl^r\oplus A_{\tors}[\ell^\infty]$; finitely generated $\Zl$-modules
are complete; completions of Noetherian rings are flat. \emph{(proved here, except the statement on completions: stated; reference)} \Cref{thm:modules:zl}.
\item \textbf{Baer's criterion.} Over a principal ideal domain divisible modules are injective, and $\Hom(-,D)$ is exact. \emph{(proved here)} \Cref{lem:modules:baer}.
\item \textbf{The module $\Qlmod$.} $\Qlmod=\Z[1/\ell]/\Z=(\Q/\Z)[\ell^\infty]=\varinjlim\Z/\ell^k$; divisible, with $(\Qlmod)[\ell^k]\simeq\Z/\ell^k$ its only proper submodules;
$\End_{\Zl}(\Qlmod)=\Zl$; $\Qlmod$ and $\Q/\Z$ are injective. \emph{(proved here)} \Cref{prop:modules:qlzl}.
\item \textbf{Topology.} Finitely generated $\Zl$-modules are compact, Hausdorff and complete; linear maps are continuous, submodules and images closed; $\GL_n(\Zl)$ is
profinite; a compact $\Zl$-module $A$ with $A/\ell A$ finite is finitely generated. \emph{(proved here, except topological Nakayama: proof sketched here)}
\Cref{prop:modules:zl-topology}.
\item \label{card:modules:jh}\textbf{Jordan-Holder and length.} Over any ring, any two composition series have the same length and factors; submodules and quotients of
finite-length modules have finite length; length is additive in exact sequences; finite length iff Noetherian and Artinian (stated). Over a local ring simple
modules are $\simeq k$, $\len M=\sum\dim_k\mf m^iM/\mf m^{i+1}M$ when $\mf m^NM=0$, and finitely generated modules over a Noetherian local ring have finite length iff killed by a
power of $\mf m$; $\len_R=f\len_{R'}$ for a finite local extension with residue degree $f$. \emph{(proved here, except ``Noetherian and Artinian'': stated; reference)}
\Cref{def:modules:composition}, \Cref{thm:modules:jordan-holder}.
\item \textbf{Length over a discrete valuation ring.} $\len\mathcal O/\varpi^a=a$; finite length iff torsion; $\#M=(\#k)^{\len M}$ for finite $k$, so $\len_{\Zl}M=v_\ell(\#M)$;
$\len M[\varpi^k]=\len M/\varpi^kM$ for finite length. \emph{(proved here)} \Cref{prop:modules:length-dvr}.
\item \textbf{Annihilators.} $\Ann(R/\mf a)=\mf a$; $\Ann(M\oplus N)=\Ann M\cap\Ann N$; $M[\mf a]=\Hom(R/\mf a,M)$; $\T[\mf a]=\Ann_\T(\mf a)$ and $L[\mf a]=\Ann_\T(\mf a)^r$ for $L\simeq\T^r$; flat base
change; $\Ann(\bigoplus\mathcal O/\varpi^{a_i})=\varpi^{\max a_i}$. \emph{(proved here)} \Cref{def:modules:annihilator}, \Cref{prop:modules:annihilator}.
\item \textbf{Fitting ideals} (definition: for generators $x_1,\dots,x_n$, the ideal generated by the determinants of $n$-tuples of relations, equivalently by the $n\times n$
minors of a presentation matrix; it does not depend on the choices). \emph{(proved here)} \Cref{def:modules:fitting}, \Cref{lem:modules:fitting-well-defined}.
\item \textbf{Properties of Fitting ideals.} $\Fitt(R/\mf a)=\mf a$ and $\Fitt(R^n)=0$; the Fitting ideal of a direct sum is the product, $\Fitt(M\oplus N)=\Fitt(M)\Fitt(N)$; $\Fitt$ grows along surjections; for $M$ with $n$
generators, $\Ann(M)^n\subseteq\Fitt(M)\subseteq\Ann(M)$;
$\Fitt$ commutes with every base change; for a short exact sequence with $M'$ finitely generated, $\Fitt(M')\Fitt(M'')\subseteq\Fitt(M)$, with strict inclusion for the module $\Z[x]/(4,x)$; over a domain,
$\Fitt(M)\ne0$ if and only if $M$ is torsion.
\emph{(proved here)} \Cref{prop:modules:fitting-properties}.
\item \textbf{Fitting equals $\varpi^{\mathrm{length}}$.} Over a principal ideal domain $\Fitt(\bigoplus R/(d_i))=(\prod d_i)$ (and $0$ with a free summand); over $\mathcal O$,
$\Fitt(M)=\varpi^{\len M}\mathcal O$ for torsion $M$, so $\Fitt_{\Zl}(M)=\#M\cdot\Zl$, multiplicative in exact sequences; $\Fitt_{\Zl}(\Z/\ell^v\oplus\Z/\ell^m)=\ell^{v+m}\Zl$, while
$\Ann=\ell^{\max(v,m)}\Zl$. \emph{(proved here)} \Cref{thm:modules:fitting-dvr}, \Cref{ex:modules:fitting-examples}.
\item \textbf{Pontryagin duality for finite groups.} $A^*=\Hom(A,\Q/\Z)\simeq A$; $A\simeq A^{**}$; duality is exact; $B^\perp\simeq(A/B)^*$, $B^*\simeq A^*/B^\perp$; for $\ell$-primary $A$,
$A^*=\Hom(A,\Qlmod)$ and $\len A^*=\len A$. \emph{(proved here)} \Cref{thm:modules:pontryagin-finite}.
\item \textbf{Pontryagin duality for $\Zl$-modules.} $X^*=\Hom(X,\Qlmod)$ is an exact duality between finitely generated $\Zl$-modules and cofinitely generated $\ell$-primary
modules, with $\Zl^*=\Qlmod$, $(\Qlmod)^*=\Zl$, $X\simeq X^{**}$, and $W^*=\varprojlim W[\ell^k]^*$ profinite; every cofinitely generated $W$ is $(\Qlmod)^r\oplus F$ with $F$ finite and $r=\rank W^*$; lengths are preserved.
\emph{(proved here; the extension to all compact and discrete modules: stated; reference)} \Cref{thm:modules:pontryagin-zl}.
\item \textbf{$T$ and $W=T\otimes\Qlmod$.} For $T$ free of rank $r$: $0\to T\to V\to W\to0$ is exact; $W\simeq(\Qlmod)^r$; $W[\ell^k]\simeq T/\ell^kT$; $T=\varprojlim W[\ell^k]$; $W^*\simeq T^\vee$;
all functorially; the same over any complete $\mathcal O$. \emph{(proved here)} \Cref{prop:modules:t-and-w}.
\item \textbf{Cyclic $\ell$-primary modules.} An $\ell$-primary torsion module with $\#A[\ell]\le\ell$ is $0$, $\Z/\ell^h$ or $\Qlmod$; if finite its length is the largest $k$ with
an element of exact order $\ell^k$. \emph{(proved here)} \Cref{cor:modules:cyclic-l-primary}.
\item \textbf{Perfect pairings and Gram determinants.} A pairing between free modules of rank $n$ over any ring is perfect iff its Gram determinant is a unit (over a domain, nondegenerate iff it is nonzero); $\det(P^{\mathsf T}GQ)=\det P\det Q\det G$,
so $\Delta R$, and over $\mathcal O$ the valuation $v(\Delta)$, are basis independent; perfectness survives base change. \emph{(proved here)} \Cref{def:modules:pairing},
\Cref{prop:modules:gram}.
\item \textbf{Orthogonals.} For a perfect pairing on a lattice of rank $n$, orthogonals are saturated, $N^\perp$ has rank $n-\rank N$, and ${}^\perp(N^\perp)=\operatorname{sat}(N)$.
\emph{(proved here)} \Cref{prop:modules:orthogonal}.
\item \textbf{Gram determinant of a saturated sublattice.} For a perfect pairing on $L$ over $\mathcal O$ and a \emph{saturated} $L'\subseteq L$ with Gram determinant $\Delta\ne0$:
$L/L'^\perp\simeq L'^\vee$, $L'\cap L'^\perp=0$, and $L/(L'\oplus L'^\perp)\simeq\coker(L'\to L'^\vee)$ has length $v(\Delta)$ (order $\#\mathcal O/\Delta$ for finite $k$); $L'^\perp$ has the same
Gram valuation; saturation is necessary. \emph{(proved here)} \Cref{thm:modules:gram-cokernel}.
\item \textbf{Alternating pairings.} Alternating implies antisymmetric (and conversely if $2$ is a nonzerodivisor); odd-rank alternating Gram determinants vanish; in
rank two $\Delta=a^2$ with $a\mathcal O$ the ideal of values; $\det=\operatorname{Pf}^2$ (stated, checked in rank $4$); symplectic normal form. \emph{(proved here, except the Pfaffian
identity: stated; reference; and the normal form: proof sketched here)} \Cref{prop:modules:alternating}.
\item \textbf{Adjoints.} For a perfect pairing $A^\dagger=G^{-1}A^{\mathsf T}G$; for a commutative algebra $\T$ of self-adjoint operators $(\mf aL)^\perp=L[\mf a]$, idempotents of $\T$ (or
$\T\otimes S$) give orthogonal decompositions on whose summands the pairing stays perfect; if $A^\dagger=wAw^{-1}$ on $\T$ then $\langle x,wy\rangle$ is perfect with $\T$ self-adjoint,
and alternating if $\langle\ ,\ \rangle$ is antisymmetric and $w^\dagger=w$. \emph{(proved here)} \Cref{prop:modules:adjoint}.
\item \textbf{An algebra paired with a module.} If $S\times\T\to R$ is perfect with $\langle ts,u\rangle=\langle s,tu\rangle$, then $S\simeq\T^\vee$ as $\T$-modules, $S[\mf a]\simeq\Hom(\T/\mf a,R)$
is saturated, $\T$ acts faithfully, and for $\mf a\cap\mf b=0$ with $\T/\mf a$, $\T/\mf b$ torsion-free and complementary ranks, $S/(S[\mf a]\oplus S[\mf b])\simeq(\T/(\mf a+\mf b))^*$ over $\Z$ or $\Zl$.
\emph{(proved here)} \Cref{prop:modules:algebra-duality}.
\item \textbf{Finite-dimensional commutative algebras.} Finitely many maximal ideals, all primes maximal, nilpotent Jacobson radical; $A\simeq\prod A_{\mf m_i}$ via idempotents;
reduced iff a product of fields; $V=\bigoplus e_iV$ with $e_iV=V_{\mf m_i}$. \emph{(proved here)} \Cref{thm:modules:artinian-algebras}.
\item \textbf{Finite free algebras over $\mathcal O$.} Noetherian; maximal ideals contain $\varpi$ and correspond to those of $\T/\varpi\T$; $\T\subseteq\T_K$; $\T$ reduced iff $\T_K$ is; for
$\pi\colon\T\onto\mathcal O$, $\T=\mathcal O\oplus\wp$, $\mf m_\pi=\wp+\varpi\T$, and $\pi_K$ factors through exactly one local factor of $\T_K$. \emph{(proved here)}
\Cref{thm:modules:finite-algebra-dvr}.
\item \textbf{Idempotent lifting.} Over a complete $\mathcal O$, idempotents of $\T/\varpi\T$ lift uniquely (Newton iteration $e\mapsto3e^2-2e^3$); $\T=\prod_{\mf m}\T_{\mf m}$ with $\T_{\mf m}=e_{\mf m}\T$
local and finite free; $L=\bigoplus e_{\mf m}L$; a map to a local ring factors through one $\T_{\mf m}$; completeness is necessary ($\Z_{(11)}[x]/(x^2+x-1)$). \emph{(proved here)}
\Cref{thm:modules:idempotent-lifting}.
\item \textbf{Hecke-type algebras.} If $\T$ is free of finite rank over $\Z$, then $\T\otimes\Zl=\prod_{\mf m\ni\ell}\T_{\mf m}$ over the maximal ideals of $\T$ containing $\ell$. Example: the
level-$23$ algebra $\Z[x]/(x^2+x-1)$ splits over $\Z_{11}$, and is local over $\Z_2$, $\Z_3$, $\Z_5$. \emph{(proved here)} \Cref{cor:modules:hecke-type}, \Cref{ex:modules:level23}.
\item \textbf{Generalised eigenspaces.} For a commutative $\T\subseteq\End_F(V)$ with split characteristic polynomials, $V=\bigoplus_\phi V(\phi)$ over the systems of eigenvalues, with
projections the idempotents of $\T$, a common eigenvector in each $V(\phi)$, and $W=\bigoplus(W\cap V(\phi))$ for stable $W$; integrally, over a complete $\mathcal O$, systems
of eigenvalues are $\mathcal O$-valued, $\mf m_\phi=\mf m_\psi$ iff $\phi\equiv\psi\bmod\varpi$, and $e_{\mf m}L\otimes K=\bigoplus_{\mf m_\phi=\mf m}V(\phi)$. \emph{(proved here)}
\Cref{thm:modules:generalised-eigenspaces}.
\item \textbf{Projecting a congruence.} $e_{\mf m}$ preserves $L$ and $\T$-stable subspaces and fixes $V(\phi)$ for $\mf m_\phi=\mf m$; so $x\equiv y\pmod{\varpi^tL}$ with $x\in V(\phi)$ gives
$x\equiv e_{\mf m}y$. \emph{(proved here)} \Cref{prop:modules:idempotent-congruence}.
\item \textbf{Semisimple operators.} $A$ is semisimple iff $\mu_A$ is squarefree iff $F[A]$ is reduced (iff diagonalisable over $\overline F$ for perfect $F$); a cyclic vector makes
$\mu_A$ the characteristic polynomial; a commutative $\T$ is reduced iff all its elements are semisimple; Jordan-Chevalley (stated). The shift operator ($U_p$ on
oldforms) has characteristic polynomial $X^{e-1}(X^2-aX+p)$, is not semisimple for $e\ge3$, yet every eigenvalue has an eigenvector. \emph{(proved here, except
Jordan-Chevalley: stated; reference)} \Cref{prop:modules:semisimple-operators}, \Cref{ex:modules:shift}.
\item \textbf{Common eigenvectors.} Over an algebraically closed field, a nonzero joint generalised eigenspace of commuting operators contains a common eigenvector of
every commutative algebra containing them; for operators acting on the factors of a tensor product, $V(\lambda)=\bigotimes V_i(\lambda_i)$. \emph{(proved here)}
\Cref{lem:modules:common-eigenvector}.
\item \textbf{Congruence ideal} (definition: for $\pi\colon\T\onto\mathcal O$, $\wp=\ker\pi$, $I=\Ann_\T(\wp)$, $\eta=\pi(I)$; congruence module $\T/(\wp+I)\simeq\mathcal O/\eta$).
\Cref{def:modules:congruence-ideal}.
\item \textbf{Congruence ideals.} $I\ne0$; $\eta\ne0$ iff $\wp\cap I=0$ iff the factor of $\pi$ in $\T_K$ is $K$ (proof: $\pi(t)x=tx+(\pi(t)-t)x=0$); then
$I\simeq\eta$, $\Ann_\T(I)=\wp$, $\eta$ is the denominator of $e_\pi$; over complete $\mathcal O$, $\Ann_\T(\wp)=\Ann_{\T_{\mf m}}(\wp_{\mf m})\times0$, so $\eta$ is computed in $\T_{\mf m}$ (the other
factors lie in $\wp$); $\eta$ commutes with flat base change. \emph{(proved here)} \Cref{prop:modules:congruence-ideal}, \Cref{ex:modules:toy-congruence}.
\item \textbf{Ingredients of Lemma lem:ddt417.} For $L$ free over $\mathcal O$ with a perfect pairing making $\T$ self-adjoint and $\eta\ne0$: $L[\wp]$, $L[I]$ are saturated, $L[\wp]\cap L[I]=0$,
$L[\wp]^\perp=L[I]$, $Q=L/(L[\wp]\oplus L[I])$ is killed by $\eta$, and $Q\simeq(\mathcal O/\eta)^r$ if $L\simeq\T^r$. \emph{(proved here)} \Cref{prop:modules:ddt-ingredients},
\Cref{rem:modules:ddt417-preview}.
\item \textbf{Gorenstein algebras} (definition: $\Hom_{\mathcal O}(\T,\mathcal O)$ free of rank one over $\T$). Equivalent to a perfect pairing $\lambda(st)$; over complete $\mathcal O$ local by local; for
local $\T$, iff the socle of $\T/\varpi\T$ is a line; $\mathcal O[x]/(f)$ and $\T_r$ are Gorenstein, $\T^\flat_3$ is not; complete intersections are Gorenstein (stated). \emph{(proved here,
except the last statement: stated; reference)} \Cref{def:modules:gorenstein}, \Cref{prop:modules:gorenstein}.
\item \textbf{Mazur-style freeness.} For $\T$ local finite free over $\mathcal O$ and $L$ free over $\mathcal O$ with a perfect pairing making $\T$ self-adjoint: if $\dim L/\mf mL=d$ and
$\rank L=d\rank\T$, then $L\simeq\T^d$ and $\T$ is Gorenstein. \emph{(proved here)} \Cref{thm:modules:mazur-freeness}.
\end{enumerate}

\section*{Where the paper uses this chapter}

\begin{itemize}
\item \textbf{Section 1.2} (the Tamagawa factor $\Tam_p(M)=\Fitt_{\Zl}\bigl((H^1(I_p,M)_{\tors})^{G_{\Q_p}}\bigr)$): Fitting ideals and $\Fitt=\ell^{\mathrm{length}}$
(\Cref{def:modules:fitting,thm:modules:fitting-dvr}).
\item \textbf{Lemma lem:tamfactor} (the Tamagawa factor at a multiplicative prime): the structure $H^1(I,M)=(\Zl/n')^2\oplus\Zl$ and its torsion
(\Cref{thm:modules:structure-pid}); the Fitting ideal of $\Z/\ell^v\oplus\Z/\ell^m$ (\Cref{thm:modules:fitting-dvr}, \Cref{pl:modules:tamagawa}).
\item \textbf{Lemma lem:fpr} ($\Fitt_{\Zl}H^1_f(\Q_p,T^D)=L_p(A_f,1)^{-1}\Tam_p(T^D)$): Fitting ideals as lengths (\Cref{thm:modules:fitting-dvr}); $\Q\cap\Zl=\Z_{(\ell)}$
(\Cref{thm:modules:zl}(iv)).
\item \textbf{Theorem thm:main, Theorem thm:dfg37, Remark rem:ns, Conjecture conj:euler} (lengths of Selmer groups and of $H^0(\Q,W^D)$): additivity of length (\Cref{thm:modules:jordan-holder}), length over $\Zl$
(\Cref{prop:modules:length-dvr}), Pontryagin duality preserves length (\Cref{thm:modules:pontryagin-zl}, \Cref{pl:modules:dfg21}).
\item \textbf{Lemma lem:dfg21} (the four-term sequence with Pontryagin duals) and the paragraph after it ($H^1(\R,T^D)$ killed by $2$): \Cref{thm:modules:pontryagin-finite,thm:modules:pontryagin-zl},
\Cref{thm:modules:zl}(iii).
\item \textbf{Lemma lem:h0} ($W^D$, its $\ell^k$-torsion, cyclicity, Nakayama): \Cref{prop:modules:t-and-w,cor:modules:cyclic-l-primary}, \Cref{thm:modules:nakayama}.
\item \textbf{Lemma lem:crit} (congruence criterion): $\T/(I+J)\simeq\Z/\phi_f(J)$ (\Cref{thm:modules:isomorphism}); $\T\into\T/I\times\T/J$ (\Cref{prop:modules:two-submodules}); dualising
it (\Cref{prop:modules:dual-inclusion,prop:modules:algebra-duality}); the congruence module of $S_2=\Z f+L$ (\Cref{prop:modules:congruence-decomposition}); $\T\otimes\Q=\Q\times\T'$ and
$I\cap J=0$ (\Cref{thm:modules:artinian-algebras}, \Cref{prop:modules:congruence-ideal}); ``the same over $\Zl$'' (\Cref{prop:modules:flat}). The lemma
itself is proved in full in \Cref{ch:modgalrep}.
\item \textbf{Lemma lem:old} (oldforms): the non-semisimple $U_p$ and its eigenvectors (\Cref{ex:modules:shift}); common eigenforms in a tensor product
(\Cref{lem:modules:common-eigenvector}).
\item \textbf{Lemma lem:IV} and Lemma lem:tamadd ($T_3E$ free of rank one over $\Z_3[\zeta_3]$): \Cref{prop:modules:torsion}(iii).
\item \textbf{Section 4.1} (the pairing $\hat\delta_!$, perfect, alternating, with self-adjoint Hecke operators; Gram determinant $(\eta^\Sigma_{f,\lambda})^2$) and \textbf{Proposition
prop:pullback}: \Cref{prop:modules:adjoint,prop:modules:alternating}.
\item \textbf{Lemma lem:ddt417} (inequality half of DDT Lemma 4.17): every ingredient, \Cref{prop:modules:congruence-ideal,prop:modules:ddt-ingredients,thm:modules:gram-cokernel},
\Cref{prop:modules:saturated}, summarised in \Cref{rem:modules:ddt417-preview}; the lemma itself is proved in \Cref{ch:congruence}.
\item \textbf{Lemma lem:hecke} and the paragraph after it ($\T_{\mf m}$, the congruence ideal computed in $\T_{\mf m}$, $e_{\mf m}$ self-adjoint): \Cref{conv:modules:tm},
\Cref{thm:modules:idempotent-lifting,cor:modules:hecke-type}, \Cref{prop:modules:congruence-ideal}(v), \Cref{prop:modules:adjoint}(iii).
\item \textbf{Remark rem:free} (Gorenstein, freeness): \Cref{prop:modules:gorenstein,thm:modules:mazur-freeness}, \Cref{thm:modules:nakayama}.
\item \textbf{Theorem thm:twist} ($g^\perp$ as a sum of generalised eigenspaces) and \textbf{Theorem thm:quant}, Steps 1 and 2 (idempotents of $\T_\ell$, generalised
eigenspaces, projecting a congruence, the $\Zl$-version of Lemma lem:crit, Lemma lem:ddt417 without freeness): \Cref{thm:modules:generalised-eigenspaces},
\Cref{prop:modules:idempotent-congruence}, \Cref{thm:modules:idempotent-lifting}, \Cref{prop:modules:flat}.
\end{itemize}

\section{Modules, submodules, quotients and homomorphisms}\label{sec:modules:modules}

This section fixes the language. Every object of the paper's commutative algebra is a module over a commutative
ring: the space $S_2(\Gamma_0(N);\Z)$ of integral cusp forms and the cohomology $H^1(X_0(N),\Zl)$ are modules over the
Hecke algebra $\T$; the Tate module $T=\Tl E$ is a module over $\Zl$, and in Lemma lem:IV over $\Z_3[\zeta_3]$; the
lattice $L$ of Lemma lem:ddt417 is a module over an $\mathcal O$-algebra $\T$. We first record the conventions of the chapter.

\begin{convention}[notation of the chapter]\label{conv:modules:rings}
The book-wide conventions of \Cref{ch:notation} apply, in particular B-C4 (one symbol for each dual: the linear dual
$M^\vee=\Hom_R(M,R)$, the Pontryagin dual $M^*$, the Kummer dual $T^D$, which is not used here, and the adjoint $A^\dagger$ of an
operator, never $A^*$; a star after the name of a map is a pullback), B-C10 (``length'' without subscript is $\len_{\Zl}$, and
$\#X=\ell^{\len X}$ for finite $\Zl$-modules), B-C13(a), (c) (the paper's $\T_{\mf m}$ and its two uses of the letter $I$) and
B-C15(a), (h) (uniformisers are $\varpi$; the paper's $\pi\colon\T\to\mathcal O$ and Gram determinant $\Delta$ keep their names). In addition:
\begin{enumerate}[label=(\alph*)]
\item Rings are commutative with $1$, ring homomorphisms preserve $1$ and modules are unital, as in \Cref{ch:rings}, except in the
Jordan-Holder theorem (\Cref{thm:modules:jordan-holder}). An $R$-algebra is a ring $A$ with a ring homomorphism $R\to A$. The
algebras $\T$ acting on modules, like the Hecke algebra, are commutative subrings of the (noncommutative) ring $\End_R(M)$.
\item $\mathcal O$ is a discrete valuation ring (\Cref{thm:modules:dvr}) with uniformiser $\varpi$, residue field $k=\mathcal O/\varpi\mathcal O$,
fraction field $K$ and valuation $v\colon K\to\Z\cup\{+\infty\}$, $v(\varpi)=1$; $v(a\mathcal O)=v(a)$ for a nonzero ideal. For $\mathcal O=\Zl$:
$\varpi=\ell$, $k=\Fl$, $v=v_\ell$. The paper writes $\lambda$ and $v_\lambda$.
\item $M[\mf a]=\{m\in M:am=0\text{ for all }a\in\mf a\}$, $M[x]=M[xR]$ and $M[\ell^\infty]=\bigcup_kM[\ell^k]$; these are the paper's $L[\wp]$ and
$L[I]$. The Pontryagin dual is $M^*=\Hom(M,\Q/\Z)$ for a finite abelian group and $M^*=\Hom_{\Zl}(M,\Qlmod)$ for a $\Zl$-module that is
finitely generated or $\ell$-primary torsion (the paper's ${}^*$ in Lemma lem:dfg21).
\item Over a principal ideal domain $R$ with fraction field $K$, a \emph{lattice} is a finitely generated free $R$-module, and a lattice
in a finite-dimensional $K$-vector space $V$ is a finitely generated $R$-submodule spanning $V$ (free of rank $\dim V$ by
\Cref{thm:modules:submodule-free}).
\item A \emph{pairing} is an $R$-bilinear map $\langle\ ,\ \rangle\colon L\times M\to R$; its Gram matrix in bases $(e_i)$, $(f_j)$ is
$G=(\langle e_i,f_j\rangle)_{i,j}$. For $S'\subseteq L$ and $S\subseteq M$, $S'^{\perp}=\{y\in M:\langle S',y\rangle=0\}$ is the \emph{right} and
${}^\perp S=\{x\in L:\langle x,S\rangle=0\}$ the \emph{left} orthogonal; for a symmetric or antisymmetric pairing on $L=M$ they agree and we
write $S^\perp$. An endomorphism $t$ of $L$ is \emph{self-adjoint} if $\langle tx,y\rangle=\langle x,ty\rangle$ for all $x,y$.
\item The paper's order $\#\mathcal O/\Delta$ is $(\#k)^{v(\Delta)}$ and makes sense only for finite $k$; the book writes the length
$\len_{\mathcal O}(\mathcal O/\Delta\mathcal O)=v(\Delta)$.
\item Relations are rows: a presentation $R^m\xrightarrow{A}R^n\to M\to0$ means $M=R^n/(\text{row space of }A)$. The letter $\varepsilon$
is not used (dual numbers are $k[x]/(x^2)$), and no Frobenius element occurs.
\item Where the text quotes Lemma lem:crit (in the list above and in paperlinks), $I=\ker\phi_f$ and $J=\Ann_\T(f^\perp)$ are that lemma's
letters ($J$ equals $\Ann_\T(I)$, \Cref{prop:modules:algebra-duality} and its paperlink); everywhere else $I=\Ann_\T(\wp)$, as in Lemma lem:ddt417.
\end{enumerate}
\end{convention}

\begin{convention}[the paper's $\T_{\mf m}$]\label{conv:modules:tm}
Let $\T$ be a commutative ring that is free of finite rank as a $\Z$-module, and $\mf m$ a maximal ideal of
$\T_\ell=\T\otimes\Zl$ (equivalently, by \Cref{cor:modules:hecke-type}, a maximal ideal of $\T$ containing $\ell$). Then
$\T_{\mf m}$ \emph{always} means $(\T\otimes\Zl)_{\mf m}$, the direct factor $e_{\mf m}\T_\ell$ of $\T_\ell$ cut out by an idempotent
$e_{\mf m}$ (\Cref{thm:modules:idempotent-lifting}); likewise $L_{\mf m}=e_{\mf m}(L\otimes\Zl)$ for a $\T$-module $L$. The algebraic
localisation of $\T$ itself at a maximal ideal is written $\T_{(\mf m)}$; it is not complete and is not a direct factor of
anything. This is the paper's usage in Lemma lem:hecke, Remark rem:free and Theorem thm:quant.
\end{convention}

\begin{definition}[modules]\label{def:modules:module}
An \emph{$R$-module} is an abelian group $M$ with a map $R\times M\to M$, $(r,m)\mapsto rm$, such that
$r(m+m')=rm+rm'$, $(r+r')m=rm+r'm$, $(rr')m=r(r'm)$ and $1m=m$. A \emph{submodule} is a subgroup $N\subseteq M$ with
$RN\subseteq N$, and the \emph{quotient} $M/N$ is the module of cosets with $r(m+N)=rm+N$. The submodule generated by a
family $(m_i)$ is $\sum_iRm_i$, the set of finite sums $\sum r_im_i$; $M$ is \emph{finitely generated} (f.g.) if
$M=\sum_{i=1}^nRm_i$ for some $m_1,\dots,m_n$, and \emph{cyclic} if $M=Rm$. A \emph{homomorphism} ($R$-linear map)
$f\colon M\to N$ is a group homomorphism with $f(rm)=rf(m)$; its kernel $\ker f$, image $f(M)=\im f$ and cokernel
$\coker f=N/f(M)$ are modules. $\Hom_R(M,N)$ is the set of homomorphisms and $\End_R(M)=\Hom_R(M,M)$.
\end{definition}

\begin{proposition}\label{prop:modules:module-basics}
\begin{enumerate}[label=(\roman*)]
\item $\Hom_R(M,N)$ is an $R$-module under $(f+g)(m)=f(m)+g(m)$ and $(rf)(m)=rf(m)$, and $\End_R(M)$ is a ring under
composition, with an $R$-algebra structure $R\to\End_R(M)$, $r\mapsto(m\mapsto rm)$.
\item For $m\in M$ let $\Ann_R(m)=\{r\in R:rm=0\}$. Then $\Ann_R(m)$ is an ideal and $R/\Ann_R(m)\simeq Rm$, $r\mapsto rm$.
\end{enumerate}
\end{proposition}

\begin{proof}
(i) $rf$ is $R$-linear because $R$ is commutative: $(rf)(sm)=rsf(m)=s(rf)(m)$. The other axioms are routine, and
$r\mapsto(m\mapsto rm)$ is a ring homomorphism into $\End_R(M)$ whose image is central.
(ii) The map $R\to Rm$, $r\mapsto rm$, is $R$-linear and surjective with kernel $\Ann_R(m)$; apply
\Cref{thm:modules:isomorphism}(i) below (whose proof does not use this proposition).
\end{proof}

\begin{example}\label{ex:modules:module-examples}
A $\Z$-module is an abelian group, with $nm=m+\dots+m$; a module over a field $K$ is a $K$-vector space. A module over
the polynomial ring $K[x]$ is a $K$-vector space $V$ together with one endomorphism (the action of $x$); submodules are
the subspaces stable under it. This dictionary underlies \Cref{sec:modules:algebras}. A module over $R/\mf a$ is the same
as an $R$-module killed by $\mf a$. The submodules of $R$ are its ideals.
\end{example}

\begin{definition}[algebras acting on modules]\label{def:modules:algebra-modules}
Let $\phi\colon R\to S$ be a ring homomorphism. Every $S$-module $M$ is an $R$-module by $rm=\phi(r)m$ (\emph{restriction
of scalars}). If $A$ is an $R$-algebra and $M$ an $R$-module, an \emph{action} of $A$ on $M$ is a homomorphism of
$R$-algebras $A\to\End_R(M)$; equivalently, an $A$-module structure on $M$ whose restriction to $R$ is the given one. The
action is \emph{faithful} if $A\to\End_R(M)$ is injective.
\end{definition}

\begin{example}[running examples]\label{ex:modules:running}
\begin{enumerate}[label=(\alph*)]
\item A module over $R_1\times R_2$ is the same as a pair of modules: with the idempotents $e_1=(1,0)$ and $e_2=(0,1)$
(\Cref{ch:rings}), $M=e_1M\oplus e_2M$, where $e_iM$ is an $R_i$-module, and every pair $(M_1,M_2)$ arises from
$M_1\times M_2$. This is how a module over $\T\otimes\Zl=\prod_{\mf m}\T_{\mf m}$ splits (\Cref{thm:modules:idempotent-lifting}).
\item The Hecke algebra $\T\subset\End_\Z S_2(\Gamma_0(N);\Z)$ acts faithfully on the integral cusp forms, by definition
(\Cref{ch:modforms,ch:hecke}); it also acts on $H^1(X_0(N),\Zl)$ (\Cref{ch:moddeg}).
\item Let $T$ be a free $\Z_3$-module of rank $2$ and $\sigma\in\Aut(T)$ with $\sigma^2+\sigma+1=0$ (the generator of the image of
inertia in Lemma lem:IV). Then $T$ is a module over $\Z_3[X]/(X^2+X+1)$, with $X$ acting as $\sigma$; this ring is
$\Z_3[\zeta_3]$, a discrete valuation ring (\Cref{ch:ant}). \Cref{prop:modules:torsion}(iii) shows that $T$ is free of rank
one over it.
\item The toy Hecke algebras of \Cref{ex:modules:toy-congruence}, which are the smallest algebras exhibiting a congruence.
\end{enumerate}
\end{example}

\begin{theorem}[isomorphism theorems]\label{thm:modules:isomorphism}
\begin{enumerate}[label=(\roman*)]
\item For $f\colon M\to N$, $m+\ker f\mapsto f(m)$ is an isomorphism $M/\ker f\simeq f(M)$.
\item For submodules $A,B\subseteq M$, $(A+B)/B\simeq A/(A\cap B)$.
\item For submodules $B\subseteq A\subseteq M$, $(M/B)/(A/B)\simeq M/A$, and $A\mapsto A/B$ is an inclusion-preserving bijection
between submodules of $M$ containing $B$ and submodules of $M/B$.
\item If $f\colon M\to N$ is surjective with kernel $K_0$ and $A\subseteq M$ is a submodule, then $f$ induces an isomorphism
$M/(K_0+A)\simeq N/f(A)$.
\end{enumerate}
\end{theorem}

\begin{proof}
(i) The map is well defined and injective because $f(m)=f(m')$ iff $m-m'\in\ker f$; it is $R$-linear and onto $f(M)$.
(ii) Apply (i) to $A\to(A+B)/B$, $a\mapsto a+B$, which is surjective (every $a+b+B$ equals $a+B$) with kernel $A\cap B$.
(iii) Apply (i) to $M/B\to M/A$, $m+B\mapsto m+A$, surjective with kernel $A/B$. A submodule $X\subseteq M/B$ is the image of
its preimage, which contains $B$, and $A=\{m:m+B\in A/B\}$ for $A\supseteq B$.
(iv) The composite $M\to N\to N/f(A)$ is surjective; its kernel is $f^{-1}(f(A))=A+K_0$. Apply (i).
\end{proof}

\begin{paperlink}\label{pl:modules:iso}
Lemma lem:crit writes ``$\T/(I+J)\simeq\Z/\phi_f(J)$''. This is \Cref{thm:modules:isomorphism}(iv) for the surjection
$\phi_f\colon\T\to\Z$ with kernel $I$ and the submodule $A=J$. In the same way the congruence module of
\Cref{def:modules:congruence-ideal} is $\T/(\wp+I)\simeq\mathcal O/\pi(I)$ for $\pi\colon\T\onto\mathcal O$ with kernel $\wp$.
\end{paperlink}

\section[Exact sequences; free and projective modules]{Exact sequences, the snake lemma, sums, free and projective modules}\label{sec:modules:exact}

Exact sequences are the bookkeeping device of the whole book. The paper's identity (eq:identity) is obtained by
adding lengths along the four-term exact sequence of Lemma lem:dfg21, and every connecting map of \Cref{ch:cohomology} is
built from the snake lemma proved here. The sequence of two submodules (\Cref{prop:modules:two-submodules}) is the exact
form of two sentences of the paper: ``the inclusion $\T\subset\T/I\times\T/J$, whose cokernel is $\T/(I+J)$'' in Lemma
lem:crit, and ``the sum $L[\wp]+L[I]$ is direct'' in Lemma lem:ddt417. Projective modules enter because the local factors
$e\T$ of a finite free $\Zl$-algebra are again free over $\Zl$.

\begin{definition}[exact sequences]\label{def:modules:exact}
A sequence of homomorphisms $\cdots\to M_{i-1}\xrightarrow{f_{i-1}}M_i\xrightarrow{f_i}M_{i+1}\to\cdots$ is \emph{exact at $M_i$} if
$\ker f_i=\im f_{i-1}$, and \emph{exact} if it is exact at every interior term. A \emph{short exact sequence} is an exact
sequence $0\to M'\xrightarrow{i}M\xrightarrow{p}M''\to0$: $i$ is injective, $p$ is surjective and $\im i=\ker p$, so
$M''\simeq M/i(M')$. It \emph{splits} if the equivalent conditions of \Cref{lem:modules:splitting} hold.
\end{definition}

\begin{lemma}[splitting lemma]\label{lem:modules:splitting}
For a short exact sequence $0\to M'\xrightarrow{i}M\xrightarrow{p}M''\to0$ the following are equivalent:
(a) there is $s\colon M''\to M$ with $ps=\id$; (b) there is $r\colon M\to M'$ with $ri=\id$; (c) there is an isomorphism
$\theta\colon M\to M'\oplus M''$ with $\theta i(x)=(x,0)$ and $p\theta^{-1}(x,y)=y$.
\end{lemma}

\begin{proof}
(a)$\Rightarrow$(c): for $m\in M$, $m-sp(m)\in\ker p=\im i$, so $m-sp(m)=i(r(m))$ for a unique $r(m)\in M'$; put
$\theta(m)=(r(m),p(m))$, with inverse $(x,y)\mapsto i(x)+s(y)$. (b)$\Rightarrow$(c): put $\theta(m)=(r(m),p(m))$. It is injective:
if $p(m)=0$ then $m=i(x)$ and $r(m)=x$, so $r(m)=0$ forces $m=0$. It is surjective: given $(x,y)$, choose $m$ with $p(m)=y$;
then $m'=m+i(x-r(m))$ has $\theta(m')=(x,y)$. (c)$\Rightarrow$(a),(b): take $s(y)=\theta^{-1}(0,y)$ and $r$ the first component of
$\theta$.
\end{proof}

\begin{example}\label{ex:modules:exact-examples}
The sequence $0\to\Z\xrightarrow{n}\Z\to\Z/n\to0$ does not split for $n\ge2$, since $\Z$ has no element of order $n$. Likewise
$0\to\Zl\xrightarrow{\ell^k}\Zl\to\Z/\ell^k\to0$ does not split for $k\ge1$. If $T\ne0$ is a free $\Zl$-module of finite rank,
$V=T\otimes\Ql$ and $W=T\otimes\Qlmod$, the sequence $0\to T\to V\to W\to0$ (\Cref{prop:modules:t-and-w}) does not split:
a splitting would embed the torsion module $W\ne0$ into $V$, which has no nonzero torsion.
\end{example}

\begin{proposition}[two submodules]\label{prop:modules:two-submodules}
For submodules $A,B$ of $M$ the sequence
\[
0\to M/(A\cap B)\xrightarrow{\ x\mapsto(x,x)\ }M/A\oplus M/B\xrightarrow{\ (y,z)\mapsto y-z\ }M/(A+B)\to0
\]
is exact. In particular, if $A\cap B=0$ then $M\into M/A\oplus M/B$ with cokernel $M/(A+B)$. The sum $A+B$ is direct (that is,
$A\oplus B\to A+B$, $(a,b)\mapsto a+b$, is an isomorphism) if and only if $A\cap B=0$.
\end{proposition}

\begin{proof}
The first map is well defined and injective: $(x+A,x+B)=0$ iff $x\in A\cap B$. The second is surjective, and the composite
is zero. If $y-z\in A+B$, say $y-z=a+b$, then $x=y-a=z+b$ satisfies $(x+A,x+B)=(y+A,z+B)$, which is exactness in the
middle. The last statement: the kernel of $(a,b)\mapsto a+b$ is $\{(c,-c):c\in A\cap B\}$. This is the module form of the
two-ideal sequence of \Cref{ch:rings}.
\end{proof}

\begin{paperlink}\label{pl:modules:two-submodules}
In Lemma lem:crit, $I\cap J=0$ in $\T$, so \Cref{prop:modules:two-submodules} gives $\T\into\T/I\times\T/J$ with cokernel
$\T/(I+J)$; \Cref{prop:modules:algebra-duality}(iii) dualises this inclusion. In Lemma lem:ddt417 the directness of
$L[\wp]+L[I]$ is the statement $L[\wp]\cap L[I]=0$, proved in \Cref{prop:modules:ddt-ingredients}(ii).
\end{paperlink}

\begin{lemma}[the snake lemma]\label{lem:modules:snake}
Consider a commutative diagram with exact rows
\[
\begin{array}{ccccccc}
 & M' & \xrightarrow{\ i\ } & M & \xrightarrow{\ p\ } & M'' & \to 0\\
 & \big\downarrow{\scriptstyle f'} & & \big\downarrow{\scriptstyle f} & & \big\downarrow{\scriptstyle f''} & \\
0\to & N' & \xrightarrow{\ j\ } & N & \xrightarrow{\ q\ } & N'' &
\end{array}
\]
(so $fi=jf'$ and $f''p=qf$). There is an $R$-linear map $\delta\colon\ker f''\to\coker f'$ defined as follows: for
$x''\in\ker f''$ choose $x\in M$ with $p(x)=x''$; then $q(f(x))=f''(x'')=0$, so $f(x)=j(y')$ for a unique $y'\in N'$; put
$\delta(x'')=y'+f'(M')$. It is well defined, and the sequence
\[
\ker f'\xrightarrow{\ i\ }\ker f\xrightarrow{\ p\ }\ker f''\xrightarrow{\ \delta\ }\coker f'\xrightarrow{\ \bar j\ }\coker f\xrightarrow{\ \bar q\ }\coker f''
\]
is exact. If $i$ is injective, so is $\ker f'\to\ker f$; if $q$ is surjective, so is $\coker f\to\coker f''$.
\end{lemma}

\begin{proof}
The restrictions of $i$, $p$ map kernels to kernels, and $j$, $q$ induce maps of cokernels, by commutativity.
\emph{$\delta$ is well defined.} The element $y'$ is unique because $j$ is injective. If $x_1$ is another lift of $x''$, then
$x_1-x=i(m')$ with $m'\in M'$, and $f(x_1)-f(x)=fi(m')=j(f'(m'))$, so $y'$ changes by $f'(m')$, which does not change its
class. Linearity follows from the same uniqueness.
\emph{Exactness at $\ker f$.} $pi=0$. If $x\in\ker f$ and $p(x)=0$, then $x=i(m')$ and $jf'(m')=f(x)=0$, so $f'(m')=0$ as $j$ is
injective: $m'\in\ker f'$.
\emph{Exactness at $\ker f''$.} If $x''=p(x)$ with $x\in\ker f$, take this $x$ as the lift: $f(x)=0$, so $\delta(x'')=0$.
Conversely, if $\delta(x'')=0$ then $f(x)=j(f'(m'))=f(i(m'))$ for some $m'$, so $x-i(m')\in\ker f$ maps to $x''$.
\emph{Exactness at $\coker f'$.} $\bar j\delta(x'')$ is the class of $j(y')=f(x)$, which is $0$ in $\coker f$. Conversely, if
$j(y')=f(x)$ for some $x\in M$, then $x''=p(x)$ satisfies $f''(x'')=qf(x)=qj(y')=0$, and $\delta(x'')=[y']$ by construction.
\emph{Exactness at $\coker f$.} $qj=0$. If $q(y)=f''(x'')$, write $x''=p(x)$; then $q(y-f(x))=0$, so $y-f(x)=j(y')$ and
$[y]=\bar j[y']$.
The last two assertions are immediate.
\end{proof}

\begin{corollary}\label{cor:modules:snake}
\begin{enumerate}[label=(\alph*)]
\item For a short exact sequence $0\to A\to B\to C\to0$ and $n\in R$ there is an exact sequence
\[
0\to A[n]\to B[n]\to C[n]\to A/nA\to B/nB\to C/nC\to0 .
\]
\item If $\phi$ is an endomorphism of a finite abelian group $A$, then $\#\ker\phi=\#\coker\phi$.
\end{enumerate}
\end{corollary}

\begin{proof}
(a) Apply \Cref{lem:modules:snake} to the diagram whose two rows are the given sequence and whose vertical maps are
multiplication by $n$. Its kernels are the $n$-torsion and its cokernels the quotients mod $n$; both end maps behave as
stated because $A\to B$ is injective and $B\to C$ surjective. (b) $\#A=\#\ker\phi\cdot\#\im\phi$ and $\#A=\#\im\phi\cdot\#\coker\phi$.
\end{proof}

\begin{lemma}[the five lemma]\label{lem:modules:five}
Consider a commutative diagram with exact rows $M_1\to M_2\to M_3\to M_4\to M_5$ and $N_1\to N_2\to N_3\to N_4\to N_5$ and
vertical maps $f_i\colon M_i\to N_i$. If $f_2$, $f_4$ are injective and $f_1$ is surjective, then $f_3$ is injective; if $f_2,f_4$
are surjective and $f_5$ is injective, then $f_3$ is surjective. In particular, if $f_1,f_2,f_4,f_5$ are isomorphisms, so is $f_3$.
\end{lemma}

\begin{proof}
Write $a_i\colon M_i\to M_{i+1}$ and $b_i\colon N_i\to N_{i+1}$. \emph{Injectivity.} Let $f_3(x)=0$. Then $f_4a_3(x)=b_3f_3(x)=0$, so
$a_3(x)=0$ and $x=a_2(y)$. Now $b_2f_2(y)=f_3(x)=0$, so $f_2(y)=b_1(z)$, and $z=f_1(w)$. Then $f_2(y-a_1(w))=b_1(z)-b_1f_1(w)=0$, so
$y=a_1(w)$ and $x=a_2a_1(w)=0$. \emph{Surjectivity.} Let $u\in N_3$. Write $b_3(u)=f_4(x)$. Then $f_5a_4(x)=b_4b_3(u)=0$, so $a_4(x)=0$
and $x=a_3(y)$. Now $b_3(u-f_3(y))=0$, so $u-f_3(y)=b_2(v)$ with $v=f_2(w)$, and $u=f_3(y+a_2(w))$.
\end{proof}

\begin{definition}[direct sums and products]\label{def:modules:sums-products}
For a family $(M_i)_{i\in I}$ the \emph{direct product} $\prod_iM_i$ is the set of all families $(m_i)$ with componentwise
operations, and the \emph{direct sum} $\bigoplus_iM_i$ is the submodule of families with finitely many nonzero terms. If the
$M_i$ are submodules of $M$, the sum $\sum M_i$ is \emph{direct}, written $\bigoplus M_i$ (internal direct sum), if the map
$\bigoplus_iM_i\to M$ is injective.
\end{definition}

\begin{proposition}\label{prop:modules:sums}
\begin{enumerate}[label=(\roman*)]
\item By restriction to the summands and composition with the projections,
\[\Hom_R\Bigl(\bigoplus_iM_i,N\Bigr)=\prod_i\Hom_R(M_i,N),\qquad\Hom_R\Bigl(N,\prod_iM_i\Bigr)=\prod_i\Hom_R(N,M_i).\] For finite $I$, $\bigoplus_iM_i=\prod_iM_i$.
\item For submodules $A,B\subseteq M$ the following are equivalent: $M=A\oplus B$ internally; $A+B=M$ and $A\cap B=0$; there is
an idempotent $e\in\End_R(M)$ with $eM=A$ and $(1-e)M=B$.
\end{enumerate}
\end{proposition}

\begin{proof}
(i) A map out of $\bigoplus M_i$ is determined by its restrictions, which may be arbitrary because each element is a finite
sum; a map into $\prod M_i$ is the family of its components. (ii) The first two conditions agree by
\Cref{prop:modules:two-submodules}. If $M=A\oplus B$, let $e(a+b)=a$; then $e^2=e$, $eM=A$, $(1-e)M=B$. Conversely, given $e$,
every $m=em+(1-e)m$, and $x\in A\cap B$ gives $x=ex'$, $x=(1-e)x''$, so $x=ex=e(1-e)x''=0$.
\end{proof}

\begin{definition}[free modules]\label{def:modules:free}
A family $(e_i)_{i\in I}$ in $M$ is a \emph{basis} if every $m\in M$ is uniquely a finite sum $\sum r_ie_i$. A module with a
basis is \emph{free}; $R^{(I)}=\bigoplus_{i\in I}R$ is free with the standard basis, and $R^n=R^{(\{1,\dots,n\})}$. A
\emph{presentation} of $M$ is an exact sequence $F_1\to F_0\to M\to0$ with $F_0$, $F_1$ free; $M$ is \emph{finitely presented} if
$F_0$ and $F_1$ can be taken of finite rank.
\end{definition}

\begin{proposition}\label{prop:modules:free}
\begin{enumerate}[label=(\roman*)]
\item (Universal property.) For any family $(m_i)_{i\in I}$ in $N$ there is a unique homomorphism $R^{(I)}\to N$ with
$e_i\mapsto m_i$. Hence every module is a quotient of a free module and has a presentation.
\item (Rank.) If $R\ne0$ and $R^m\simeq R^n$, then $m=n$.
\item A matrix $A\in\Mat_n(R)$ is invertible if and only if $\det A\in R^\times$; an endomorphism of $R^n$ is an automorphism if and
only if the determinant of its matrix is a unit.
\end{enumerate}
\end{proposition}

\begin{proof}
(i) Send $\sum r_ie_i$ to $\sum r_im_i$. Taking all elements of $M$ as the family gives $R^{(M)}\onto M$; its kernel is again a
quotient of a free module, which gives $F_1$.
(ii) Let $\mf m$ be a maximal ideal of $R$; it exists by Zorn's lemma (\Cref{ch:rings}), or by \Cref{prop:modules:noetherian}(iii)
when $R$ is Noetherian. By \Cref{thm:modules:tensor}(ii), $R^m\simeq R^n$ gives $(R/\mf m)^m\simeq(R/\mf m)^n$ as vector spaces over the
field $R/\mf m$, so $m=n$ by linear algebra. (We use the tensor product before proving its properties only to keep the
statement here; the proof of \Cref{thm:modules:tensor} does not use (ii).)
(iii) If $AB=1$ then $\det A\det B=1$. Conversely, the adjugate satisfies $A\operatorname{adj}(A)=\operatorname{adj}(A)A=\det(A)\cdot1$ over any
commutative ring (the identity is a polynomial identity in the entries, true over $\Z[x_{ij}]$), so $A^{-1}=\det(A)^{-1}\operatorname{adj}(A)$.
\end{proof}

\begin{definition}[projective modules]\label{def:modules:projective}
A module $P$ is \emph{projective} if the equivalent conditions of \Cref{prop:modules:projective} hold.
\end{definition}

\begin{proposition}\label{prop:modules:projective}
For an $R$-module $P$ the following are equivalent:
(a) every surjection $M\onto P$ splits;
(b) $P$ is a direct summand of a free module;
(c) $\Hom_R(P,-)$ is exact, that is, carries short exact sequences to short exact sequences;
(d) for every surjection $g\colon M\onto N$ and every $h\colon P\to N$ there is $\tilde h\colon P\to M$ with $g\tilde h=h$.
Free modules are projective, and so is $eR$ for every idempotent $e\in R$; so each factor $R_i$ of $R=\prod R_i$ is a projective
$R$-module. The $\Z/6$-module $\Z/2\simeq3\Z/6$ is projective and not free.
\end{proposition}

\begin{proof}
(a)$\Rightarrow$(b): split a surjection $R^{(P)}\onto P$ and use \Cref{lem:modules:splitting}.
(b)$\Rightarrow$(d): let $F=P\oplus Q$ be free with basis $(e_i)$; extend $h$ by $0$ on $Q$, lift each image $h(e_i)$ to $M$
and use the universal property of $F$; restrict the lift to $P$.
(d)$\Rightarrow$(a): lift $\id_P$ along $M\onto P$.
(c)$\Leftrightarrow$(d): $\Hom(P,-)$ is always left exact (\Cref{prop:modules:hom}(ii)), and (d) is exactly surjectivity of
$\Hom(P,M)\to\Hom(P,N)$.
For an idempotent $e$, $R=eR\oplus(1-e)R$. Under $\Z/6\simeq\Z/2\times\Z/3$ the module $\Z/2$ is the first factor, so it is projective;
it is not free because a nonzero free $\Z/6$-module has at least $6$ elements.
\end{proof}

\begin{remark}\label{rem:modules:projective-free}
Over a local ring every finitely generated projective module is free (\Cref{thm:modules:nakayama}(v)), and over a principal
ideal domain every submodule of a free module is free (\Cref{thm:modules:submodule-free}). Consequently, if $\T$ is a
commutative $\Zl$-algebra that is finite and free as a $\Zl$-module and $e\in\T$ is an idempotent, then $e\T$ is a direct
summand of $\T$, hence a finitely generated free $\Zl$-module. This is why the local factors $\T_{\mf m}$ in Lemma lem:hecke and
Theorem thm:quant are again finite free $\Zl$-algebras, to which Lemma lem:ddt417 applies.
\end{remark}

\section{Hom and tensor product; flatness}\label{sec:modules:tensor}

The paper changes coefficients constantly: $\T\otimes\Zl$ and $\T\otimes\Q$ in Lemmas lem:crit and lem:hecke and in Theorem
thm:quant, $W=\ad T\otimes\Qlmod$ in Theorem thm:main and Lemma lem:h0, and the sentence ``the same statements hold with $\Z$
replaced by $\Zl$'' at the end of Lemma lem:crit. This section proves what makes such changes harmless: $\Hom$ is left exact,
the tensor product is right exact, and tensoring with a flat module such as $\Zl$ or $\Q$ (over $\Z$) commutes with kernels,
annihilators and saturation.

\begin{proposition}[Hom]\label{prop:modules:hom}
\begin{enumerate}[label=(\roman*)]
\item $\Hom_R(R,M)\simeq M$ by $f\mapsto f(1)$; $\Hom_R(R^n,M)\simeq M^n$; and $\Hom_R(R/\mf a,M)\simeq M[\mf a]$ by $f\mapsto f(1)$.
\item $\Hom_R(N,-)$ and $\Hom_R(-,N)$ are left exact: if $0\to M'\xrightarrow{i}M\xrightarrow{p}M''$ is exact then so is
$0\to\Hom(N,M')\to\Hom(N,M)\to\Hom(N,M'')$, and if $M'\xrightarrow{i}M\xrightarrow{p}M''\to0$ is exact then so is
$0\to\Hom(M'',N)\to\Hom(M,N)\to\Hom(M',N)$.
\item Neither is right exact in general: applying $\Hom_\Z(\Z/2,-)$ to $0\to\Z\xrightarrow{2}\Z\to\Z/2\to0$ gives $0\to0\to0\to\Z/2$, and
applying $\Hom_\Z(-,\Z)$ gives $0\to0\to\Z\xrightarrow{2}\Z$; in both the last map is not surjective.
\item For $M$ free with basis $e_1,\dots,e_n$, $M^\vee=\Hom_R(M,R)$ is free with the \emph{dual basis} $e_i^\vee$,
$e_i^\vee(e_j)=\delta_{ij}$, and the evaluation map $M\to M^{\vee\vee}$, $m\mapsto(\varphi\mapsto\varphi(m))$, is an isomorphism. Over a
domain $M^\vee$ is torsion-free for every $M$.
\end{enumerate}
\end{proposition}

\begin{proof}
(i) A homomorphism out of $R$ is determined by the image of $1$, which may be arbitrary; out of $R/\mf a$ the image of $1$
must be killed by $\mf a$, and every such element occurs. (ii) For the first sequence: if $if=0$ then $f=0$ since $i$ is
injective; if $pf=0$ then $f(N)\subseteq\ker p=i(M')$ and $f=ig$ with $g=i^{-1}f$. For the second: if $fp=0$ then $f=0$ as $p$ is
surjective; if $fi=0$ then $f$ vanishes on $\ker p$ and factors through $M/\ker p\simeq M''$. (iii) $\Hom(\Z/2,\Z)=0$, and
$\Hom(\Z,\Z)=\Z$ with the induced map multiplication by $2$. (iv) $\varphi=\sum\varphi(e_i)e_i^\vee$, and the evaluation map sends $e_i$ to
the dual basis of $(e_i^\vee)$. Over a domain, $r\varphi=0$ with $r\ne0$ gives $r\varphi(m)=0$, so $\varphi(m)=0$ for all $m$.
\end{proof}

\begin{theorem}[tensor product]\label{thm:modules:tensor}
\begin{enumerate}[label=(\roman*)]
\item For $R$-modules $M,N$ there is an $R$-module $M\otimes_RN$ with an $R$-bilinear map $M\times N\to M\otimes_RN$,
$(m,n)\mapsto m\otimes n$, such that every bilinear map $b\colon M\times N\to P$ factors uniquely as $b(m,n)=\beta(m\otimes n)$ with
$\beta$ linear. It is unique up to unique isomorphism, and generated by the elements $m\otimes n$.
\item $R\otimes M\simeq M$; $(\bigoplus_iM_i)\otimes N\simeq\bigoplus_i(M_i\otimes N)$; $R^m\otimes R^n\simeq R^{mn}$; $(R/\mf a)\otimes M\simeq M/\mf aM$;
$M\otimes N\simeq N\otimes M$ and $(M\otimes N)\otimes P\simeq M\otimes(N\otimes P)$.
\item (Adjunction.) $\Hom_R(M\otimes N,P)\simeq\Hom_R(M,\Hom_R(N,P))$, $\beta\mapsto(m\mapsto(n\mapsto\beta(m\otimes n)))$.
\item (Base change.) For an $R$-algebra $S$, $S\otimes_RM$ is an $S$-module by $s(s'\otimes m)=ss'\otimes m$; $S\otimes_RR^n\simeq S^n$ and
$S\otimes_R(R/\mf a)\simeq S/\mf aS$. If $A$, $B$ are $R$-algebras, $A\otimes_RB$ is an $R$-algebra with $(a\otimes b)(a'\otimes b')=aa'\otimes bb'$.
\end{enumerate}
\end{theorem}

\begin{proof}
(i) Let $F$ be the free module on the set $M\times N$ and $F_0$ its submodule generated by the elements
$(m+m',n)-(m,n)-(m',n)$, $(m,n+n')-(m,n)-(m,n')$, $(rm,n)-r(m,n)$ and $(m,rn)-r(m,n)$. Put $M\otimes N=F/F_0$ and $m\otimes n$ the
class of $(m,n)$; the map is bilinear by construction. A bilinear $b$ defines a linear map $F\to P$ (\Cref{prop:modules:free}(i))
that kills $F_0$ because $b$ is bilinear, hence $\beta$; it is unique because the $m\otimes n$ generate. Uniqueness up to unique
isomorphism is the usual consequence of a universal property.
(ii) Each isomorphism is produced by the universal property from an evident bilinear map, with inverse written down
directly: $r\otimes m\mapsto rm$ with inverse $m\mapsto1\otimes m$; $(m_i)\otimes n\mapsto(m_i\otimes n)$ with inverse given on each summand;
$R^m\otimes R^n$ follows, with basis $e_i\otimes e_j$; for $R/\mf a$, the bilinear map $(\bar r,m)\mapsto rm+\mf aM$ gives
$(R/\mf a)\otimes M\to M/\mf aM$, with inverse $m+\mf aM\mapsto\bar1\otimes m$, well defined because $\bar1\otimes am=\bar a\otimes m=0$ for $a\in\mf a$.
Commutativity and associativity are proved in the same way.
(iii) Bilinear maps $M\times N\to P$ are exactly the linear maps $M\to\Hom(N,P)$.
(iv) For fixed $s$, multiplication by $s$ on $S$ is $R$-linear and induces the action on $S\otimes M$; the rest follows from (ii).
The product on $A\otimes B$ is induced by the $R$-multilinear map $(a,b,a',b')\mapsto aa'\otimes bb'$.
\end{proof}

\begin{example}\label{ex:modules:tensor-table}
Over $\Z$ (all isomorphisms by \Cref{thm:modules:tensor}(ii)):
\begin{center}
\begin{tabular}{ll}
\toprule
tensor product & value\\
\midrule
$\Z/m\otimes\Z/n$ & $\Z/(m\Z+n\Z)=\Z/\gcd(m,n)$\\
$\Q\otimes\Z/n$ & $\Q/n\Q=0$\\
$\Zl\otimes\Z/n$ & $\Zl/n\Zl\simeq\Z/\ell^{v_\ell(n)}$ (\Cref{thm:modules:zl}(iii): $n=\ell^{v_\ell(n)}u$ with $u\in\Zl^\times$)\\
$\Qlmod\otimes A$, $A$ torsion & $0$\\
$\Qlmod\otimes\Qlmod$ & $0$\\
\bottomrule
\end{tabular}
\end{center}
For the fourth row, if $na=0$ write $x=ny$ in the divisible group $\Qlmod$; then $x\otimes a=y\otimes na=0$. The fifth row is a
case of the fourth. The same argument shows $D\otimes A=0$ whenever $D$ is divisible and $A$ is torsion.
\end{example}

\begin{proposition}[right exactness]\label{prop:modules:right-exact}
If $M'\xrightarrow{i}M\xrightarrow{p}M''\to0$ is exact, so is $M'\otimes N\to M\otimes N\to M''\otimes N\to0$. The functor $-\otimes N$ is not left
exact: tensoring the injection $\Z\xrightarrow{2}\Z$ with $\Z/2$ gives the zero map $\Z/2\to\Z/2$. A module $N$ is \emph{flat} if
$-\otimes_RN$ is exact (carries injections to injections). Free modules, and more generally projective modules, are flat.
\end{proposition}

\begin{proof}
First a criterion: a sequence $A\xrightarrow{a}B\xrightarrow{b}C\to0$ is exact if $0\to\Hom(C,P)\to\Hom(B,P)\to\Hom(A,P)$ is exact for every
module $P$. Indeed, with $P=\coker b$ the projection $C\to P$ restricts to $0$ on $B$, so it is $0$, and $b$ is surjective; with
$P=C$, the image of $\id_C$ in $\Hom(A,C)$ is $ba$, which is $0$ because the composite of the two maps of $\Hom$ sequences is
$0$; with $P=B/a(A)$, the projection $B\to P$ kills $a(A)$, so it factors through $b$, and $\ker b\subseteq a(A)$.
By the adjunction (\Cref{thm:modules:tensor}(iii)) the sequence $0\to\Hom(M''\otimes N,P)\to\Hom(M\otimes N,P)\to\Hom(M'\otimes N,P)$ is
the sequence $0\to\Hom(M'',Q)\to\Hom(M,Q)\to\Hom(M',Q)$ with $Q=\Hom(N,P)$, which is exact by \Cref{prop:modules:hom}(ii).
The criterion gives right exactness. $R\otimes-$ is the identity, a direct sum of injections is an injection, and a direct
summand of a flat module is flat; so free and projective modules are flat.
\end{proof}

\begin{proposition}[flatness]\label{prop:modules:flat}
\begin{enumerate}[label=(\roman*)]
\item For a multiplicative set $S\subseteq R$, the ring $S^{-1}R$ is a flat $R$-module.
\item Over a principal ideal domain a module is flat if and only if it is torsion-free. Hence $\Q$, $\Z_{(\ell)}$, $\Zl$ and $\Ql$
are flat over $\Z$, and $\Ql$ is flat over $\Zl$.
\item Let $R\to S$ be flat. Then $S\otimes_R-$ commutes with kernels, images, cokernels and finite intersections of submodules.
If $\mf a=(a_1,\dots,a_k)$ is finitely generated, $M[\mf a]\otimes_RS=(M\otimes_RS)[\mf aS]$ inside $M\otimes_RS$. If $M$ is finitely
generated, $\Ann_S(M\otimes_RS)=\Ann_R(M)S$.
\item If $R\subseteq S$ are principal ideal domains with $S$ flat over $R$, and $L'\subseteq L$ with $L$ a lattice over $R$, then
the saturation of $L'\otimes S$ in $L\otimes S$ is $\operatorname{sat}(L')\otimes S$ (\Cref{def:modules:saturated}).
\end{enumerate}
\end{proposition}

\begin{proof}
(i) By \Cref{thm:modules:localisation}(ii) below, $S^{-1}R\otimes_RM\simeq S^{-1}M$ and $M\mapsto S^{-1}M$ is exact; that proof does not use
this proposition.
(ii) If $N$ is flat and $r\ne0$, tensoring the injection $R\xrightarrow{r}R$ with $N$ shows that $r$ is injective on $N$. Conversely let
$N$ be torsion-free. If $N$ is finitely generated it is free (\Cref{prop:modules:torsion}(ii)), hence flat. In general let
$f\colon M'\to M$ be injective and $z=\sum_im'_i\otimes n_i\in M'\otimes N$ with $(f\otimes1)(z)=0$. In the construction of
\Cref{thm:modules:tensor}(i), this says that $\sum_i(f(m'_i),n_i)$ is a finite combination of the defining relations of
$M\otimes N$; these involve finitely many elements of $N$. Let $N_0\subseteq N$ be the submodule generated by them and the
$n_i$. The same relations show that $z_0=\sum m'_i\otimes n_i\in M'\otimes N_0$ maps to $0$ in $M\otimes N_0$. As $N_0$ is finitely
generated and torsion-free, it is flat, so $z_0=0$, and $z$, the image of $z_0$, is $0$.
(iii) Kernels, images and cokernels: tensor the exact sequences $0\to\ker g\to M\to M''$ and $M\to M''\to\coker g\to0$ with
$S$. Intersections: $A\cap B=\ker(M\to M/A\oplus M/B)$. Torsion: $M[\mf a]$ is the kernel of $M\to M^k$, $m\mapsto(a_im)$, and
after tensoring the kernel of the same map is $(M\otimes S)[\mf aS]$. Annihilators: for generators $m_1,\dots,m_n$,
$\Ann_R(M)$ is the kernel of $R\to M^n$, $r\mapsto(rm_i)$, and after tensoring the kernel of $S\to(M\otimes S)^n$ is
$\Ann_S(M\otimes S)$, because the $m_i\otimes1$ generate $M\otimes S$ over $S$.
(iv) Write $L/L'\simeq F\oplus C$ with $F$ free and $C$ torsion (\Cref{thm:modules:structure-pid}); then $\operatorname{sat}(L')$ is the
kernel of $L\to L/L'\to F$. After tensoring, $\operatorname{sat}(L')\otimes S$ is the kernel of $L\otimes S\to F\otimes S$, a free target, so it is
saturated in $L\otimes S$; it contains $L'\otimes S$ with quotient $C\otimes S$, which is killed by the nonzero element of $R$ killing
$C$, hence torsion. A saturated submodule containing $L'\otimes S$ with torsion quotient is its saturation
(\Cref{prop:modules:saturated}).
\end{proof}

\begin{paperlink}\label{pl:modules:flat}
``The same statements hold with $\Z$ replaced by $\Zl$'' (Lemma lem:crit) and ``every step holds verbatim over $\Zl$'' are
\Cref{prop:modules:flat}: $\Zl$ is flat over $\Z$, so $\ker\phi_f\otimes\Zl$ is the kernel of $\phi_f\otimes\Zl$, the annihilator
$\Ann_\T(\ker\phi_f)\otimes\Zl$ is the annihilator in $\T\otimes\Zl$, and the congruence ideal $\phi_f(\Ann_\T\ker\phi_f)$ becomes
$\phi_f(\Ann_\T\ker\phi_f)\Zl$, whose generator is the $\ell$-part of $r$. The same proposition is used in Step 1 of the proof of
Theorem thm:quant, which works over $\Zl$ from the start.
\end{paperlink}

\section[Noetherian rings, Nakayama, localisation]{Noetherian rings, Nakayama's lemma, local rings and localisation}\label{sec:modules:local}

Every ring in the paper is Noetherian: $\Z$, $\Zl$, the Hecke algebra $\T$ (a finitely generated $\Z$-module) and its
completions $\T_{\mf m}$. Nakayama's lemma is used twice in the paper: in Remark rem:free, where ``$J_0(N)[\mf m]$ is
two-dimensional'' becomes ``the Tate module is generated by two elements over $\T_{\mf m}$'' (\Cref{thm:modules:mazur-freeness}),
and in Lemma lem:h0, where a homomorphism $T/\ell^k\to(T/\ell^k)(1)$ that is nonzero modulo $\ell$ is an isomorphism. Local rings
and localisation give the meaning of $\T_{\mf m}$.

\begin{proposition}[Noetherian modules]\label{prop:modules:noetherian}
\begin{enumerate}[label=(\roman*)]
\item For an $R$-module $M$ the following are equivalent: every ascending chain of submodules is eventually constant; every
nonempty set of submodules has a maximal element; every submodule is finitely generated. Such $M$ is \emph{Noetherian}; $R$
is a Noetherian ring if it is Noetherian as an $R$-module, that is, if every ideal is finitely generated.
\item Submodules and quotients of Noetherian modules are Noetherian, and if $0\to M'\to M\to M''\to0$ is exact with $M'$, $M''$
Noetherian then $M$ is Noetherian. Over a Noetherian ring every finitely generated module is Noetherian; hence submodules of
finitely generated modules are finitely generated, and finitely generated modules are finitely presented.
\item In a nonzero Noetherian ring every proper ideal is contained in a maximal ideal.
\item (Hilbert's basis theorem.) If $R$ is Noetherian, so is $R[x]$.
\item Fields, principal ideal domains, $\Zl$ (\Cref{thm:modules:zl}), and every ring that is finitely generated as a module over a
Noetherian ring (such as the Hecke algebra $\T$ over $\Z$, or $\T\otimes\Zl$ over $\Zl$) are Noetherian.
\end{enumerate}
\end{proposition}

\begin{proof}
(i) If chains stabilise and a nonempty set $\Sigma$ of submodules had no maximal element, we could choose a strictly
increasing infinite chain in $\Sigma$. If maximal elements exist and $N\subseteq M$, a maximal element of the set of finitely
generated submodules of $N$ is $N$ itself (otherwise add an element of $N$ outside it). If every submodule is finitely
generated and $N_1\subseteq N_2\subseteq\cdots$, then $\bigcup N_i$ is generated by finitely many elements, all in some $N_k$, so the chain
is constant from $k$ on.
(ii) Submodules: by the third condition. Quotients: by the first, via \Cref{thm:modules:isomorphism}(iii). Extensions: if
$N_1\subseteq N_2\subseteq\cdots$ in $M$, both $N_i\cap M'$ and the images of $N_i$ in $M''$ stabilise, and then so does $N_i$ (if
$N\subseteq N'$ have the same intersection with $M'$ and the same image in $M''$, then $N=N'$). By induction $R^n$ is Noetherian,
and a finitely generated module is a quotient of some $R^n$. The kernel of $R^n\onto M$ is then finitely generated.
(iii) The set of proper ideals containing a given proper ideal is nonempty and has a maximal element by (i), which is a maximal
ideal.
(iv) We do not prove this; see \cite[Theorem 7.5]{AtiyahMacdonald} (number not verified against the printed edition).
(v) Fields and principal ideal domains: every ideal is generated by one element. A ring $A$ finite over a Noetherian $R$ is a
Noetherian $R$-module by (ii), and its ideals are $R$-submodules.
\end{proof}

\begin{lemma}[the determinant trick]\label{lem:modules:determinant-trick}
Let $M$ be generated by $n$ elements, $\mf a\subseteq R$ an ideal and $\phi\in\End_R(M)$ with $\phi(M)\subseteq\mf aM$. Then
$\phi^n+a_1\phi^{n-1}+\cdots+a_n=0$ in $\End_R(M)$ for some $a_i\in\mf a^i$. Consequently:
\begin{enumerate}[label=(\alph*)]
\item if $M$ is finitely generated and $\mf aM=M$, there is $x\in R$ with $x\equiv1\pmod{\mf a}$ and $xM=0$;
\item a surjective endomorphism of a finitely generated module is an isomorphism.
\end{enumerate}
\end{lemma}

\begin{proof}
Let $m_1,\dots,m_n$ generate $M$ and write $\phi(m_i)=\sum_ja_{ij}m_j$ with $a_{ij}\in\mf a$. Work in the commutative ring $R[\phi]\subseteq\End_R(M)$
and let $B=(\delta_{ij}\phi-a_{ij})$, a matrix over it. Then $B$ applied to the column $(m_j)$ is $0$. Multiplying by the adjugate of
$B$ gives $\det(B)m_j=0$ for all $j$, so $\det B=0$ in $\End_R(M)$; expanding the determinant gives the stated polynomial with
$a_i\in\mf a^i$. (a) Take $\phi=\id$: then $x=1+a_1+\cdots+a_n$ kills $M$. (b) Let $\phi$ be surjective and regard $M$ as an
$R[X]$-module, $X$ acting as $\phi$; it is finitely generated over $R[X]$, and $(X)M=M$. By (a) there is $g\in R[X]$ with
$(1-Xg(X))M=0$, that is, $\phi\circ g(\phi)=\id$; so $\phi$ is injective.
\end{proof}

\begin{definition}[local rings]\label{def:modules:local-ring}
A ring is \emph{local} if it has exactly one maximal ideal $\mf m$; we write $(R,\mf m,k)$ with $k=R/\mf m$ the \emph{residue field}.
The \emph{Jacobson radical} $\mf J(R)$ is the intersection of all maximal ideals.
\end{definition}

\begin{proposition}\label{prop:modules:local-ring}
\begin{enumerate}[label=(\roman*)]
\item $R$ is local if and only if its non-units form an ideal; then this ideal is $\mf m$, and $R^\times=R\setminus\mf m$.
\item If $x\in\mf J(R)$ then $1+x$ is a unit.
\item Examples of local rings: fields; $\Z_{(p)}$ (\Cref{ch:rings}); discrete valuation rings (\Cref{thm:modules:dvr}); $\Z/p^n$;
$k[x]/(x^2)$; $\Zl[x]/(x^2)$; each factor $\T_{\mf m}$ of a finite free $\Zl$-algebra (\Cref{thm:modules:idempotent-lifting}).
Non-examples: $\Z$; $\Zl\times\Zl$ (maximal ideals $\ell\Zl\times\Zl$ and $\Zl\times\ell\Zl$); $\Z_{(11)}[x]/(x^2+x-1)$, which has two
maximal ideals (\Cref{thm:modules:idempotent-lifting}(v)).
\end{enumerate}
\end{proposition}

\begin{proof}
(i) A non-unit $x$ generates a proper ideal, which lies in a maximal ideal (\Cref{ch:rings}, by Zorn's lemma; or
\Cref{prop:modules:noetherian}(iii)). If $R$ is local this ideal is $\mf m$, so the non-units are exactly $\mf m$. Conversely if
the non-units form an ideal $\mf n$, every proper ideal consists of non-units and lies in $\mf n$, which is therefore the unique
maximal ideal. (ii) If $1+x$ were not a unit it would lie in a maximal ideal $\mf m'$; so would $x$, and then $1\in\mf m'$.
(iii) In $\Z/p^n$ and $k[x]/(x^2)$, $\Zl[x]/(x^2)$ the non-units are the multiples of $p$, of $x$, of $(\ell,x)$; these form
ideals. The others are proved where cited.
\end{proof}

\begin{theorem}[Nakayama's lemma]\label{thm:modules:nakayama}
Let $(R,\mf m,k)$ be a local ring and $M$ a finitely generated $R$-module.
\begin{enumerate}[label=(\roman*)]
\item If $\mf mM=M$ then $M=0$. More generally, for any ring $R$ and any ideal $\mf a\subseteq\mf J(R)$, $\mf aM=M$ implies $M=0$.
\item If $N\subseteq M$ is a submodule with $N+\mf mM=M$, then $N=M$.
\item Elements $x_1,\dots,x_r$ generate $M$ if and only if their images span the $k$-vector space $M/\mf mM$. The minimal number of
generators of $M$ is $\dim_kM/\mf mM$.
\item A homomorphism $f\colon N'\to M$ is surjective if and only if $f\otimes k\colon N'/\mf mN'\to M/\mf mM$ is. A matrix in $\Mat_n(R)$ is
invertible if and only if its reduction modulo $\mf m$ is; an endomorphism of $R^n$ is an automorphism if and only if it is so
modulo $\mf m$.
\item Every finitely generated projective $R$-module is free.
\end{enumerate}
\end{theorem}

\begin{proof}
(i) By \Cref{lem:modules:determinant-trick}(a) some $x=1+a$ with $a\in\mf a$ kills $M$; $x$ is a unit by
\Cref{prop:modules:local-ring}(ii), so $M=0$. (ii) Apply (i) to $M/N$, which satisfies $\mf m(M/N)=M/N$. (iii) Apply (ii) to
$N=\sum Rx_i$; conversely generators of $M$ span $M/\mf mM$. A minimal generating set therefore has $\dim_kM/\mf mM$ elements.
(iv) The first statement is (ii) for $N=f(N')$. A matrix is invertible iff its determinant is a unit
(\Cref{prop:modules:free}(iii)), iff the determinant is not in $\mf m$, iff the reduced determinant is nonzero.
(v) Let $P$ be finitely generated projective and $n=\dim_kP/\mf mP$. By (iii) there is a surjection $R^n\onto P$; it splits, so
$R^n\simeq P\oplus Q$ with $Q$ finitely generated (a quotient of $R^n$). Reducing modulo $\mf m$ gives $k^n\simeq P/\mf mP\oplus Q/\mf mQ$,
so $Q/\mf mQ=0$ and $Q=0$ by (i).
\end{proof}

\begin{example}\label{ex:modules:nakayama}
A finitely generated $\Zl$-module $M$ with $M/\ell M\simeq\Fl^2$ is generated by two elements. An endomorphism of $(\Z/\ell^k)^n$ whose
reduction modulo $\ell$ is invertible is invertible (apply (iv) over the local ring $\Z/\ell^k$).
\end{example}

\begin{paperlink}\label{pl:modules:nakayama}
In the proof of Lemma lem:h0 the paper writes: ``An element of exact order $\ell^k$ is a map $\phi\colon T/\ell^k\to(T/\ell^k)(1)$ whose
reduction modulo $\ell$ is nonzero, hence an isomorphism, so $\phi$ is an isomorphism by Nakayama's lemma''. The underlying
module map is an endomorphism of $(\Z/\ell^k)^2$ (the twist does not change the module), invertible modulo $\ell$ because a nonzero
$G_\Q$-map between irreducible representations is an isomorphism, its kernel and image being subrepresentations (\Cref{ch:reps}); it
is therefore invertible by \Cref{ex:modules:nakayama}. In Remark rem:free, the two-dimensionality of $J_0(N)[\mf m]$ is turned into
two generators by \Cref{thm:modules:nakayama}(iii); see \Cref{thm:modules:mazur-freeness}.
\end{paperlink}

\begin{theorem}[localisation]\label{thm:modules:localisation}
Let $S\subseteq R$ be multiplicative ($1\in S$, $SS\subseteq S$). Let $S^{-1}M$ be the set of fractions $m/s$ ($m\in M$, $s\in S$), where
$m/s=m'/s'$ if and only if $u(s'm-sm')=0$ for some $u\in S$; with the evident operations $S^{-1}R$ is a ring and $S^{-1}M$ an
$S^{-1}R$-module.
\begin{enumerate}[label=(\roman*)]
\item The map $R\to S^{-1}R$, $r\mapsto r/1$, sends $S$ to units, and every ring homomorphism $R\to A$ sending $S$ to units factors
uniquely through it.
\item $S^{-1}R\otimes_RM\simeq S^{-1}M$, $(a/s)\otimes m\mapsto am/s$, and $M\mapsto S^{-1}M$ is exact.
\item $\mf b\mapsto\mf b\cap R$ (preimage) and $\mf p\mapsto S^{-1}\mf p$ are inverse bijections between the prime ideals of $S^{-1}R$ and the
prime ideals of $R$ not meeting $S$.
\item For a prime $\mf p$, $R_{\mf p}=(R\setminus\mf p)^{-1}R$ is local with maximal ideal $\mf pR_{\mf p}$; we write $M_{\mf p}$ for the localised module.
\item $M=0$ if and only if $M_{\mf m}=0$ for every maximal ideal $\mf m$; a homomorphism is injective (surjective) if and only if it is
so after localising at every maximal ideal.
\item Localisation commutes with finite direct sums, quotients, and annihilators of finitely generated modules:
$S^{-1}\Ann_R(M)=\Ann_{S^{-1}R}(S^{-1}M)$. It commutes with $\Hom$ out of finitely presented modules (stated; see
\cite[Chapter 3]{AtiyahMacdonald}).
\item Examples: $\Z_{(p)}$; $\Q=\Z_{(0)}$; $\Z[1/\ell]$; $(\Z/n)_{(p)}\simeq\Z/p^{v_p(n)}$; for $R=R_1\times R_2$ and the prime
$\mf p=\mf p_1\times R_2$, $R_{\mf p}\simeq(R_1)_{\mf p_1}$.
\end{enumerate}
\end{theorem}

\begin{proof}
The relation is an equivalence relation (transitivity uses the factor $u$ and that $S$ is closed under products), and the
operations are well defined by the usual computations.
(i) $s/1$ has inverse $1/s$. If $\phi\colon R\to A$ inverts $S$, put $\psi(r/s)=\phi(r)\phi(s)^{-1}$; if $u(s'r-sr')=0$ then
$\phi(u)$ is a unit and $\phi(r)\phi(s)^{-1}=\phi(r')\phi(s')^{-1}$. Uniqueness is forced by $\psi(r/s)\psi(s/1)=\psi(r/1)$.
(ii) The bilinear map $(a/s,m)\mapsto am/s$ gives $\alpha\colon S^{-1}R\otimes M\to S^{-1}M$. Every element of $S^{-1}R\otimes M$ has the form
$(1/s)\otimes m$ (bring the finitely many fractions to a common denominator). Define $\beta(m/s)=(1/s)\otimes m$; it is well
defined: if $u(s'm-sm')=0$ then $(1/s)\otimes m=(us'/uss')\otimes m=(1/uss')\otimes us'm=(1/uss')\otimes usm'=(1/s')\otimes m'$. Clearly
$\alpha\beta=\id$ and $\beta\alpha=\id$. Exactness: let $M'\xrightarrow{f}M\xrightarrow{g}M''$ be exact at $M$ and $g(m)/s=0$. Then $ug(m)=0$ for some
$u\in S$, so $um=f(m')$ and $m/s=f(m')/us$.
(iii) If $\mf b$ is prime in $S^{-1}R$, its preimage is prime and misses $S$ (it contains no unit), and $\mf b=S^{-1}(\mf b\cap R)$
because $r/s\in\mf b$ iff $r/1\in\mf b$. If $\mf p$ is prime with $\mf p\cap S=\emptyset$, then $S^{-1}\mf p$ is a proper ideal, it is prime
(if $(a/s)(b/t)\in S^{-1}\mf p$ then $uab\in\mf p$ for some $u\in S$, and $u\notin\mf p$), and its preimage is $\mf p$ (if $r/1=p/s$ then
$u(sr-p)=0$, so $usr\in\mf p$ and $r\in\mf p$).
(iv) By (iii) the primes of $R_{\mf p}$ correspond to the primes contained in $\mf p$, so $\mf pR_{\mf p}$ is the unique maximal ideal.
(v) If $0\ne m\in M$, then $\Ann(m)$ is proper and lies in a maximal ideal $\mf m$; then $m/1\ne0$ in $M_{\mf m}$, since $um=0$ would put
$u\in\Ann(m)\subseteq\mf m$. Apply this to kernels and cokernels, which commute with localisation by (ii).
(vi) Direct sums and quotients: by (ii) and \Cref{thm:modules:tensor}(ii). Annihilators: $S^{-1}R$ is flat by (ii), so
\Cref{prop:modules:flat}(iii) applies.
(vii) The first three are the definitions (with $S=\Z\setminus p\Z$, $\Z\setminus\{0\}$, $\{\ell^k\}$). For $\Z/n=\Z/p^a\times\Z/m$ with $p\nmid m$: $m$
is a unit of $\Z_{(p)}$ and kills $\Z/m$, so $(\Z/m)_{(p)}=0$, while every element of $\Z\setminus p\Z$ acts invertibly on $\Z/p^a$. For the
product, $(1,0)\notin\mf p$ becomes a unit in $R_{\mf p}$ and $(1,0)(0,1)=0$, so $(0,1)=0$ in $R_{\mf p}$, and $R_{\mf p}=(R_1\times0)_{\mf p}\simeq(R_1)_{\mf p_1}$.
\end{proof}

\begin{remark}\label{rem:modules:localisation-factor}
The last example of (vii) is the reason why, for a finite algebra over a \emph{complete} discrete valuation ring, the
localisation at a maximal ideal is a direct factor (\Cref{thm:modules:idempotent-lifting}): completeness produces the
decomposition into a product, and localisation then picks out one factor. Without completeness there is in general no
such decomposition (\Cref{thm:modules:idempotent-lifting}(v)), which is why the paper always works with $\T\otimes\Zl$ (\Cref{conv:modules:tm}).
\end{remark}

\section[Modules over a principal ideal domain]{Finitely generated modules over a principal ideal domain; lattices and saturation}\label{sec:modules:pid}

All the lattices of the paper live over the principal ideal domains $\Z$ and $\Zl$: $S_2(\Gamma_0(N);\Z)$ and its sublattices
$\Z f$ and $L=f^\perp\cap S_2(\Gamma_0(N);\Z)$ in Lemma lem:crit; $T_3E$ in Lemma lem:IV; the lattice $L$ and its saturated
sublattices $L[\wp]$, $L[I]$ in Lemma lem:ddt417; the module $H^1(I,M)\simeq(\Zl/n')^2\oplus\Zl$ in Lemma lem:tamfactor. This section proves the
structure theorem through the Smith normal form, introduces saturation, and ends with the two results that carry the
second paragraph of the proof of Lemma lem:crit: the dual of an inclusion of lattices (\Cref{prop:modules:dual-inclusion}) and
the congruence module of a decomposition (\Cref{prop:modules:congruence-decomposition}). Throughout, $R$ is a principal ideal
domain with fraction field $K$.

\begin{theorem}[submodules of free modules]\label{thm:modules:submodule-free}
A submodule of a free $R$-module of rank $n$ is free of rank at most $n$. Consequently finitely generated $R$-modules are finitely
presented, and finitely generated projective $R$-modules are free.
\end{theorem}

\begin{proof}
Induction on $n$, the case $n=0$ being trivial. Let $N\subseteq R^n$ and let $\pi_n\colon R^n\to R$ be the last coordinate. Then
$\pi_n(N)$ is an ideal, $aR$ say. If $a=0$, $N\subseteq R^{n-1}$ and we are done. Otherwise choose $x\in N$ with $\pi_n(x)=a$. Then
$N=(N\cap R^{n-1})\oplus Rx$: every $y\in N$ has $\pi_n(y)=ra$, and $y-rx\in N\cap R^{n-1}$; the sum is direct because $\pi_n(rx)=ra=0$
forces $r=0$ in the domain $R$. By induction $N\cap R^{n-1}$ is free of rank at most $n-1$, and $Rx\simeq R$. The consequences: the
kernel of $R^n\onto M$ is free of finite rank, and a projective module is a direct summand, hence a submodule, of a free module.
\end{proof}

\begin{theorem}[Smith normal form]\label{thm:modules:smith}
For $A\in\Mat_{m\times n}(R)$ there are $U\in\GL_m(R)$ and $V\in\GL_n(R)$ such that $UAV$ is the matrix with diagonal entries
$d_1,\dots,d_s,0,\dots,0$ and zeros elsewhere, where $d_i\ne0$ and $d_1\mid d_2\mid\cdots\mid d_s$. For each $k$ the ideal $D_k(A)$ generated by the
$k\times k$ minors of $A$ (the $k$-th \emph{determinantal ideal}) satisfies $D_k(UAV)=D_k(A)$; hence $D_k(A)=(d_1\cdots d_k)$ for $k\le s$
and $0$ for $k>s$, and the ideals $(d_i)$, the \emph{invariant factors} of $A$, are uniquely determined.
\end{theorem}

\begin{proof}
For $a\ne0$ let $\lambda(a)$ be the number of prime factors of $a$ counted with multiplicity ($R$ is a unique factorisation domain,
\Cref{ch:rings}); units have $\lambda=0$. We use row and column operations by invertible matrices: permutations, adding a
multiple of one row (column) to another, and the following \emph{Bezout operation}: if $g=\gcd(a,b)=xa+yb$ with $a,b$ not both $0$,
then $\Mtwo{x}{y}{-b/g}{a/g}$ has determinant $1$ and sends the column $(a,b)^{\mathsf T}$ to $(g,0)^{\mathsf T}$; acting on two rows (or, transposed, on
two columns) it replaces $a$ by $g$.
If $A=0$ there is nothing to do. Otherwise move a nonzero entry with $\lambda$ minimal to position $(1,1)$; call it $a$. If $a$
does not divide some entry $b$ of the first column, the Bezout operation on the two rows replaces $a$ by $g$ with $\lambda(g)<\lambda(a)$;
similarly for the first row. Since $\lambda$ cannot decrease forever, after finitely many steps $a$ divides every entry of the first row
and column, which we clear by elementary operations. Now $A=\operatorname{diag}(a,A')$. If $a$ does not divide an entry $b$ of $A'$, add the row of
$b$ to the first row: the first row now contains $b$, and the previous step lowers $\lambda(a)$ again. Eventually $a$ divides every
entry of $A'$, and this persists under operations on $A'$, whose entries are combinations of the old ones. Induction on the size
finishes the existence, with $d_1=a$.
For invariance: each $k\times k$ minor of $UA$ is an $R$-combination of $k\times k$ minors of $A$ (Cauchy-Binet, or expand the rows of $UA$ as
combinations of rows of $A$ by multilinearity), so $D_k(UA)\subseteq D_k(A)$, and equality follows by applying the same to $U^{-1}$; columns
likewise. For a diagonal matrix with $d_1\mid\cdots\mid d_s$ the nonzero $k\times k$ minors are products of $k$ of the $d_i$, all divisible by
$d_1\cdots d_k$. Hence $(d_1\cdots d_k)$ is determined by $A$, and so is $(d_k)=(d_1\cdots d_k)/(d_1\cdots d_{k-1})$.
\end{proof}

\begin{example}\label{ex:modules:smith}
Computed by hand and checked with PARI's \texttt{matsnf}, over $\Z$:
$\Mtwo2468\sim\operatorname{diag}(2,4)$; the relation matrix with rows $(2,4,4)$, $(-6,6,12)$, $(10,-4,-16)$ (from \cite[Section 12.1]{DummitFoote})
$\sim\operatorname{diag}(2,6,12)$; $\Mtwo3693\sim\operatorname{diag}(3,15)$ (its determinant is $-45$); the $3\times2$ matrix with rows $(6,0)$, $(0,4)$, $(2,2)$
$\sim\operatorname{diag}(2,2)$ (the gcd of its entries is $2$ and its $2\times2$ minors are $24$, $12$, $-8$, with gcd $4$). Over $\Z_3$ the matrix
$\Mtwo3693$ is equivalent to $\operatorname{diag}(3,3)$, since $15=3\cdot5$ and $5\in\Z_3^\times$.
\begin{lstlisting}[language=GP]
matsnf([2,4;6,8])                       \\ [4, 2]
matsnf([2,4,4;-6,6,12;10,-4,-16])       \\ [12, 6, 2]
matsnf([3,6;9,3])                       \\ [15, 3]
matsnf([6,0,2;0,4,2])                   \\ [2, 2] (transpose of the 3x2 matrix)
\end{lstlisting}
PARI lists the invariant factors in decreasing order.
\end{example}

\begin{theorem}[structure of finitely generated modules over a principal ideal domain]\label{thm:modules:structure-pid}
Let $M$ be a finitely generated $R$-module.
\begin{enumerate}[label=(\roman*)]
\item (Invariant factors.) $M\simeq R^r\oplus R/(d_1)\oplus\cdots\oplus R/(d_s)$ with $d_i$ nonzero non-units and $d_1\mid\cdots\mid d_s$; the integer $r$ and the
ideals $(d_i)$ are uniquely determined.
\item (Elementary divisors.) $M\simeq R^r\oplus\bigoplus_jR/(p_j^{a_j})$ with $p_j$ primes and $a_j\ge1$; the multiset of the ideals $(p_j^{a_j})$ is
uniquely determined.
\item Special cases: for $R=\Z$ this is the structure of finitely generated abelian groups (for finite groups, \Cref{ch:groups});
for a discrete valuation ring, $M\simeq\mathcal O^r\oplus\bigoplus_i\mathcal O/\varpi^{a_i}$ (\Cref{prop:modules:dvr-modules}); for $R=F[x]$ and $M$ a
finite-dimensional vector space with an endomorphism, the rational canonical form, and over an algebraically closed $F$
the Jordan form.
\end{enumerate}
\end{theorem}

\begin{proof}
\emph{Existence.} By \Cref{thm:modules:submodule-free} there is a presentation $R^m\xrightarrow{A}R^n\to M\to0$. Replacing $A$ by $UAV$ changes
the basis of the relations by $U$ and the basis of $R^n$ by $V^{-1}$, so $M\simeq R^n/\text{(row space of }UAV)$. For
$UAV=\operatorname{diag}(d_1,\dots,d_s,0,\dots)$ this is $\bigoplus_{i\le s}R/(d_i)\oplus R^{n-s}$; drop the summands with $d_i$ a unit. (ii) follows from
(i) by the Chinese remainder theorem $R/(d)\simeq\bigoplus_pR/(p^{v_p(d)})$ (\Cref{ch:rings}).
\emph{Uniqueness.} The torsion submodule $M_{\tors}$ (\Cref{def:modules:torsion}) is intrinsic, and $M/M_{\tors}\simeq R^r$, so
$r=\dim_K(M\otimes K)$ is determined. For a prime $p$ the $p$-primary part $M(p)=M_{\tors}[p^\infty]$ is intrinsic and is the sum of
the summands $R/(p^{a_j})$ with $p_j=p$ (up to units). Since $p^kR/(p^a)\big/p^{k+1}R/(p^a)$ is $R/(p)$ if $k<a$ and $0$ otherwise,
$\dim_{R/(p)}\bigl(p^kM(p)/p^{k+1}M(p)\bigr)=\#\{j:p_j=p,\ a_j>k\}$, and these numbers determine the exponents $a_j$ belonging to $p$. So the
elementary divisors are unique, and the invariant factors, which are recovered from them ($d_s$ is the product of the largest prime
powers for each prime, $d_{s-1}$ of the next ones, and so on), are unique too.
\end{proof}

\begin{example}\label{ex:modules:structure}
The abelian group with relation rows $(2,4,4)$, $(-6,6,12)$, $(10,-4,-16)$ is $\Z/2\oplus\Z/6\oplus\Z/12$ by \Cref{ex:modules:smith}; its elementary
divisors are $2,2,4$ and $3,3$, that is, it is $\Z/2\oplus\Z/2\oplus\Z/4\oplus\Z/3\oplus\Z/3$.
\end{example}

\begin{definition}[torsion]\label{def:modules:torsion}
Let $R$ be a domain with fraction field $K$. An element $m\in M$ is \emph{torsion} if $rm=0$ for some $r\ne0$. The torsion elements
form a submodule $M_{\tors}$; $M$ is \emph{torsion-free} if $M_{\tors}=0$ and \emph{torsion} if $M_{\tors}=M$. The \emph{rank} of $M$ is
$\rank M=\dim_K(M\otimes_RK)$.
\end{definition}

\begin{proposition}\label{prop:modules:torsion}
\begin{enumerate}[label=(\roman*)]
\item $M_{\tors}$ is the kernel of $M\to M\otimes_RK$, $m\mapsto m\otimes1$, and $M/M_{\tors}$ is torsion-free.
\item Over a principal ideal domain a finitely generated torsion-free module is free, and $M\simeq M_{\tors}\oplus M/M_{\tors}$ for every
finitely generated $M$ (the splitting is not canonical).
\item Let $R\subseteq R'$ be domains with $R'$ a principal ideal domain that is free of rank $e$ as an $R$-module. A finitely generated
$R'$-module $X$ that is torsion-free over $R$, of $R$-rank $n$, is free over $R'$ of rank $n/e$. For instance, $\Z_3[\zeta_3]$ is free of
rank $2$ over $\Z_3$, with basis $1,\zeta_3$.
\end{enumerate}
\end{proposition}

\begin{proof}
(i) $M\otimes K=S^{-1}M$ with $S=R\setminus\{0\}$ (\Cref{thm:modules:localisation}(ii)), and $m/1=0$ iff $um=0$ for some $u\ne0$. If $rm\in M_{\tors}$
with $r\ne0$, then $r'rm=0$ with $r'\ne0$, so $m\in M_{\tors}$.
(ii) By \Cref{thm:modules:structure-pid}, a finitely generated torsion-free module has no summands $R/(d_i)$; and $M/M_{\tors}$ is free, so
$M\onto M/M_{\tors}$ splits (\Cref{prop:modules:projective}).
(iii) First, $X$ is torsion-free over $R'$. Let $0\ne r'\in R'$ and let $\mu$ be multiplication by $r'$ on $R'\simeq R^e$; it is
injective ($R'$ is a domain), so $\delta=\det\mu\ne0$ (it is injective on $R'\otimes K\simeq K^e$, hence invertible there). By
Cayley-Hamilton, $\mu^e+c_1\mu^{e-1}+\cdots+c_e=0$ with $c_e=\pm\delta$, so $\delta=r'y$ for some $y\in R'$. If $r'x=0$ then $\delta x=0$ with
$0\ne\delta\in R$, so $x=0$. By (ii) over the principal ideal domain $R'$, $X\simeq R'^{\,m}$, which is free of rank $me$ over $R$; so $n=me$.
\end{proof}

\begin{paperlink}\label{pl:modules:IV}
The remark after Lemma lem:IV says: ``$T_3E$ is then free of rank one over $\Z_3[X]/(X^2+X+1)\simeq\Z_3[\zeta_3]$, with $X$ acting as
$\sigma$''. Apply \Cref{prop:modules:torsion}(iii) with $R=\Z_3$, $R'=\Z_3[\zeta_3]$ (a discrete valuation ring, hence a principal ideal
domain, by \Cref{ch:ant}), $e=2$, and $X=T_3E$, which is a free $\Z_3$-module of rank $2$ and an $R'$-module through $\sigma$
(\Cref{ex:modules:running}(c)): it is free of rank $2/2=1$ over $R'$. Lemma lem:tamadd uses the same statement. Exercise
\ref{ex:modules:ex-IV} works out the consequences for $T/3T$.
\end{paperlink}

\begin{definition}[saturation]\label{def:modules:saturated}
Let $M$ be a lattice and $L\subseteq M$ a submodule. $L$ is \emph{saturated} in $M$ if $M/L$ is torsion-free. The \emph{saturation} of $L$ is
$\operatorname{sat}(L)=M\cap KL$, computed inside $M\otimes K$ (where $KL$ is the $K$-span of $L$).
\end{definition}

\begin{proposition}\label{prop:modules:saturated}
Let $M$ be a lattice and $L\subseteq M$ a submodule.
\begin{enumerate}[label=(\roman*)]
\item The following are equivalent: (a) $L$ is saturated; (b) $L=M\cap KL$; (c) $L$ is a direct summand of $M$; and, when
$R=\mathcal O$ is a discrete valuation ring, (d) $L\cap\varpi M=\varpi L$, that is, $L/\varpi L\to M/\varpi M$ is injective.
\item $\operatorname{sat}(L)$ is the smallest saturated submodule containing $L$; it has the same rank as $L$, and $\operatorname{sat}(L)/L=(M/L)_{\tors}$. If
$L\subseteq N\subseteq M$ with $N$ saturated and $N/L$ torsion, then $N=\operatorname{sat}(L)$.
\item The following are saturated: kernels of homomorphisms $M\to N$ with $N$ torsion-free; $M\cap W$ for a $K$-subspace
$W\subseteq M\otimes K$; $M[\mf a]$ for any action of a ring $\T$ on $M$ and any ideal $\mf a\subseteq\T$; and orthogonals for a pairing
(\Cref{prop:modules:orthogonal}). Non-examples: $3\Z\subset\Z$; $\Z(3,3)\subset\Z^2$; while $\Z(1,3)\subset\Z^2$ is saturated.
\end{enumerate}
\end{proposition}

\begin{proof}
(i) (a)$\Rightarrow$(c): $M/L$ is finitely generated and torsion-free, hence free (\Cref{prop:modules:torsion}(ii)), so $M\onto M/L$ splits.
(c)$\Rightarrow$(a): $M/L$ is isomorphic to a complement, a submodule of the torsion-free $M$. (a)$\Leftrightarrow$(b): $x\in M\cap KL$ iff $rx\in L$ for
some $r\ne0$, iff the class of $x$ in $M/L$ is torsion. (a)$\Rightarrow$(d): if $\varpi m\in L$ then $m\in L$ by torsion-freeness of $M/L$.
(d)$\Rightarrow$(a): suppose $rx\in L$ with $r=u\varpi^k$, $u$ a unit; we show $x\in L$ by induction on $k$. If $k\ge1$, then
$\varpi(\varpi^{k-1}x)\in L\cap\varpi M=\varpi L$, so $\varpi(\varpi^{k-1}x-l)=0$ for some $l\in L$, and $\varpi^{k-1}x=l\in L$ because $M$ is torsion-free.
(ii) $M/\operatorname{sat}(L)$ embeds in $(M\otimes K)/KL$, a $K$-vector space, so it is torsion-free. A saturated $N\supseteq L$ satisfies $N=M\cap KN\supseteq M\cap KL$.
$KL$ and $K\operatorname{sat}(L)$ coincide, so the ranks agree, and $\operatorname{sat}(L)/L$ consists of the classes killed by some $r\ne0$. If $N/L$ is torsion, $KN=KL$
and $N=M\cap KN=\operatorname{sat}(L)$.
(iii) $M/\ker f\simeq f(M)\subseteq N$; $M/(M\cap W)$ embeds in $(M\otimes K)/W$; $M/M[\mf a]$ embeds in $\prod_{a\in\mf a}M$ by $x\mapsto(ax)_a$
(the $\T$-action is $R$-linear). In $\Z/3\Z$ and $\Z^2/\Z(3,3)\simeq\Z\oplus\Z/3$ there is torsion; $\Z(1,3)$ is a direct summand with
complement $\Z(0,1)$.
\end{proof}

\begin{theorem}[stacked bases]\label{thm:modules:stacked-bases}
Let $M$ be free of rank $n$ and $L\subseteq M$ a submodule of rank $d$. There are a basis $e_1,\dots,e_n$ of $M$ and nonzero
$d_1\mid\cdots\mid d_d$ such that $d_1e_1,\dots,d_de_d$ is a basis of $L$. The ideals $(d_i)$, the \emph{elementary divisors of $L$ in $M$},
are unique. Hence $M/L\simeq R^{n-d}\oplus\bigoplus_iR/(d_i)$; $L$ is saturated if and only if every $d_i$ is a unit; and if $d=n$ and $B$ is
the matrix whose rows are the coordinates of a basis of $L$ in a basis of $M$, then $M/L$ is torsion, of order $\abs{\det B}$ (the
\emph{index} $[M:L]$) when $R=\Z$, and of length $v(\det B)$ when $R$ is a discrete valuation ring.
\end{theorem}

\begin{proof}
$L$ is free of rank $d$ (\Cref{thm:modules:submodule-free}); let $B\in\Mat_{d\times n}(R)$ be the matrix of a basis of $L$ in a basis $f$ of
$M$. Its rows are independent, so in \Cref{thm:modules:smith}, $UBV=\operatorname{diag}(d_1,\dots,d_d)$ with all $d_i\ne0$. The rows of $UB$ form another
basis of $L$, and in the basis $e=V^{-1}f$ of $M$ (that is, $e_i=\sum_j(V^{-1})_{ij}f_j$) their coordinates are the rows of $UBV$. So
$d_ie_i$ is a basis of $L$. Then $M/L\simeq R^{n-d}\oplus\bigoplus R/(d_i)$, and uniqueness follows from \Cref{thm:modules:structure-pid}
(the number of unit $d_i$ is $d$ minus the number of nonzero torsion summands). The other statements follow, with
$\det B=\det(U^{-1})\det(V^{-1})\prod d_i$ and \Cref{prop:modules:length-dvr}.
\end{proof}

\begin{proposition}[the dual of an inclusion of lattices]\label{prop:modules:dual-inclusion}
Let $L\subseteq M$ be lattices of the same rank $n$, so $M/L$ is torsion. Then restriction $M^\vee\to L^\vee$ is injective and
$L^\vee/M^\vee\simeq\bigoplus_iR/(d_i)$, with the same invariant factors as $M/L$. More precisely, there is a canonical isomorphism
\[
L^\vee/M^\vee\ \xrightarrow{\ \sim\ }\ \Hom_R(M/L,K/R),\qquad\varphi\mapsto\bigl(m+L\mapsto\varphi_K(m)+R\bigr),
\]
where $\varphi_K$ is the $K$-linear extension of $\varphi$ to $M\otimes K=L\otimes K$. For $R=\Z$ or $\Zl$ the right-hand side is the Pontryagin dual
$(M/L)^*$ (\Cref{thm:modules:pontryagin-finite}), so $[L^\vee:M^\vee]=[M:L]$. Equivalently: if $A\into B$ is an injection of lattices of the same
rank with finite cokernel $C$, then $B^\vee\into A^\vee$ with cokernel $\simeq C^*$.
\end{proposition}

\begin{proof}
The displayed map is well defined: $\varphi_K(L)=\varphi(L)\subseteq R$, and if $\varphi\in M^\vee$ the image is $0$. Its kernel is $\{\varphi:\varphi_K(M)\subseteq R\}=M^\vee$.
For surjectivity take stacked bases (\Cref{thm:modules:stacked-bases}): $M=\bigoplus Re_i$, $L=\bigoplus Rd_ie_i$. A homomorphism
$\psi\colon M/L\to K/R$ is determined by the values $\psi(e_i)$, which are killed by $d_i$, so $\psi(e_i)=c_i/d_i+R$ with $c_i\in R$; the functional
$\varphi=\sum_ic_ie_i^\vee/d_i$ is integral on $L$ and maps to $\psi$. In these bases $L^\vee$ has basis $e_i^\vee/d_i$ and $M^\vee$ has basis $e_i^\vee$, which gives
the invariant factors. For a finite group, $\Hom(X,\Q/\Z)$ and, for an $\ell$-primary one, $\Hom(X,\Qlmod)$ are the Pontryagin duals, and
$\#X^*=\#X$ (\Cref{thm:modules:pontryagin-finite}).
\end{proof}

\begin{remark}\label{rem:modules:ext}
In the language of \Cref{ch:cohomology}, $\Hom_R(M/L,K/R)\simeq\Ext^1_R(M/L,R)$, and \Cref{prop:modules:dual-inclusion} is the long exact
$\Ext$ sequence of $0\to L\to M\to M/L\to0$. We do not need this.
\end{remark}

\begin{proposition}[the congruence module of a decomposition]\label{prop:modules:congruence-decomposition}
Let $M$ be a lattice, $M\otimes K=V_1\oplus V_2$ a decomposition into $K$-subspaces with projections $p_1,p_2$, and $L_i=M\cap V_i$ (saturated).
\begin{enumerate}[label=(\roman*)]
\item $L_1\oplus L_2\subseteq M$ has torsion quotient, and $x\mapsto p_i(x)$ induces isomorphisms $M/(L_1\oplus L_2)\simeq p_1(M)/L_1\simeq p_2(M)/L_2$.
\item Suppose $V_1=Kf$ with $f\in M$ and $L_1=Rf$ (that is, $f$ is primitive). Then $M/(L_1\oplus L_2)$ embeds in $Kf/Rf\simeq K/R$, so it is cyclic,
$\simeq R/(r)$, generated by the class of any $h\in M$ with $p_1(h)=f/r$; then $rh=f+g_0$ with $g_0\in L_2$, so $f\equiv-g_0\pmod{rM}$. For
$0\ne r'\in R$ there is $g\in L_2$ with $f\equiv g\pmod{r'M}$ if and only if $r'\mid r$. So $(r)$ is the \emph{congruence modulus} of $f$
against $L_2$: the largest modulus of a congruence between $f$ and an element of $L_2$.
\item Example: $M=\Z^2$, $f=(1,a)$, $g=(1,b)$ with $a\ne b$, $V_1=\Q f$, $V_2=\Q g$. Then $L_1=\Z f$, $L_2=\Z g$, $M/(\Z f\oplus\Z g)\simeq\Z/(a-b)$, and
$f\equiv g\pmod{(a-b)}$ coordinatewise.
\end{enumerate}
\end{proposition}

\begin{proof}
(i) Every $v\in V_1$ has a nonzero multiple in $M$ (which spans $M\otimes K$), hence in $L_1$; so $L_1$ spans $V_1$, likewise $L_2$, and the
quotient has rank $0$. The map $M\to p_1(M)/L_1$ is onto; $x$ lies in its kernel iff $p_1(x)\in L_1\subseteq M$, and then
$p_2(x)=x-p_1(x)\in M\cap V_2=L_2$; so the kernel is $L_1\oplus L_2$. Symmetrically for $p_2$.
(ii) $p_1(M)$ is a finitely generated submodule of $Kf$ containing $Rf$, hence free of rank one, $p_1(M)=R(f/r)$ for some $r\ne0$ (a generator
$u$ satisfies $f=su$, so $u=f/s$). By (i) the quotient is $R(f/r)/Rf\simeq R/(r)$, generated by any $h$ with $p_1(h)=f/r$. Then
$rh-f=rp_2(h)\in M\cap V_2=L_2$. If $f-g=r'm$ with $g\in L_2$ and $m\in M$, then $p_1(m)=f/r'\in R(f/r)$, so $r=sr'$; conversely
$f\equiv-g_0\pmod{rM}$ implies the congruence modulo every divisor $r'$ of $r$.
(iii) $f$ and $g$ are primitive, so $L_1=\Z f$ and $L_2=\Z g$, and the index of $\Z f\oplus\Z g$ is $\abs{\det\Mtwo1a1b}=\abs{b-a}$ (\Cref{thm:modules:stacked-bases});
the quotient is cyclic by (ii), and $f-g=(0,a-b)$.
\end{proof}

\begin{paperlink}\label{pl:modules:crit-second}
The second paragraph of the proof of Lemma lem:crit is \Cref{prop:modules:dual-inclusion,prop:modules:congruence-decomposition} with
$M=S_2(\Gamma_0(N);\Z)$, $V_1=\Q f$, $V_2=f^\perp$, $L_1=\Z f$ (as $a_1(f)=1$, $f$ is primitive) and $L_2=L=f^\perp\cap S_2(\Gamma_0(N);\Z)$.
``Dualising the inclusion $\T\subset\T/I\times\T/J$, whose cokernel is $\T/(I+J)$, shows that $\Z f\oplus L$ has index $r'$'' is
\Cref{prop:modules:algebra-duality}(iii), which rests on \Cref{prop:modules:dual-inclusion}. ``The quotient embeds in $\Q f/\Z f$ by
projection onto $\Q f$ along $f^\perp$, so it is cyclic, generated by the class of a form $h$ with projection $f/r'$; then $r'h=f+g_0$ with
$g_0\in L$'' is \Cref{prop:modules:congruence-decomposition}(ii), whose last clause is the paper's ``$r'\le r$'' and ``$r\mid r'$'' in one
statement: the congruence number is the congruence modulus. Lemma lem:ddt417 uses the same structure, with
$L\otimes K=L[\wp]_K\oplus L[I]_K$ and $Q=L/(L[\wp]\oplus L[I])$.
\end{paperlink}

\section[Discrete valuation rings and \texorpdfstring{$\Zl$}{Zl}]{Discrete valuation rings, \texorpdfstring{$\Zl$}{Zl}, \texorpdfstring{$\Ql$}{Ql} and \texorpdfstring{$\Qlmod$}{Ql/Zl}}\label{sec:modules:dvr}

The coefficient rings of the paper are discrete valuation rings: $\Zl$ everywhere, $\mathcal O$ in Lemma lem:ddt417 and Section 4.1,
$\Z_3[\zeta_3]$ in Lemma lem:IV. Over such a ring finitely generated modules are classified by a rank and a list of exponents, which is
what makes lengths and Fitting ideals computable (\Cref{sec:modules:length}). The module $\Qlmod$ carries the discrete coefficients
$W=\ad T\otimes\Qlmod$ of the paper's Selmer groups, and its injectivity makes Pontryagin duality exact (\Cref{sec:modules:pontryagin}).
Here $\Zl$ is built algebraically; the analytic theory is \Cref{ch:local} (\Cref{warn:modules:analytic}).

\begin{theorem}[discrete valuation rings]\label{thm:modules:dvr}
Let $\mathcal O$ be a Noetherian local domain that is not a field, with maximal ideal $\mf m$ and fraction field $K$. The following are equivalent:
\begin{enumerate}[label=(\alph*)]
\item $\mf m$ is principal;
\item $\mathcal O$ is a principal ideal domain;
\item every nonzero ideal is a power of $\mf m$;
\item there is a surjective homomorphism $v\colon K^\times\to\Z$ with $v(x+y)\ge\min(v(x),v(y))$ and $\mathcal O=\{x\in K^\times:v(x)\ge0\}\cup\{0\}$;
\item $\mathcal O$ is integrally closed in $K$ (and of dimension one).
\end{enumerate}
Such a ring is a \emph{discrete valuation ring}. A generator $\varpi$ of $\mf m$ is a \emph{uniformiser}; every $x\in K^\times$ is uniquely
$x=u\varpi^{v(x)}$ with $u\in\mathcal O^\times$; $\mathcal O^\times=\{v=0\}$; and the nonzero ideals are the $\varpi^k\mathcal O$, $k\ge0$.
\end{theorem}

\begin{proof}
(a)$\Rightarrow$(c),(d),(b). Let $\mf m=\varpi\mathcal O$. First, $\bigcap_k\mf m^k=0$: if $x=\varpi^ky_k$ for every $k$, then $\varpi^ky_k=\varpi^{k+1}y_{k+1}$ gives
$y_k=\varpi y_{k+1}$ ($\mathcal O$ is a domain), so $(y_0)\subseteq(y_1)\subseteq\cdots$; this chain stabilises, $y_{k+1}=uy_k$ for some $k$ and $u\in\mathcal O$, so
$(1-\varpi u)y_k=0$, and $1-\varpi u$ is a unit (\Cref{prop:modules:local-ring}); hence $y_k=0$ and $x=0$. So every $x\ne0$ lies in $\mf m^k$ but not
$\mf m^{k+1}$ for a unique $k$; then $x=\varpi^ku$ with $u\notin\mf m$, a unit. The exponent is unique: $\varpi^ku=\varpi^ju'$ with $k<j$ would give
$u=\varpi^{j-k}u'\in\mf m$. Put $v(x)=k$ and $v(x/y)=v(x)-v(y)$; this is well defined and multiplicative, and
$v(x+y)\ge\min(v(x),v(y))$ because $\varpi^{\min}$ divides both. A nonzero ideal $\mf a$ contains an element of least valuation $k$, and then
$\mf a=\varpi^k\mathcal O$. This proves (c), (d) and (b).
(b)$\Rightarrow$(a): trivial. (c)$\Rightarrow$(a): $\mf m\ne\mf m^2$, since otherwise $\mf m=0$ by Nakayama ($\mf m$ is finitely generated); pick
$\varpi\in\mf m\setminus\mf m^2$; then $\varpi\mathcal O=\mf m^k$ with $k=1$. (d)$\Rightarrow$(a): pick $\varpi$ with $v(\varpi)=1$. If $x\in\mf m$ then $v(x)\ge1$ (if
$v(x)=0$ then $v(x^{-1})=0$ and $x^{-1}\in\mathcal O$), so $x/\varpi\in\mathcal O$.
(e): we do not prove the equivalence with (e); see \cite[Proposition 9.2]{AtiyahMacdonald} (number not verified against the printed
edition) or \cite[Chapter I, Section 2]{SerreLocalFields}. It is used only in \Cref{ch:ant}.
The last statements were proved along the way.
\end{proof}

A discrete valuation ring $\mathcal O$ is \emph{complete} if the map $\mathcal O\to\varprojlim_n\mathcal O/\varpi^n\mathcal O$ is an isomorphism. Then every
finitely generated $\mathcal O$-module $M$, in particular every free one of finite rank, is complete: $M\to\varprojlim_nM/\varpi^nM$ is an
isomorphism, by the proof of \Cref{thm:modules:zl}(vi), which uses only completeness of the ring. $\Zl$ is complete (\Cref{thm:modules:zl}(vi));
$\Z_{(\ell)}$ is not, since $\varprojlim\Z_{(\ell)}/\ell^n=\varprojlim\Z/\ell^n=\Zl$ is uncountable.

\begin{example}\label{ex:modules:dvr-examples}
Discrete valuation rings: $\Z_{(p)}$ (\Cref{ch:rings}); $F[[x]]$ for a field $F$ (every nonzero series is $x^ku$ with $u$ a unit, since a
series with nonzero constant term is invertible, so the ideals are the $(x^k)$); $\Zl$ (\Cref{thm:modules:zl}); and, with proofs in
\Cref{ch:ant,ch:local}, $\Z_3[\zeta_3]$, $\Z_\ell[\zeta_3]$ for $\ell\equiv2\pmod3$, and the completions $\mathcal O_\lambda$ of rings of integers.
Not discrete valuation rings: $\Z$ (not local); $\Z/p^2$ and $\mathcal O[x]/(x^2)$ (not domains); $\Zl\times\Zl$ (not local); $\Zl[x]$ (not local:
$(\ell,x)$ and $(x-1,\ell)$ are distinct maximal ideals); $\Z_{(11)}[x]/(x^2+x-1)$ (a domain with two maximal ideals,
\Cref{thm:modules:idempotent-lifting}(v)).
\end{example}

\begin{proposition}[modules over a discrete valuation ring]\label{prop:modules:dvr-modules}
Let $M$ be a finitely generated $\mathcal O$-module.
\begin{enumerate}[label=(\roman*)]
\item $M\simeq\mathcal O^r\oplus\bigoplus_{i=1}^s\mathcal O/\varpi^{a_i}$ with $a_i\ge1$; $r$ and the multiset $\{a_i\}$ are unique.
\item If $M$ is torsion-free, it is free.
\item $\dim_kM/\varpi M=r+s$ is the minimal number of generators, and $\dim_kM[\varpi]=s$.
\item If $r=0$ then $M[\varpi^k]\simeq\bigoplus_i\mathcal O/\varpi^{\min(k,a_i)}$; in general $M/\varpi^kM\simeq(\mathcal O/\varpi^k)^r\oplus\bigoplus_i\mathcal O/\varpi^{\min(k,a_i)}$.
\item $M$ is cyclic if and only if $r+s\le1$; a torsion $M$ with $\dim_kM[\varpi]\le1$ is cyclic.
\item A submodule $L$ of a lattice $M$ is saturated if and only if $L\cap\varpi M=\varpi L$.
\end{enumerate}
\end{proposition}

\begin{proof}
(i) is \Cref{thm:modules:structure-pid}, as $\varpi$ is the only prime up to units; (ii) is \Cref{prop:modules:torsion}(ii). For (iii) and (iv)
compute summand by summand. For the summand $\mathcal O$: $\mathcal O/\varpi\mathcal O=k$ and $\mathcal O[\varpi]=0$. For $\mathcal O/\varpi^a$: the quotient by $\varpi$ is $k$, the
$\varpi$-torsion is $\varpi^{a-1}\mathcal O/\varpi^a\simeq k$, the $\varpi^k$-torsion is $\varpi^{\max(a-k,0)}\mathcal O/\varpi^a\simeq\mathcal O/\varpi^{\min(k,a)}$, and the quotient by $\varpi^k$
is $\mathcal O/\varpi^{\min(k,a)}$. The minimal
number of generators is \Cref{thm:modules:nakayama}(iii). (v) follows from (iii). (vi) is \Cref{prop:modules:saturated}(i)(d).
\end{proof}

\begin{theorem}[the ring $\Zl$]\label{thm:modules:zl}
Let $\Zl=\varprojlim_n\Z/\ell^n$, the ring of compatible sequences $x=(x_n)_{n\ge1}$, $x_n\in\Z/\ell^n$, $x_{n+1}\equiv x_n\pmod{\ell^n}$, with
componentwise operations (as a profinite group it is constructed in \Cref{ch:groups}).
\begin{enumerate}[label=(\roman*)]
\item $x$ is a unit if and only if $x_1\ne0$ in $\Fl$.
\item Every nonzero $x$ is uniquely $\ell^ku$ with $k\ge0$ and $u\in\Zl^\times$. Hence $\Zl$ is a discrete valuation ring with uniformiser $\ell$,
residue field $\Fl$ (via $x\mapsto x_1$) and ideals $\ell^k\Zl$, and its valuation extends $v_\ell$ on $\Z$.
\item $\Z\to\Zl$, $n\mapsto(n\bmod\ell^k)_k$, is injective; $\Z_{(\ell)}\subseteq\Zl$; and $x\mapsto x_n$ induces $\Zl/\ell^n\Zl\simeq\Z/\ell^n$. An integer prime to
$\ell$ is a unit of $\Zl$, so a $\Zl$-module killed by an integer prime to $\ell$ is $0$.
\item $\Ql:=\Frac(\Zl)=\Zl[1/\ell]$, and $\Q\cap\Zl=\Z_{(\ell)}$ inside $\Ql$.
\item $\Zl$ is flat over $\Z$. For a finitely generated abelian group $A\simeq\Z^r\oplus A_{\tors}$, $A\otimes_\Z\Zl\simeq\Zl^r\oplus A_{\tors}[\ell^\infty]$; for finite
$A$, $\#(A\otimes\Zl)=\ell^{v_\ell(\#A)}$, the $\ell$-part of $\#A$.
\item Every finitely generated $\Zl$-module $M$ is complete: $M\to\varprojlim_nM/\ell^nM$ is an isomorphism.
\item More generally, for a Noetherian ring $R$ and an ideal $I$, the $I$-adic completion $\hat R=\varprojlim R/I^n$ is Noetherian and
flat over $R$, and $\hat M=\varprojlim M/I^nM\simeq M\otimes_R\hat R$ for finitely generated $M$.
\end{enumerate}
\end{theorem}

\begin{proof}
(i) If $x_1\ne0$, then $\ell\nmid x_n$ for all $n$, so each $x_n$ is a unit of $\Z/\ell^n$, and the inverses are compatible (inverses are unique);
conversely $x_1$ is the image of $x$ in the field $\Fl$.
(ii) Let $x\ne0$ and let $k$ be the largest index with $x_k=0$ (with $x_0=0$ by convention); it exists because $x_n=0$ forces $x_j=0$ for
$j\le n$. For $n\ge1$, $x_{n+k}\equiv x_k=0\pmod{\ell^k}$, so $x_{n+k}=\ell^ky_n$ with $y_n\in\Z/\ell^n$ well defined; the $y_n$ are compatible, $y_1\ne0$ because
$x_{k+1}\ne0$, so $u=(y_n)$ is a unit with $x=\ell^ku$. Then $k=v(x)$ is determined by $x$. Since $\ell^k\ne0$ (its component at $k+1$ is
nonzero), the product $\ell^{a+b}uu'$ of two nonzero elements is nonzero: $\Zl$ is a domain. A nonzero ideal is $\ell^k\Zl$ for its least
valuation $k$, so $\Zl$ is a principal ideal domain, local with maximal ideal $\ell\Zl$ (the non-units, by (i)): a discrete valuation ring by
\Cref{thm:modules:dvr}. The kernel of $x\mapsto x_1$ is $\{v\ge1\}=\ell\Zl$.
(iii) $n\ne0$ has nonzero image in $\Z/\ell^k$ for $\ell^k>\abs n$. By (ii), $x_n=0$ iff $v(x)\ge n$, so the kernel of $x\mapsto x_n$ is $\ell^n\Zl$; it is
onto because $\Z$ maps onto $\Z/\ell^n$. Integers prime to $\ell$ are units by (i), which gives $\Z_{(\ell)}\subseteq\Zl$.
(iv) $x/y=x\ell^{-v(y)}u^{-1}$ for $y=\ell^{v(y)}u$. By (iii), $\Q\subseteq\Ql$ and $\Z_{(\ell)}\subseteq\Zl$. Conversely, if $a/b$ is in lowest terms
with $\ell\mid b$, then $\ell\nmid a$ and $v(a/b)=-v_\ell(b)<0$ by (ii), so $a/b\notin\Zl$.
(v) $\Zl$ is torsion-free over $\Z$, hence flat (\Cref{prop:modules:flat}(ii)). $\Z/n\otimes\Zl=\Zl/n\Zl=\Zl/\ell^{v_\ell(n)}\Zl\simeq\Z/\ell^{v_\ell(n)}$ by
(ii), (iii); sum over the cyclic summands of $A$.
(vi) For $M=\Zl$ the map $\Zl\to\varprojlim\Zl/\ell^n\Zl=\varprojlim\Z/\ell^n=\Zl$ is the identity; for $M=\Z/\ell^a$ the system is eventually constant;
and both sides commute with finite direct sums. Conclude by \Cref{prop:modules:dvr-modules}(i).
(vii) We do not prove this; see \cite[Chapter 10]{AtiyahMacdonald} (the Artin-Rees lemma) or \cite[Chapter 7]{Eisenbud}.
\end{proof}

\begin{lemma}[Baer's criterion over a principal ideal domain]\label{lem:modules:baer}
Let $R$ be a principal ideal domain and $D$ a \emph{divisible} $R$-module ($dD=D$ for every $d\ne0$). Then $D$ is \emph{injective}: for every
submodule $N\subseteq M$, every homomorphism $N\to D$ extends to $M$. Consequently $\Hom_R(-,D)$ is exact.
\end{lemma}

\begin{proof}
Let $f\colon N\to D$. By Zorn's lemma there is a maximal extension $f_0\colon N_0\to D$ ($N\subseteq N_0\subseteq M$), the union of a chain of
extensions being an extension. Suppose $m\in M\setminus N_0$. The ideal $\{r:rm\in N_0\}$ is $dR$. If $d=0$, extend by $f(n+rm)=f_0(n)$. If
$d\ne0$, choose $y\in D$ with $dy=f_0(dm)$ and put $f(n+rm)=f_0(n)+ry$. This is well defined: if $n+rm=0$ then $rm\in N_0$, so $r=sd$, and
$f_0(n)+ry=f_0(n)+sf_0(dm)=f_0(n+rm)=0$. This contradicts maximality, so $N_0=M$. Exactness of $\Hom(-,D)$: by \Cref{prop:modules:hom}(ii) only
surjectivity of $\Hom(M,D)\to\Hom(N,D)$ for $N\subseteq M$ remains, and that is injectivity of $D$.
\end{proof}

\begin{proposition}[the module $\Qlmod$]\label{prop:modules:qlzl}
\begin{enumerate}[label=(\roman*)]
\item $\Qlmod\simeq\Z[1/\ell]/\Z\simeq(\Q/\Z)[\ell^\infty]\simeq\varinjlim_k\Z/\ell^k$, the transition maps being multiplication by $\ell$. As an abstract group
it is isomorphic to $\mu_{\ell^\infty}$ (no Galois action is involved here).
\item It is divisible and $\ell$-primary torsion; $(\Qlmod)[\ell^k]=\ell^{-k}\Zl/\Zl\simeq\Z/\ell^k$; its proper $\Zl$-submodules are exactly these.
\item The ring $\End_{\Zl}(\Qlmod)$ is $\Zl$, acting by multiplication; moreover $\Hom_{\Zl}(\Z/\ell^k,\Qlmod)\simeq\Z/\ell^k$, and
$\Hom_{\Zl}(\Zl,\Qlmod)=\Qlmod$.
\item $\Qlmod$ is an injective $\Zl$-module and $\Q/\Z$ an injective $\Z$-module; so $\Hom(-,\Qlmod)$ and $\Hom(-,\Q/\Z)$ are exact.
\end{enumerate}
\end{proposition}

\begin{proof}
(i) Every element of $\Ql$ is $x=\ell^{-k}u$ with $u\in\Zl$; choosing $a\in\Z$ with $u\equiv a\pmod{\ell^k}$ gives $x\equiv a/\ell^k$ modulo $\Zl$. So
$\Z[1/\ell]\to\Qlmod$ is onto, with kernel $\Z[1/\ell]\cap\Zl=\Z$. A class $a/b+\Z$ in reduced form is killed by a power of $\ell$ iff $b$ is a power of
$\ell$, which identifies $\Z[1/\ell]/\Z$ with $(\Q/\Z)[\ell^\infty]$; and $\Z/\ell^k\simeq\ell^{-k}\Z/\Z$. The map $\Z[1/\ell]/\Z\to\mu_{\ell^\infty}$,
$x\mapsto e^{2\pi ix}$, is an isomorphism.
(ii) Divisibility: $x=\ell^{-k}u$ is $\ell^ju'$ times $\ell^{-k-j}u'^{-1}u$ for any $j$ and unit $u'$. The $\ell^k$-torsion is $\ell^{-k}\Zl/\Zl$, cyclic of order
$\ell^k$. A submodule containing elements of every order contains every $(\Qlmod)[\ell^k]$, since an element of exact order $\ell^k$
generates this cyclic group; otherwise it lies in some $(\Qlmod)[\ell^k]$ and is one of its subgroups $(\Qlmod)[\ell^j]$.
(iii) An endomorphism $\phi$ preserves each $(\Qlmod)[\ell^k]\simeq\Z/\ell^k$, where it is multiplication by some $a_k\in\Z/\ell^k$; the $a_k$ are compatible,
so $\phi$ is multiplication by $a=(a_k)\in\Zl$, and $a\ne0$ acts nontrivially. The other two are \Cref{prop:modules:hom}(i).
(iv) Both are divisible (for $\Q/\Z$ over $\Z$, $x/n$ divides $x$), so \Cref{lem:modules:baer} applies.
\end{proof}

\begin{proposition}[the topology of finitely generated $\Zl$-modules]\label{prop:modules:zl-topology}
Give every finitely generated $\Zl$-module $M$ the $\ell$-adic topology, in which the $\ell^nM$ form a basis of neighbourhoods of $0$.
\begin{enumerate}[label=(\roman*)]
\item $M$ is compact, Hausdorff and complete.
\item Every $\Zl$-linear map between finitely generated modules is continuous; every submodule of $M$ is finitely generated and
closed, and every linear map has closed image.
\item $\GL_n(\Zl)=\Aut_{\Zl}(\Zl^n)=\varprojlim_k\GL_n(\Z/\ell^k)$ is profinite (\Cref{ch:groups}).
\item (Topological Nakayama.) A compact Hausdorff $\Zl$-module $A$ on which $\Zl$ acts continuously and with $A/\ell A$ finite is finitely
generated.
\end{enumerate}
\end{proposition}

\begin{proof}
(i) $M\simeq\Zl^r\oplus F$ with $F$ finite, and $\ell^nM$ corresponds to $(\ell^n\Zl)^r\oplus\ell^nF$, so the $\ell$-adic topology is the product topology.
$\Zl=\varprojlim\Z/\ell^n$ is a closed subset of the compact product $\prod\Z/\ell^n$ of finite discrete sets, and its subspace topology has the
basis $\ell^n\Zl$ (the kernels of the projections). Completeness is \Cref{thm:modules:zl}(vi).
(ii) $f(\ell^nM)\subseteq\ell^nN$, so $f$ is continuous. A submodule $N$ is finitely generated (\Cref{prop:modules:noetherian}), hence the image of a
continuous map $\Zl^s\to M$, hence compact, hence closed. The same applies to $f(M)$.
(iii) A matrix over $\Zl$ is invertible iff it is invertible modulo $\ell$ (\Cref{thm:modules:nakayama}(iv)), and $\Zl^{n\times n}=\varprojlim(\Z/\ell^k)^{n\times n}$.
\end{proof}

\begin{proof}[Sketch of proof of (iv)]
By Pontryagin duality for compact and discrete modules, which is \Cref{thm:modules:pontryagin-zl}(v), the dual $A^*$ is a discrete $\ell$-primary module with
$A^*[\ell]=(A/\ell A)^*$ finite, hence cofinitely generated, so $A=A^{**}$ is finitely generated by \Cref{thm:modules:pontryagin-zl}(ii). The
duality theory for compact modules is in \cite[Chapter I, Section 1, and Chapter V, Section 2]{NSW} and \cite[Appendix B]{Rubin}
(section numbers not verified against the printed editions).
\end{proof}

\begin{warning}\label{warn:modules:analytic}
\Cref{thm:modules:zl} constructs $\Zl$ as a ring and proves only what is needed for modules. The absolute value, the description of $\Q_p$ as
the completion of $\Q$, Hensel's lemma and the structure of $\Z_p^\times$ and $\Q_p^\times$ are proved in \Cref{ch:local}. Idempotent lifting
(\Cref{thm:modules:idempotent-lifting}) is proved below by a Newton iteration that uses only completeness, not Hensel's lemma.
\end{warning}

\begin{paperlink}\label{pl:modules:killed-by-2}
After Lemma lem:dfg21 the paper says: ``At $\infty$, $H^1(\R,T^D)$ is killed by $2$, so $H^1_f(\R,T^D)=0$ for $\ell$ odd''. The group
$H^1(\R,T^D)$ is a $\Zl$-module (\Cref{ch:cohomology}) killed by the integer $2$, which is a unit of $\Zl$ for $\ell$ odd
(\Cref{thm:modules:zl}(iii)); so it is $0$, and so is its submodule $H^1_f(\R,T^D)$. In Lemma lem:fpr, ``a nonzero rational number which is an
$\ell$-adic integer'' means an element of $\Q\cap\Zl=\Z_{(\ell)}$ (\Cref{thm:modules:zl}(iv)).
\end{paperlink}

\section{Length, annihilators and Fitting ideals}\label{sec:modules:length}

The paper measures its finite modules in two ways. Selmer groups and $H^0$ are measured by their length over $\Zl$
(Theorem thm:main, Theorem thm:dfg37, Remark rem:ns, Lemma lem:h0). Tamagawa factors are \emph{Fitting ideals}: Section 1.2 defines
$\Tam_p(M)=\Fitt_{\Zl}\bigl((H^1(I_p,M)_{\tors})^{G_{\Q_p}}\bigr)$, and Lemma lem:fpr reads $\Fitt_{\Zl}H^1_f(\Q_p,T^D)=L_p(A_f,1)^{-1}\Tam_p(T^D)$. The main
result of this section, \Cref{thm:modules:fitting-dvr}, says that over a discrete valuation ring the two measurements agree:
$\Fitt=\varpi^{\len}$. It is the reason why the paper can add valuations of Tamagawa factors to lengths of Selmer groups. The
Jordan-Holder theorem is proved here for modules over any ring; \Cref{ch:reps} uses it for composition factors of representations.

\begin{definition}[composition series and length]\label{def:modules:composition}
In this definition and in \Cref{thm:modules:jordan-holder}, $R$ is any ring with $1$ (not necessarily commutative) and modules are
left modules. A module is \emph{simple} if it is nonzero and has no submodules other than $0$ and itself. A \emph{composition series}
of $M$ is a chain $0=M_0\subsetneq M_1\subsetneq\cdots\subsetneq M_n=M$ with every $M_i/M_{i-1}$ simple; $n$ is its \emph{length} and the $M_i/M_{i-1}$ are its
\emph{factors}. The \emph{length} $\len_R(M)$ is the length of any composition series (well defined by \Cref{thm:modules:jordan-holder}), and
$\len_R(M)=\infty$ if $M$ has none.
\end{definition}

\begin{theorem}[Jordan-Holder; length]\label{thm:modules:jordan-holder}
Let $R$ be any ring.
\begin{enumerate}[label=(\roman*)]
\item If $M$ has a composition series, then every chain of submodules of $M$ is finite and can be refined to a composition series, and
any two composition series of $M$ have the same length and the same factors up to order and isomorphism (Jordan-Holder).
Submodules and quotients of $M$ have composition series.
\item A module has finite length if and only if it is Noetherian and Artinian (descending chains stabilise).
\item (Additivity.) For an exact sequence $0\to M'\to M\to M''\to0$, $\len(M)=\len(M')+\len(M'')$ (in $\Z_{\ge0}\cup\{\infty\}$). For an exact sequence
$0\to M_1\to\cdots\to M_k\to0$ of modules of finite length, $\sum_i(-1)^i\len(M_i)=0$.
\end{enumerate}
Now let $R$ be commutative.
\begin{enumerate}[label=(\roman*),resume]
\item Let $R$ be local, with maximal ideal $\mf m$ and residue field $k$. Every simple module is isomorphic to $k$. If $\mf m^NM=0$, then
$\len_R(M)=\sum_{i<N}\dim_k\mf m^iM/\mf m^{i+1}M$. If $R$ is Noetherian local, a finitely generated module has finite length if and only if it is
killed by a power of $\mf m$.
\item If $R\to R'$ is a local homomorphism of local rings ($\mf m_R$ maps into $\mf m_{R'}$), $R'$ is finitely generated as an $R$-module and
$f=[k':k]$ is the degree of the residue field extension, then $\len_R(M)=f\cdot\len_{R'}(M)$ for every $R'$-module $M$ of finite length.
\end{enumerate}
\end{theorem}

\begin{proof}
(i) For a module $X$ with a composition series let $\lambda(X)$ be the least length of a composition series of $X$.
\emph{Step 1: if $N\subsetneq M$ then $N$ has a composition series and $\lambda(N)<\lambda(M)$.} Let $(M_i)_{i\le n}$ be a composition series with
$n=\lambda(M)$ and put $N_i=N\cap M_i$. Then $N_i/N_{i-1}=N_i/(N_i\cap M_{i-1})\simeq(N_i+M_{i-1})/M_{i-1}\subseteq M_i/M_{i-1}$ (\Cref{thm:modules:isomorphism}(ii),
which holds verbatim over any ring), so it is $0$ or simple. Deleting repetitions gives a composition series of $N$ of length at
most $n$. If the length were $n$, every $N_i/N_{i-1}$ would be nonzero, hence equal to $M_i/M_{i-1}$ inside it; by induction on $i$,
$N_{i-1}=M_{i-1}$ and $N_i+M_{i-1}=M_i$ give $N_i=M_i$, so $N=M$, a contradiction.
\emph{Step 2: lengths.} A strictly decreasing chain $M=K_0\supsetneq K_1\supsetneq\cdots\supsetneq K_j$ gives $\lambda(M)>\lambda(K_1)>\cdots>\lambda(K_j)\ge0$ by
Step 1, so $j\le\lambda(M)$. Hence every chain is finite of length at most $\lambda(M)$, every composition series has length exactly
$\lambda(M)$ (it is at most $\lambda(M)$ and at least $\lambda(M)$ by minimality), and a chain can be refined by inserting terms until
all quotients are simple, which must stop.
\emph{Step 3: factors.} Induction on $n=\lambda(M)$, the case $n\le1$ being clear. Let $(M_i)$ and $(M'_i)$ be composition series, both of length $n$
by Step 2. If $M_{n-1}=M'_{n-1}$, apply the induction hypothesis to $M_{n-1}$, of length $n-1$. Otherwise $M_{n-1}+M'_{n-1}=M$ (it strictly
contains $M_{n-1}$, and $M/M_{n-1}$ is simple). Put $K=M_{n-1}\cap M'_{n-1}$. By \Cref{thm:modules:isomorphism}(ii),
$M/M_{n-1}\simeq M'_{n-1}/K$ and $M/M'_{n-1}\simeq M_{n-1}/K$, both simple. $K$ has a composition series $(K_j)$ (Step 1), of length $n-2$, because
$(K_j)$ followed by $M_{n-1}$ is a composition series of $M_{n-1}$, of length $n-1$. By induction, the factors of $(M_i)_{i\le n-1}$ are those of
$(K_j)$ together with $M_{n-1}/K\simeq M/M'_{n-1}$; so the factors of $(M_i)$ are those of $(K_j)$ together with $M/M'_{n-1}$ and $M/M_{n-1}$.
The same computation for $(M'_i)$ gives the factors of $(K_j)$ together with $M'_{n-1}/K\simeq M/M_{n-1}$ and $M/M'_{n-1}$: the same list.
Quotients: the image in $M/N$ of a composition series of $M$ has successive quotients that are quotients of simple modules, so $0$ or simple.
(ii) We do not prove this; see \cite[Proposition 6.8]{AtiyahMacdonald} (number not verified). We use only the direction ``finite
length implies Noetherian'', which is Step 2.
(iii) If $M$ has finite length so do $M'$ and $M''$ by (i); conversely, the preimage in $M$ of a composition series of $M''$, placed after a
composition series of $M'$, is a composition series of $M$. Its length is $\len M'+\len M''$. For the long sequence, cut it into short exact
sequences $0\to Z_i\to M_i\to Z_{i+1}\to0$ with $Z_i$ the images, and add.
(iv) A simple module is $Rs$ for any $s\ne0$, so $\simeq R/\Ann(s)$ with $\Ann(s)$ maximal (\Cref{prop:modules:module-basics}), hence $\simeq k$.
A module killed by $\mf m$ is a $k$-vector space whose $R$-submodules are its subspaces, so its length is its dimension; apply (iii) to
the filtration $M\supseteq\mf mM\supseteq\cdots\supseteq\mf m^NM=0$. If $R$ is Noetherian and $M$ finitely generated, the $\mf m^iM/\mf m^{i+1}M$ are finitely
generated over $k$, which gives one direction. Conversely if $\len M=n$, the chain $(\mf m^iM)_i$ has at most $n$ strict steps, so
$\mf m^jM=\mf m^{j+1}M$ for some $j$, and $\mf m^jM=0$ by Nakayama.
(v) By (iii) along an $R'$-composition series it suffices to treat the simple $R'$-module $k'$. It is killed by $\mf m_R$, so by (iv)
its $R$-length is $\dim_kk'=f$, which is finite as $k'$ is a quotient of the finitely generated $R$-module $R'$.
\end{proof}

\begin{proposition}[length over a discrete valuation ring]\label{prop:modules:length-dvr}
Let $\mathcal O$ be a discrete valuation ring.
\begin{enumerate}[label=(\roman*)]
\item $\len(\mathcal O/\varpi^a)=a$. A finitely generated $\mathcal O$-module has finite length if and only if it is torsion, and then
$\len\bigl(\bigoplus_i\mathcal O/\varpi^{a_i}\bigr)=\sum_ia_i$.
\item If $k$ is finite, $\#M=(\#k)^{\len M}$. Over $\Zl$, $\len_{\Zl}(M)=v_\ell(\#M)$ for finite $M$; over a ring $\mathcal O_\lambda$ with residue field $\F_{\ell^f}$
(finite over $\Zl$), $v_\ell(\#M)=f\len_{\mathcal O_\lambda}(M)$.
\item For a finite-length $\mathcal O$-module $M$ and $k\ge0$, $\len M[\varpi^k]=\len M/\varpi^kM$.
\end{enumerate}
For example, $\len_{\Zl}(\Z/\ell^v\oplus\Z/\ell^m)=v+m$ and $\len_{\Zl}(\Zl/n'\Zl)=v_\ell(n')$ for $n'\in\Zl\setminus\{0\}$.
\end{proposition}

\begin{proof}
(i) $0\subsetneq\varpi^{a-1}\mathcal O/\varpi^a\subsetneq\cdots\subsetneq\mathcal O/\varpi^a$ has factors $\simeq k$. $\mathcal O$ itself has the infinite chain $\varpi^i\mathcal O$, so it has infinite
length (\Cref{thm:modules:jordan-holder}(i)); now use \Cref{prop:modules:dvr-modules}(i) and additivity. (ii) Orders are multiplicative in
exact sequences and $\#k$ is the order of each factor; for $\mathcal O_\lambda$ use \Cref{thm:modules:jordan-holder}(v). (iii) The sequence
$0\to M[\varpi^k]\to M\xrightarrow{\varpi^k}M\to M/\varpi^kM\to0$ is exact; apply \Cref{thm:modules:jordan-holder}(iii).
\end{proof}

\begin{definition}[annihilators]\label{def:modules:annihilator}
$\Ann_R(M)=\{r\in R:rM=0\}$, an ideal; $\Ann_R(m)$ was defined in \Cref{prop:modules:module-basics}; for an ideal $\mf b$ of a ring $\T$,
$\Ann_\T(\mf b)$ is the annihilator of $\mf b$ viewed as a $\T$-module. ``$M$ is killed by $\eta$'' means $\eta\subseteq\Ann(M)$.
\end{definition}

\begin{proposition}\label{prop:modules:annihilator}
\begin{enumerate}[label=(\roman*)]
\item $\Ann(R/\mf a)=\mf a$; $\Ann(M\oplus N)=\Ann(M)\cap\Ann(N)$; for generators $m_1,\dots,m_n$, $\Ann(M)$ is the kernel of $R\to M^n$, $r\mapsto(rm_i)$.
\item $M[\mf a]\simeq\Hom_R(R/\mf a,M)$; $M[\mf a]\cap M[\mf b]=M[\mf a+\mf b]$; $M[\mf a]+M[\mf b]\subseteq M[\mf a\cap\mf b]$.
\item For a ring $\T$ and an ideal $\mf a$, $\T[\mf a]=\Ann_\T(\mf a)$; if $L\simeq\T^r$ as $\T$-modules, then $L[\mf a]\simeq\Ann_\T(\mf a)^r$.
\item For a flat $R\to S$, finitely generated $M$ and finitely generated $\mf a$: $\Ann_S(M\otimes S)=\Ann_R(M)S$ and $M[\mf a]\otimes S=(M\otimes S)[\mf aS]$.
\item Over a discrete valuation ring, $\Ann\bigl(\bigoplus_{i=1}^s\mathcal O/\varpi^{a_i}\bigr)=\varpi^{\max_ia_i}\mathcal O$ (for $s\ge1$), and $\Ann(M)=0$ if $M$ has a free summand.
\end{enumerate}
\end{proposition}

\begin{proof}
(i), (ii), (iii) are immediate from the definitions and \Cref{prop:modules:hom}(i) (componentwise for $\T^r$). (iv) is
\Cref{prop:modules:flat}(iii). (v) $\varpi^a$ kills $\mathcal O/\varpi^b$ iff $a\ge b$; and $\mathcal O$ is torsion-free.
\end{proof}

\begin{definition}[Fitting ideals]\label{def:modules:fitting}
Let $M$ be a finitely generated $R$-module with generators $x=(x_1,\dots,x_n)$, and let $\operatorname{Rel}(x)=\{r\in R^n:\sum r_ix_i=0\}$ be the module of
relations (row vectors). The (zeroth) \emph{Fitting ideal} $\Fitt_R(M)$ is the ideal generated by the determinants $\det(r^{(1)},\dots,r^{(n)})$ of all
$n$-tuples of relations (the $n\times n$ matrices with these rows). If $R^m\xrightarrow{A}R^n\to M\to0$ is a presentation (\Cref{conv:modules:rings}(g)),
$\Fitt_R(M)$ is the ideal generated by the $n\times n$ minors of $A$ (and is $0$ if $m<n$); see \Cref{lem:modules:fitting-well-defined}. More generally
$\Fitt_i(M)$ is generated by the $(n-i)\times(n-i)$ minors; we use only $\Fitt=\Fitt_0$.
\end{definition}

\begin{lemma}[the Fitting ideal is well defined]\label{lem:modules:fitting-well-defined}
\begin{enumerate}[label=(\roman*)]
\item If the rows of $A\in\Mat_{m\times n}(R)$ generate $\operatorname{Rel}(x)$, the ideal of $n\times n$ minors of $A$ equals the ideal $F(x)$ generated by the
determinants of $n$-tuples of relations.
\item $F(x)$ does not depend on the generating family $x$.
\end{enumerate}
\end{lemma}

\begin{proof}
(i) Every $n\times n$ minor of $A$ is the determinant of $n$ relations. Conversely, if $r^{(i)}=\sum_jc_{ij}a_j$ with $a_j$ the rows of $A$, then by
multilinearity
\[\det(r^{(1)},\dots,r^{(n)})=\sum_{j_1,\dots,j_n}c_{1j_1}\cdots c_{nj_n}\det(a_{j_1},\dots,a_{j_n}),\]
and each term is $0$ (two equal rows) or $\pm$ a minor of $A$.
(ii) Permuting the generators permutes the columns of the relation matrices and changes determinants by a sign. Next let
$x'=(x_1,\dots,x_n,y)$ with $y=\sum_ia_ix_i$. If the rows of $A$ generate $\operatorname{Rel}(x)$, then $\operatorname{Rel}(x')$ is generated by the rows $(r,0)$, $r$ a row of
$A$, and the row $(-a_1,\dots,-a_n,1)$: given $(s,t)\in\operatorname{Rel}(x')$, subtracting $t$ times the last row leaves $(s+ta,0)$ with $s+ta\in\operatorname{Rel}(x)$. An
$(n+1)\times(n+1)$ minor of this matrix that omits the last row has a zero last column; one that contains it equals, expanding along the
last column, $\pm$ an $n\times n$ minor of $A$, and every such minor occurs. By (i), $F(x')=F(x)$. Finally, for two generating families $x$
and $z=(z_1,\dots,z_p)$, add the $z_j$ to $x$ one at a time (each is a combination of the $x_i$, hence of the current family), and the $x_i$
to $z$ likewise; both processes end at the same family up to order. So $F(x)=F(x\cup z)=F(z)$.
\end{proof}

\begin{proposition}[properties of Fitting ideals]\label{prop:modules:fitting-properties}
Let $M$, $N$ be finitely generated.
\begin{enumerate}[label=(\roman*)]
\item $\Fitt(R/\mf a)=\mf a$; $\Fitt(0)=R$; $\Fitt(R^n)=0$ for $n\ge1$.
\item $\Fitt(M\oplus N)=\Fitt(M)\Fitt(N)$.
\item If $M\onto N$ then $\Fitt(M)\subseteq\Fitt(N)$.
\item $\Fitt(M)\subseteq\Ann(M)$, and $\Ann(M)^n\subseteq\Fitt(M)$ if $M$ is generated by $n$ elements.
\item $\Fitt_S(S\otimes_RM)=\Fitt_R(M)S$ for every ring homomorphism $R\to S$ (no flatness is needed).
\item For an exact sequence $0\to M'\to M\to M''\to0$ with $M'$ finitely generated (automatic if $R$ is Noetherian,
\Cref{prop:modules:noetherian}(ii)), $\Fitt(M')\Fitt(M'')\subseteq\Fitt(M)$. Equality can fail: over $R=\Z[x]$, with $M=R/(4,x)\simeq\Z/4$,
$M'=2M\simeq R/(2,x)$ and $M''=M/2M\simeq R/(2,x)$, one has $\Fitt(M)=(4,x)$ but $\Fitt(M')\Fitt(M'')=(2,x)^2=(4,2x,x^2)$.
\item Over a domain, $\Fitt(M)\ne0$ if and only if $M$ is torsion.
\end{enumerate}
\end{proposition}

\begin{proof}
(i) $R/\mf a$ has the generator $1$ with relations $\mf a$, and its $1\times1$ ``determinants'' generate $\mf a$. For $M=0$ use the generator $0$ with
the relation $(1)$. $R^n$ with its basis has no nonzero relation.
(ii) Use generators $x$ of $M$ and $y$ of $N$; $\operatorname{Rel}(x,y)=\operatorname{Rel}(x)\oplus\operatorname{Rel}(y)$ is generated by the rows of $\operatorname{diag}(A,B)$, with $A$ of size $m\times n$ and $B$ of
size $m'\times n'$. Consider $n+n'$ of these rows, $j$ from the first block and $j'$ from the second. If $j>n$, the Laplace expansion along the first $j$
rows involves only $j\times j$ minors of these rows that use one of the last $n'$ columns, which vanish there; so the determinant is $0$, and
likewise if $j'>n'$. If $j=n$ and $j'=n'$ the determinant is the product of an $n\times n$ minor of $A$ and an $n'\times n'$ minor of $B$. So the ideal is
$\Fitt(M)\Fitt(N)$.
(iii) Generators of $M$ map to generators of $N$, and relations among them in $M$ remain relations in $N$.
(iv) Let $B$ be an $n\times n$ matrix of relations, so $B$ kills the column $x$ of generators; multiplying by the adjugate gives
$\det(B)x_i=0$ for all $i$. For $a_1,\dots,a_n\in\Ann(M)$ the rows $a_ie_i$ are relations, with determinant $a_1\cdots a_n$.
(v) Tensoring the presentation gives a presentation of $S\otimes M$ with the image of the same matrix (\Cref{prop:modules:right-exact}).
(vi) Let $x'$ generate $M'$ (viewed in $M$) and let $y\in M^{n''}$ lift generators of $M''$. For $r'\in\operatorname{Rel}(x')$, $(r',0)$ is a relation of $(x',y)$; for
$r''\in\operatorname{Rel}(\bar y)$, $\sum r''_jy_j\in M'$, say $=\sum c_ix'_i$, and $(-c,r'')$ is a relation. The determinant of $n'$ relations of the first kind followed by $n''$
of the second is block triangular, equal to $\det(r'\text{s})\det(r''\text{s})$. In the example: $M'=2\Z/4$ is killed by $2$ and $x$, so it is $R/(2,x)$, as is
$M/2M$; $\Fitt(R/(4,x))=(4,x)$ by (i); $(2,x)^2=(4,2x,x^2)$; and $x\notin(4,2x,x^2)$ (every element of that ideal has the coefficient of $x$ even).
(vii) If $M$ is torsion, $\Ann(M)\ne0$ (the product of nonzero annihilators of the generators), so $\Fitt(M)\supseteq\Ann(M)^n\ne0$. If not, $M\otimes K\simeq K^r$
with $r\ge1$ (\Cref{prop:modules:torsion}(i)), and $\Fitt(M)K=\Fitt_K(K^r)=0$ by (i) and (v).
\end{proof}

\begin{theorem}[Fitting ideals and length]\label{thm:modules:fitting-dvr}
\begin{enumerate}[label=(\roman*)]
\item Over a principal ideal domain, the module $M=R^r\oplus\bigoplus_iR/(d_i)$ has $\Fitt(M)=0$ if $r>0$, and $\Fitt(M)=(\prod_id_i)$ if $r=0$.
\item Over a discrete valuation ring $\mathcal O$, every finitely generated torsion module $M$ has $\Fitt_{\mathcal O}(M)=\varpi^{\len(M)}\mathcal O$, so $v(\Fitt M)=\len M$.
Over $\Zl$: $\Fitt_{\Zl}(M)=\#M\cdot\Zl=\ell^{v_\ell(\#M)}\Zl$ for finite $M$.
\item Consequently, over a discrete valuation ring Fitting ideals of finite-length modules are multiplicative in short exact sequences,
and $\Fitt(N)\supseteq\Fitt(M)$ for every submodule or quotient $N$ of $M$.
\item The Fitting ideal is not the annihilator: $(\Z/\ell)^2$ and $\Z/\ell^2$ both have Fitting ideal $\ell^2\Zl$, with annihilators $\ell\Zl$ and $\ell^2\Zl$.
\end{enumerate}
\end{theorem}

\begin{proof}
(i) By \Cref{prop:modules:fitting-properties}(i), (ii). (ii) By (i), $\Fitt(\bigoplus\mathcal O/\varpi^{a_i})=\varpi^{\sum a_i}\mathcal O$, and $\sum a_i=\len M$
(\Cref{prop:modules:length-dvr}). Over $\Zl$, $\#M=\ell^{\len M}$. (iii) Lengths are additive (\Cref{thm:modules:jordan-holder}(iii)), and submodules
and quotients have smaller length. (iv) By (ii) and \Cref{prop:modules:annihilator}(v).
\end{proof}

\begin{warning}\label{warn:modules:fitt-ann}
Over $\Zl$ the Fitting ideal of a finite module sees its whole order, while the annihilator sees only its exponent. The Tamagawa factor is a
Fitting ideal; reading it as an annihilator would replace $v_\ell(n)+\min(\dots)$ in Lemma lem:tamfactor by $\max(v_\ell(n),\min(\dots))=v_\ell(n)$. Over a general ring the Fitting ideal is not determined by lengths either
(\Cref{prop:modules:fitting-properties}(vi)).
\end{warning}

\begin{example}\label{ex:modules:fitting-examples}
Over $\Z$ (relations as rows), and then over $\Z_3$ and $\Zl$; all entries were computed by hand and the Smith forms checked with PARI
(\Cref{ex:modules:smith}).
\begin{center}\small
\begin{tabular}{llll}
\toprule
relations & module & $\Fitt$ & $\Ann$\\
\midrule
rows $(2,4)$, $(6,8)$ & $\Z/2\oplus\Z/4$ & $(8)$ & $(4)$\\
rows $(2,4,4)$, $(-6,6,12)$, $(10,-4,-16)$ & $\Z/2\oplus\Z/6\oplus\Z/12$ & $(144)$ & $(12)$\\
rows $(3,6)$, $(9,3)$ & $\Z/3\oplus\Z/15$ & $(45)$ & $(15)$\\
the same over $\Z_3$ & $\Z/3\oplus\Z/3$ & $9\Z_3$ & $3\Z_3$\\
rows $(6,0)$, $(0,4)$, $(2,2)$ & $\Z/2\oplus\Z/2$ & $(4)$ & $(2)$\\
one generator, rows $(2)$, $(4)$ & $\Z/2$ & $(2)$ & $(2)$\\
over $\Z[x]$: one generator, rows $(2)$, $(x)$ & $\Z[x]/(2,x)$ & $(2,x)$ & $(2,x)$\\
over $\Z[x]$: the square of the previous & $(\Z[x]/(2,x))^2$ & $(4,2x,x^2)$ & $(2,x)$\\
over $\Zl$: the paper's module & $\Z/\ell^v\oplus\Z/\ell^m$ & $\ell^{v+m}\Zl$ & $\ell^{\max(v,m)}\Zl$\\
$\ell=3$, $v=2$, $m=1$ & $\Z/9\oplus\Z/3$ & $27\Z_3$ & $9\Z_3$\\
\bottomrule
\end{tabular}
\end{center}
In the fifth row the three $2\times2$ minors are $24$, $12$, $-8$, with gcd $4$, although there are more relations than generators. In the
eighth row the four relation rows $(2,0)$, $(x,0)$, $(0,2)$, $(0,x)$ have $2\times2$ minors $0,4,2x,2x,x^2,0$, and the Fitting ideal $(2,x)^2$ is strictly
smaller than the annihilator, as \Cref{prop:modules:fitting-properties}(iv) allows.
\end{example}

\begin{paperlink}\label{pl:modules:tamagawa}
Section 1.2 defines $\Tam_p(M)=\Fitt_{\Zl}\bigl((H^1(I_p,M)_{\tors})^{G_{\Q_p}}\bigr)$, the Fitting ideal \emph{of the invariants} (Appendix B, convention
B-C10). This is a finite $\Zl$-module, so by
\Cref{thm:modules:fitting-dvr}(ii) $v_\ell(\Tam_p(M))$ is its length. In Lemma lem:tamfactor the invariants are the $x\bar e+y\bar h$ in
$(\Z/\ell^{v_\ell(n)})\bar e\oplus(\Z/\ell^{v_\ell(n)})\bar h$ with $(q-1)y\equiv0$ and $c'y\equiv0$, a module isomorphic to $\Z/\ell^{v_\ell(n)}\oplus\Z/\ell^{\min(v_\ell(n),v_\ell(q-1),v_\ell(c'))}$, of
length $v_\ell(n)+\min(v_\ell(n),v_\ell(q-1),v_\ell(c'))$: this is the paper's formula, and the Fitting ideal of $\Z/\ell^v\oplus\Z/\ell^m$ is
$\ell^{v+m}\Zl$ by the last rows of \Cref{ex:modules:fitting-examples}. In Lemma lem:fpr the identity of Fitting ideals is an identity of
lengths, $\len H^1_f(\Q_p,T^D)=v_\ell(L_p(A_f,1)^{-1})+v_\ell(\Tam_p(T^D))$, which is how the proof of Theorem thm:main uses it. The definitions of
$\Tam_p$ and $H^1_f$ are \Cref{ch:blochkato}.
\end{paperlink}

\section{Pontryagin duality}\label{sec:modules:pontryagin}

Lemma lem:dfg21 is an exact sequence
\[
0\to H^1_f(\Q,W)\to H^1_\Sigma(\Q,W)\to\bigoplus_{p\in\Sigma\cup\{\infty\}}H^1_f(\Q_p,T^D)^*\to H^0(\Q,W^D)^*\to0
\]
in which ${}^*$ is the Pontryagin dual, and the proof of Theorem thm:main adds lengths along it. This needs two facts proved here: Pontryagin
duality is exact, and it preserves the length of finite modules. The coefficient module $W=\ad T\otimes\Qlmod$ is the basic example of a
discrete module dual to a finitely generated one (\Cref{prop:modules:t-and-w}), and Lemma lem:h0 identifies $H^0(\Q,W^D)$ by its
$\ell^k$-torsion, using $W[\ell^k]=T/\ell^kT$ and the cyclicity criterion \Cref{cor:modules:cyclic-l-primary}.

\begin{theorem}[Pontryagin duality for finite abelian groups]\label{thm:modules:pontryagin-finite}
For a finite abelian group $A$ let $A^*=\Hom(A,\Q/\Z)$.
\begin{enumerate}[label=(\roman*)]
\item $(\Z/n)^*\simeq\Z/n$, and $A^*\simeq A$ (non-canonically).
\item The evaluation map $A\to A^{**}$, $a\mapsto(\xi\mapsto\xi(a))$, is an isomorphism.
\item $A\mapsto A^*$ is an exact contravariant functor, and $\#A^*=\#A$.
\item For a subgroup $B\subseteq A$ let $B^\perp=\{\xi\in A^*:\xi(B)=0\}$. Then $B^\perp\simeq(A/B)^*$ and $B^*\simeq A^*/B^\perp$, and $B\mapsto B^\perp$ is an
inclusion-reversing bijection between the subgroups of $A$ and those of $A^*$.
\item If $A$ is $\ell$-primary, then $A^*=\Hom_{\Zl}(A,\Qlmod)$ and $\len_{\Zl}A^*=\len_{\Zl}A$; dualising an exact sequence of finite $\Zl$-modules
gives an exact sequence with the arrows reversed.
\item Via $x\mapsto e^{2\pi ix}$, $A^*$ is the group of characters $A\to\C^\times$ (\Cref{ch:rings}); in the language of Riley, Hobson and Bence,
Pontryagin duality is the discrete Fourier transform on $A$.
\end{enumerate}
\end{theorem}

\begin{proof}
(i) $\Hom(\Z/n,\Q/\Z)\simeq(\Q/\Z)[n]=\frac1n\Z/\Z\simeq\Z/n$ by \Cref{prop:modules:hom}(i); and $A$ is a finite direct sum of cyclic groups
(\Cref{ch:groups}), and $\Hom$ turns finite sums into sums. (ii) Injectivity: for $a\ne0$ of order $d$, the map $\langle a\rangle\to\Q/\Z$, $a\mapsto1/d$,
extends to $A$ because $\Q/\Z$ is injective (\Cref{prop:modules:qlzl}(iv)); so some $\xi$ has $\xi(a)\ne0$. Both sides have order $\#A$ by (i).
(iii) Exactness is the injectivity of $\Q/\Z$, and $\#A^*=\#A$ is (i). (iv) The exact sequence $0\to B\to A\to A/B\to0$ dualises to
$0\to(A/B)^*\to A^*\to B^*\to0$, with the image of the first map equal to $B^\perp$. Hence $\#B^\perp=\#A/\#B$. Under $A=A^{**}$, $B\subseteq(B^\perp)^\perp$ and
$\#(B^\perp)^\perp=\#A^*/\#B^\perp=\#B$, so $B=(B^\perp)^\perp$; and symmetrically for subgroups of $A^*$.
(v) $A$ is killed by some $\ell^k$, so every homomorphism lands in $(\Q/\Z)[\ell^k]\subseteq(\Q/\Z)[\ell^\infty]=\Qlmod$ (\Cref{prop:modules:qlzl}(i)); it is
$\Zl$-linear because $\Zl$ acts on a module killed by $\ell^k$ through $\Z/\ell^k$. Lengths agree because orders do
(\Cref{prop:modules:length-dvr}(ii)). (vi) $x\mapsto e^{2\pi ix}$ is an isomorphism from $\Q/\Z$ onto the roots of unity in $\C^\times$.
\end{proof}

\begin{theorem}[Pontryagin duality for $\Zl$-modules]\label{thm:modules:pontryagin-zl}
Let $\mathcal C$ be the category of finitely generated $\Zl$-modules and $\mathcal D$ that of $\ell$-primary torsion $\Zl$-modules $W$ with $W[\ell]$ finite
(\emph{cofinitely generated} modules), with the discrete topology. For $X$ in either, put $X^*=\Hom_{\Zl}(X,\Qlmod)$. (For $M\in\mathcal C$ every such
homomorphism is continuous for the $\ell$-adic topology, since its image is a finitely generated, hence finite, submodule of $\Qlmod$, killed by
some $\ell^k$; so the homomorphism kills the open submodule $\ell^kM$.) For $W\in\mathcal D$, $W^*=\varprojlim_k(W[\ell^k])^*$ along the restriction
maps, and we give it this inverse-limit topology, which is profinite; by (ii) below $W^*$ is finitely generated, and the topology is its
$\ell$-adic one, because $\ell^kW^*$ kills $W[\ell^k]$ and a continuous bijection from a compact space onto a Hausdorff space is a homeomorphism.
\begin{enumerate}[label=(\roman*)]
\item $\Zl^*=\Qlmod$, $(\Qlmod)^*=\Zl$, and $(\Z/\ell^a)^*\simeq\Z/\ell^a$.
\item $X\mapsto X^*$ is exact, maps $\mathcal C$ to $\mathcal D$ and $\mathcal D$ to $\mathcal C$, and the evaluation maps $M\to M^{**}$, $W\to W^{**}$ are isomorphisms.
\item Every $W\in\mathcal D$ is isomorphic to $(\Qlmod)^r\oplus F$ with $F$ finite; $r$, the \emph{corank} of $W$, equals $\rank_{\Zl}W^*$. $W$ is finite iff
$r=0$, and divisible iff $F=0$.
\item Lengths of finite modules are preserved: $\len X^*=\len X$.
\item More generally, $X\mapsto\Hom_{\mathrm{cont}}(X,\Qlmod)$ is an exact anti-equivalence between compact and discrete $\ell$-primary $\Zl$-modules (stated;
see \cite[Chapter I, Section 1]{NSW}).
\end{enumerate}
\end{theorem}

\begin{proof}
(iii) first. By induction $W[\ell^{k+1}]$ is finite: multiplication by $\ell$ maps it to $W[\ell^k]$ with kernel $W[\ell]$. Let $D=\bigcap_k\ell^kW$.
\emph{$D$ is divisible.} Let $x\in D$; for each $k$ write $x=\ell^{k+1}w_k$ and put $y_k=\ell^kw_k$, so $\ell y_k=x$ and $y_k\in\ell^kW$. The $y_k$ all lie in
the finite set $y_0+W[\ell]$, so one value $y$ occurs for infinitely many $k$; then $y\in\ell^kW$ for all $k$, so $y\in D$ and $\ell y=x$. Divisibility by
units is automatic.
\emph{$D\simeq(\Qlmod)^r$, $r=\dim_{\Fl}D[\ell]$.} Choose an isomorphism $(\Qlmod)^r[\ell]\simeq D[\ell]$ and extend it to $\phi\colon(\Qlmod)^r\to D$ using the
injectivity of $D$ (\Cref{lem:modules:baer}). The kernel of $\phi$ has no $\ell$-torsion, and an $\ell$-primary module without $\ell$-torsion is $0$; so $\phi$ is
injective. Its image is divisible, hence injective, hence a direct summand: $D=\phi((\Qlmod)^r)\oplus C$, where $C[\ell]=0$ because $D[\ell]$ lies in the
image; so $C=0$.
\emph{$W=D\oplus F$ with $F$ finite.} $D$ is injective, so $W=D\oplus F$. Then $\ell^kW=D\oplus\ell^kF$, so $\bigcap_k\ell^kF=0$. The subgroups
$\ell^kF\cap F[\ell]$ of the finite group $F[\ell]$ decrease and have intersection $0$, so $\ell^kF\cap F[\ell]=0$ for some $k$; then $\ell^kF$ has no
$\ell$-torsion, so $\ell^kF=0$, and $F=F[\ell^k]$ is finite. $W$ is finite iff $r=0$; it is divisible iff $F=0$ (a finite divisible $\ell$-primary module is $0$).
(i) $\Hom(\Zl,X)=X$; $\End(\Qlmod)=\Zl$ and $\Hom(\Z/\ell^a,\Qlmod)\simeq\Z/\ell^a$ (\Cref{prop:modules:qlzl}(iii)).
(ii) Exactness: $\Qlmod$ is injective (\Cref{prop:modules:qlzl}(iv)). By \Cref{prop:modules:dvr-modules}(i) and (iii) above, every object of $\mathcal C$ or
$\mathcal D$ is a finite direct sum of copies of $\Zl$, $\Qlmod$ and $\Z/\ell^a$, and by (i) duality exchanges $\Zl$ and $\Qlmod$ and preserves $\Z/\ell^a$. The
evaluation map is natural and compatible with finite sums, so it suffices to check it on these three modules: for $\Zl$, identify $\Zl^*=\Qlmod$ by
$\xi\mapsto\xi(1)$; then the evaluation at $a\in\Zl$ is multiplication by $a$ on $\Qlmod$, which is the element $a$ of $\End(\Qlmod)=\Zl$. For $\Qlmod$, identify
$(\Qlmod)^*=\Zl$; the evaluation at $x$ is $a\mapsto ax$, which is the element $x$ of $\Hom(\Zl,\Qlmod)=\Qlmod$. For $\Z/\ell^a$ it is
\Cref{thm:modules:pontryagin-finite}(ii). Finally $\rank_{\Zl}W^*=r$ by the decomposition.
(iv) is \Cref{thm:modules:pontryagin-finite}(v). (v) We do not prove this.
\end{proof}

\begin{proposition}[the duality between $T$ and $T\otimes\Qlmod$]\label{prop:modules:t-and-w}
Let $T$ be a free $\Zl$-module of rank $r$, $V=T\otimes_{\Zl}\Ql$ and $W=T\otimes_{\Zl}\Qlmod$.
\begin{enumerate}[label=(\roman*)]
\item $0\to T\to V\to W\to0$ is exact, and $W\simeq(\Qlmod)^r$ is divisible of corank $r$.
\item $W[\ell^k]\simeq T/\ell^kT$, by $t\otimes\ell^{-k}\mapsto t\bmod\ell^k$. Under these isomorphisms the inclusion $W[\ell^k]\subseteq W[\ell^{k+1}]$ becomes multiplication by
$\ell\colon T/\ell^k\to T/\ell^{k+1}$, and multiplication by $\ell\colon W[\ell^{k+1}]\to W[\ell^k]$ becomes reduction $T/\ell^{k+1}\to T/\ell^k$.
\item $T\simeq\varprojlim_kW[\ell^k]$, the transition maps being multiplication by $\ell$: $T$ is the \emph{Tate module} of $W$.
\item $W^*=\Hom(T\otimes\Qlmod,\Qlmod)\simeq\Hom_{\Zl}(T,\End(\Qlmod))=\Hom_{\Zl}(T,\Zl)=T^\vee$, and $T^*\simeq T^\vee\otimes\Qlmod$.
\item Everything is functorial in $T$: a group acting $\Zl$-linearly on $T$ acts on $V$, $W$, $W[\ell^k]=T/\ell^kT$ compatibly with the maps above, and on
duals by $(g\xi)(x)=\xi(g^{-1}x)$.
\item The same holds over any complete discrete valuation ring $\mathcal O$ with fraction field $K$, with $K/\mathcal O$ in place of $\Qlmod$.
\end{enumerate}
\end{proposition}

\begin{proof}
(i) Tensor the exact sequence $0\to\Zl\to\Ql\to\Qlmod\to0$ with the free, hence flat, module $T$. (ii) Apply \Cref{cor:modules:snake}(a) with $n=\ell^k$:
$V[\ell^k]=0$ and $V/\ell^kV=0$, so the connecting map $W[\ell^k]\to T/\ell^kT$ is an isomorphism. It sends the class of $v\in V$ (with $\ell^kv\in T$) to
$\ell^kv\bmod\ell^kT$, which on $t\otimes\ell^{-k}$ is $t\bmod\ell^k$. If $w$ is the class of $v$ with $\ell^kv=t$, then in $W[\ell^{k+1}]$ it corresponds to
$\ell^{k+1}v=\ell t$; and $\ell w$, the class of $\ell v$, corresponds in $W[\ell^k]$ to $\ell^k(\ell v)=\ell^{k+1}v$ reduced modulo $\ell^k$.
(iii) By (ii) the system $(W[\ell^k],\cdot\ell)$ is $(T/\ell^kT,\text{reduction})$, whose limit is $T$ (\Cref{thm:modules:zl}(vi)).
(iv) The adjunction (\Cref{thm:modules:tensor}(iii)) and \Cref{prop:modules:qlzl}(iii). For the second, the canonical map
$T^\vee\otimes X\to\Hom(T,X)$, $\varphi\otimes x\mapsto(t\mapsto\varphi(t)x)$, is an isomorphism for $T$ free of finite rank (check on a basis).
(v) All maps above are natural in $T$. (vi) The proofs use only that $\mathcal O$ is a discrete valuation ring, that $K/\mathcal O$ is divisible (hence injective,
\Cref{lem:modules:baer}), and, for (iii) and (iv), completeness of $\mathcal O$: $\End_{\mathcal O}(K/\mathcal O)=\varprojlim\mathcal O/\varpi^k=\mathcal O$ as in
\Cref{prop:modules:qlzl}(iii).
\end{proof}

\begin{corollary}\label{cor:modules:cyclic-l-primary}
An $\ell$-primary torsion $\Zl$-module $A$ with $\#A[\ell]\le\ell$ is $0$, $\Z/\ell^h$ for some $h\ge1$, or $\Qlmod$. If it is finite, its length is the largest
$k$ such that $A$ has an element of exact order $\ell^k$.
\end{corollary}

\begin{proof}
$A\in\mathcal D$, so $A\simeq(\Qlmod)^r\oplus F$ by \Cref{thm:modules:pontryagin-zl}(iii), and $\dim_{\Fl}A[\ell]=r+\dim F[\ell]\le1$. If $r=1$ then $F[\ell]=0$, so $F=0$.
If $r=0$, $F$ is finite with $\dim F[\ell]\le1$, hence cyclic (\Cref{prop:modules:dvr-modules}(v)), $F\simeq\Z/\ell^h$, of length $h$, and an element of exact
order $\ell^k$ exists iff $k\le h$.
\end{proof}

\begin{paperlink}\label{pl:modules:dfg21}
In the proof of Theorem thm:main the paper takes lengths in the sequence of Lemma lem:dfg21. All four terms are finite there:
$H^1_\Sigma(\Q,W)$ by Theorem thm:dfg37, the local terms $H^1_f(\Q_p,T^D)$ as torsion submodules of finitely generated $\Zl$-modules (the
paragraph after Lemma lem:dfg21), and $H^0(\Q,W^D)$ by Lemma lem:h0. By \Cref{thm:modules:pontryagin-zl}(iv) the duals have the same lengths, and by
\Cref{thm:modules:jordan-holder}(iii)
\[
\len H^1_\Sigma(\Q,W)=\len H^1_f(\Q,W)+\sum_{p\in\Sigma\cup\{\infty\}}\len H^1_f(\Q_p,T^D)-\len H^0(\Q,W^D),
\]
which is the paper's display once $H^0(\Q,W^D)=0$ and $H^1_f(\R,T^D)=0$ (\Cref{pl:modules:killed-by-2}); Remark rem:ns keeps the term
$\len H^0(\Q,W^D)=h_\ell(E)$. In Lemma lem:h0, $W^D=\ad T(1)\otimes\Qlmod$, so $W^D[\ell^k]\simeq\ad T(1)/\ell^k$ by \Cref{prop:modules:t-and-w}(ii)
(functorially, hence as Galois modules), which is the paper's ``$\ell^k$-torsion of $H^0(\Q,W^D)$ is $H^0(\Q,\ad(T/\ell^k)(1))$''. The $\ell$-torsion of
$H^0(\Q,W^D)$ has $\Fl$-dimension at most one, so \Cref{cor:modules:cyclic-l-primary} shows that $H^0(\Q,W^D)$ is $0$, cyclic of order $\ell^h$, or
$\Qlmod$; the last is excluded by the determinant argument of the lemma, and $h_\ell(E)$ is the largest $k$ for which an element of exact order $\ell^k$
exists.
\end{paperlink}

\section[Pairings]{Pairings: Gram determinants, saturated sublattices, alternating pairings, adjoints}\label{sec:modules:pairings}

The heart of Lemma lem:ddt417 is a statement about a perfect pairing on a lattice $L$ and a saturated sublattice $L[\wp]$: ``the pairing
identifies $L/L[\wp]^\perp$ with $\Hom(L[\wp],\mathcal O)$, because $L[\wp]$ is saturated'', and ``$Q$ is the cokernel of $L[\wp]\to\Hom(L[\wp],\mathcal O)$, of
order $\#\mathcal O/\Delta$''. Section 4.1 of the paper adds that the pairing $\hat\delta_!$ is perfect and alternating, that the Hecke operators are
self-adjoint for it, and that on the rank-two lattice $M_{f^\Sigma,\lambda}$ ``the Gram determinant is $(\eta^\Sigma_{f,\lambda})^2$ up to a unit''. Lemma lem:crit
uses the pairing $S_2(\Gamma_0(N);\Z)\times\T\to\Z$. This section proves each of these statements in the abstract. Throughout, except in \Cref{prop:modules:gram}, which holds over any commutative ring, $R$ is a principal
ideal domain (in practice $\Z$, $\Zl$ or $\mathcal O$) with fraction field $K$, and pairings follow \Cref{conv:modules:rings}(e). To a reader of Riley,
Hobson and Bence a Gram matrix is a metric tensor $g_{ij}=\langle e_i,e_j\rangle$; the novelty is that over $\Z$ or $\Zl$ the determinant is meaningful
only up to units, and its valuation is what is measured.

\begin{definition}[perfect pairings]\label{def:modules:pairing}
For a pairing $\langle\ ,\ \rangle\colon L\times M\to R$ let $\alpha\colon L\to M^\vee$, $x\mapsto\langle x,-\rangle$, and $\beta\colon M\to L^\vee$, $y\mapsto\langle-,y\rangle$. The pairing is
\emph{perfect} if $\alpha$ and $\beta$ are isomorphisms and \emph{nondegenerate} if they are injective. For free modules $L$, $M$ of the same finite rank
with bases $(e_i)$, $(f_j)$, the \emph{Gram determinant} is $\Delta=\det(\langle e_i,f_j\rangle)$.
\end{definition}

\begin{proposition}\label{prop:modules:gram}
Let $R$ be any commutative ring, $L$, $M$ free $R$-modules of rank $n$ and $G$ the Gram matrix of a pairing $L\times M\to R$.
\begin{enumerate}[label=(\roman*)]
\item The pairing is perfect iff $\det G\in R^\times$, iff one of $\alpha$, $\beta$ is an isomorphism. If $R$ is a domain, it is nondegenerate iff
$\det G\ne0$. For example,
$\langle x,y\rangle=2xy$ on $\Z$ is nondegenerate and not perfect; over $\Zl$ with $\ell$ odd it is perfect.
\item If the bases are changed by $e'=eP$, $f'=fQ$ (that is, $e'_j=\sum_iP_{ij}e_i$), then $G'=P^{\mathsf T}GQ$ and $\det G'=\det P\det Q\det G$; for $L=M$ with the
same change of basis on both sides, $\det G'=(\det P)^2\det G$. So the ideal $\Delta R$ does not depend on the bases; over a discrete valuation ring
$v(\Delta)$ is a basis-independent invariant.
\item Perfectness is preserved by base change: if $R\to S$ is a ring homomorphism, the $S$-bilinear extension to
$(L\otimes S)\times(M\otimes S)$ has Gram matrix the image of $G$, so it is perfect if the original is.
\end{enumerate}
\end{proposition}

\begin{proof}
(i) In the dual basis $(f_j^\vee)$, $\alpha(e_i)=\sum_j\langle e_i,f_j\rangle f_j^\vee$, so $\alpha$ has matrix $G$ and $\beta$ has matrix $G^{\mathsf T}$, both of determinant $\det G$.
\Cref{prop:modules:free}(iii), valid over any commutative ring, gives the first claim; over a domain $R$ with fraction field $K$, a square matrix is injective iff its determinant is nonzero (compute over
$K$). (ii) $\langle e'_i,f'_j\rangle=\sum_{a,b}P_{ai}\langle e_a,f_b\rangle Q_{bj}$, and $\det P$, $\det Q$ are units. (iii) Clear.
\end{proof}

\begin{proposition}[orthogonals]\label{prop:modules:orthogonal}
Let $\langle\ ,\ \rangle\colon L\times L\to R$ be a perfect pairing on a lattice of rank $n$, not necessarily symmetric.
\begin{enumerate}[label=(\roman*)]
\item For every subset $S$, $S^\perp$ and ${}^\perp S$ are saturated.
\item For a submodule $N$ of rank $d$, $N^\perp$ and ${}^\perp N$ have rank $n-d$, and ${}^\perp(N^\perp)=\operatorname{sat}(N)=({}^\perp N)^\perp$. In particular
${}^\perp(N^\perp)=N=({}^\perp N)^\perp$ if $N$ is saturated.
\item If the pairing is symmetric or antisymmetric, left and right orthogonals agree.
\end{enumerate}
\end{proposition}

\begin{proof}
(i) If $ry\in S^\perp$ with $r\ne0$, then $r\langle s,y\rangle=0$, so $\langle s,y\rangle=0$ for $s\in S$. (ii) Since $\langle rx,y\rangle=0$ iff $\langle x,y\rangle=0$ for $r\ne0$,
$N^\perp=\operatorname{sat}(N)^\perp$, and we may assume $N$ saturated, hence a direct summand (\Cref{prop:modules:saturated}). Then the map $L\to N^\vee$,
$y\mapsto\langle-,y\rangle|_N$, is surjective: a functional on $N$ extends by $0$ on a complement to an element of $L^\vee$, which is $\langle-,y\rangle$ for some $y$ because
$\beta$ is onto. Its kernel is $N^\perp$, so $\rank N^\perp=n-d$; likewise for ${}^\perp N$. Now ${}^\perp(N^\perp)\supseteq N$, both are saturated (by (i)), and both
have rank $d$; a saturated module containing another of the same rank equals its saturation (\Cref{prop:modules:saturated}(ii)). (iii) Clear.
\end{proof}

\begin{theorem}[the Gram determinant of a saturated sublattice]\label{thm:modules:gram-cokernel}
Let $\mathcal O$ be a discrete valuation ring, $L$ a lattice of rank $n$ with a perfect pairing, $L'\subseteq L$ a \emph{saturated} sublattice of rank $d$, and
$\Delta$ the Gram determinant of the restriction of the pairing to $L'\times L'$ (in any basis).
\begin{enumerate}[label=(\roman*)]
\item The map $\rho\colon L\to\Hom_{\mathcal O}(L',\mathcal O)=L'^\vee$, $y\mapsto\langle-,y\rangle|_{L'}$, is surjective with kernel $L'^\perp$; so $L/L'^\perp\simeq L'^\vee$.
\item If $\Delta\ne0$, then $L'\cap L'^\perp=0$, the restriction $\rho|_{L'}\colon L'\to L'^\vee$ is injective, and
\[
L/(L'\oplus L'^\perp)\ \simeq\ \coker\bigl(L'\xrightarrow{\rho}L'^\vee\bigr)\ \simeq\ \bigoplus_{i=1}^d\mathcal O/(\delta_i),
\]
where the $(\delta_i)$ are the invariant factors of the Gram matrix of $L'$. Its length is $v(\Delta)$; when $k$ is finite its order is
$\#\mathcal O/\Delta\mathcal O=(\#k)^{v(\Delta)}$.
\item In (ii), $L'^\perp$ is saturated of rank $n-d$, and the Gram determinant $\Delta^\perp$ of the restriction of the pairing to $L'^\perp$ satisfies
$v(\Delta^\perp)=v(\Delta)$.
\item Saturation is necessary: in $L=\Zl^2$ with the standard pairing $\langle x,y\rangle=x_1y_1+x_2y_2$, the sublattice $L'=\Zl\cdot\ell e_1$ has $\Delta=\ell^2$, while
$L/(L'\oplus L'^\perp)=\Zl^2/(\ell\Zl\oplus\Zl)\simeq\Z/\ell$ has length $1$.
\end{enumerate}
\end{theorem}

\begin{proof}
(i) is the surjectivity shown in the proof of \Cref{prop:modules:orthogonal}(ii) (saturation makes $L'$ a direct summand); the kernel is $L'^\perp$ by
definition.
(ii) In a basis of $L'$ and the dual basis of $L'^\vee$, $\rho|_{L'}$ has matrix the transpose of the Gram matrix of $L'$, of determinant $\Delta\ne0$; so it is
injective, and its kernel $L'\cap L'^\perp$ is $0$. By (i), $\rho$ induces $L/L'^\perp\simeq L'^\vee$, under which $(L'+L'^\perp)/L'^\perp$ corresponds to
$\rho(L')$; hence $L/(L'\oplus L'^\perp)\simeq L'^\vee/\rho(L')$. By \Cref{thm:modules:smith} this cokernel is $\bigoplus\mathcal O/(\delta_i)$ with $\prod\delta_i=\Delta$ up to a unit,
of length $\sum v(\delta_i)=v(\Delta)$ (\Cref{prop:modules:length-dvr}).
(iii) $L'^\perp$ is saturated of rank $n-d$ by \Cref{prop:modules:orthogonal}. If $\Delta^\perp=0$, some $0\ne y\in L'^\perp\otimes K$ is right-orthogonal to
$L'^\perp\otimes K$; it is also right-orthogonal to $L'\otimes K$ (as $y\in L'^\perp\otimes K$); and $L'\otimes K+L'^\perp\otimes K=L\otimes K$ by (ii) (the quotient is
torsion), so $y$ is right-orthogonal to everything, contradicting perfectness. So $\Delta^\perp\ne0$. Apply (ii) to the transposed pairing
$\langle x,y\rangle^{\mathsf T}=\langle y,x\rangle$, which is perfect, and to the saturated $N=L'^\perp$: its Gram determinant is $\Delta^\perp$ (transposed), its right orthogonal
for the transposed pairing is the left orthogonal ${}^\perp(L'^\perp)=L'$ (\Cref{prop:modules:orthogonal}(ii)), so $L/(L'^\perp\oplus L')$ has length
$v(\Delta^\perp)$. It also has length $v(\Delta)$.
(iv) $L'^\perp=\Zl e_2$ and the computation is immediate.
\end{proof}

\begin{example}\label{ex:modules:gram}
On $\Z^2$ take the Gram matrix $\Mtwo9110$ (perfect: determinant $-1$) and $L'=\Z e_1$, saturated, with $\Delta=9$. Then
$L'^\perp=\{y:9y_1+y_2=0\}=\Z(1,-9)$, and $L'\oplus L'^\perp$ has basis $(1,0)$, $(1,-9)$, so index $\abs{\det\Mtwo101{-9}}=9=\Delta$, as
\Cref{thm:modules:gram-cokernel}(ii) predicts; the Gram determinant on $L'^\perp$ is $\langle(1,-9),(1,-9)\rangle=9-9-9=-9$, of the same valuation. With
the standard pairing and the non-saturated $L'=\Z\cdot3e_1$ the Gram determinant is $9$ but the index of $L'\oplus L'^\perp=3\Z\oplus\Z$ is $3$.
\end{example}

\begin{paperlink}\label{pl:modules:gram}
In Lemma lem:ddt417, ``the pairing identifies $L/L[\wp]^\perp$ with $\Hom(L[\wp],\mathcal O)$, because $L[\wp]$ is saturated'' is
\Cref{thm:modules:gram-cokernel}(i) with $L'=L[\wp]$, saturated by \Cref{prop:modules:saturated}(iii); ``$Q$ is the cokernel of $L[\wp]\to\Hom(L[\wp],\mathcal O)$,
of order $\#\mathcal O/\Delta$'' is part (ii), once $L[\wp]^\perp=L[I]$ is known (\Cref{prop:modules:ddt-ingredients}). Part (iv) shows that the word
``saturated'' cannot be dropped. For a residue field that is not finite the order must be read as the length $v(\Delta)$
(\Cref{conv:modules:rings}(f)); the revised paper assumes $k$ finite.
\end{paperlink}

\begin{proposition}[alternating pairings]\label{prop:modules:alternating}
A pairing $L\times L\to R$ is \emph{alternating} if $\langle x,x\rangle=0$ for all $x$. An alternating pairing is antisymmetric; conversely an antisymmetric
pairing is alternating if $2$ is a nonzerodivisor in $R$ (for instance $R$ a domain of characteristic $0$).
\begin{enumerate}[label=(\roman*)]
\item The Gram matrix of an alternating pairing is antisymmetric with zero diagonal; its determinant vanishes in odd rank, so a perfect alternating
pairing on a nonzero lattice has even rank. In even rank $\det G=\operatorname{Pf}(G)^2$ for a polynomial $\operatorname{Pf}$ in the entries (the Pfaffian); in rank $4$,
$\operatorname{Pf}(G)=g_{12}g_{34}-g_{13}g_{24}+g_{14}g_{23}$.
\item Rank two: $G=\Mtwo0a{-a}0$ with $a=\langle e_1,e_2\rangle$, so $\Delta=a^2$; and $\langle x,y\rangle=a(x_1y_2-x_2y_1)$, so the ideal of values $\langle L,L\rangle=\{\langle x,y\rangle\}$ is
$aR$. Hence, without reference to a basis, $\Delta R=(\text{ideal of values})^2$.
\item (Symplectic normal form.) Over a principal ideal domain an alternating pairing has a basis in which $G$ is block diagonal with blocks
$\Mtwo0{d_i}{-d_i}0$, $d_1\mid d_2\mid\cdots$, followed by zeros; it is perfect iff all $d_i$ are units and there is no zero block.
\end{enumerate}
\end{proposition}

\begin{proof}
$0=\langle x+y,x+y\rangle=\langle x,y\rangle+\langle y,x\rangle$; conversely antisymmetry gives $2\langle x,x\rangle=0$.
(i) The diagonal entries are $\langle e_i,e_i\rangle=0$. For odd $n$, let $X$ be the generic alternating matrix over the polynomial ring
$\Z[x_{ij}:i<j]$; then $\det X=\det X^{\mathsf T}=\det(-X)=-\det X$, and as $\Z[x_{ij}]$ is a domain of characteristic $0$, $\det X=0$; specialising the
variables gives $\det G=0$ over any $R$. A perfect pairing has $\det G$ a unit, and units are nonzero when $R\ne0$. The identity $\det G=\operatorname{Pf}(G)^2$
in general is classical and we do not prove it (see \cite[Chapter XV, Section 9]{Lang}, section number not verified); in rank $4$ it is a
direct expansion, which we checked with PARI.
(ii) Expand $\langle x_1e_1+x_2e_2,y_1e_1+y_2e_2\rangle$; the value $x_1y_2-x_2y_1$ is $1$ for $x=e_1$, $y=e_2$.
\end{proof}

\begin{proof}[Sketch of proof of (iii)]
Choose $x,y$ with $\langle x,y\rangle=d_1$ a generator of the ideal of values (if it is $0$ the pairing is zero). The map
$z\mapsto z-\frac{\langle z,y\rangle}{d_1}x+\frac{\langle z,x\rangle}{d_1}y$ is defined over $R$ because $d_1$ divides every value, and it projects $L$ onto the orthogonal of
$\{x,y\}$, giving $L=(Rx\oplus Ry)\oplus\{x,y\}^\perp$; induct, noting that the ideal of values only shrinks. See \cite[Chapter XV, Section 8]{Lang}
(section number not verified).
\end{proof}

\begin{paperlink}\label{pl:modules:alternating}
Section 4.1: ``the restriction of $\hat\delta_!$ to the rank-two lattice $M_{f^\Sigma,\lambda}$ takes its values in an ideal of $\Zl$, \dots, which is
$\eta^\Sigma_{f,\lambda}$; equivalently the Gram determinant of the pairing on $M_{f^\Sigma,\lambda}$ is $(\eta^\Sigma_{f,\lambda})^2$ up to a unit''. This is
\Cref{prop:modules:alternating}(ii), as $\hat\delta_!$ is alternating. The same proposition is behind ``the Gram determinant of a rank-two alternating
form is the square of its ideal of values'' in the proof of Proposition prop:pullback, and behind $d=2$, $\Delta=(\eta^\emptyset_{f,\lambda})^2$ when Lemma
lem:ddt417 is applied in Step 2 of the proof of Theorem thm:quant.
\end{paperlink}

\begin{proposition}[adjoints]\label{prop:modules:adjoint}
Let $\langle\ ,\ \rangle$ be a perfect pairing on a lattice $L$, with Gram matrix $G$ in a basis, coordinates written as columns so that
$\langle x,y\rangle=x^{\mathsf T}Gy$.
\begin{enumerate}[label=(\roman*)]
\item Every $A\in\End_R(L)$ has a unique \emph{adjoint} $A^\dagger$ with $\langle Ax,y\rangle=\langle x,A^\dagger y\rangle$; in matrices $A^\dagger=G^{-1}A^{\mathsf T}G$. Moreover
$(AB)^\dagger=B^\dagger A^\dagger$, and $(A^\dagger)^\dagger=A$ if the pairing is symmetric or antisymmetric.
\item Let $\T\subseteq\End_R(L)$ be a commutative subring of self-adjoint endomorphisms. For every ideal $\mf a\subseteq\T$, $(\mf aL)^\perp=L[\mf a]$.
\item In (ii), every idempotent $e$ of $\T$, or of $\T\otimes_RS$ acting on $L\otimes_RS$ for an $R$-algebra $S$ that is a principal ideal domain, is self-adjoint;
$L=eL\oplus(1-e)L$ is an orthogonal decomposition (on both sides), and the restriction of the pairing to $eL$ is perfect.
\item (Twisted pairings.) Let $w\in\Aut_R(L)$ satisfy $A^\dagger=wAw^{-1}$ for all $A\in\T$, and put $[x,y]=\langle x,wy\rangle$. Then $[\ ,\ ]$ is perfect and every
$A\in\T$ is self-adjoint for it. If $\langle\ ,\ \rangle$ is antisymmetric and $w^\dagger=w$ (for instance if $w$ is an involution preserving the pairing, so that
$w^\dagger=w^{-1}=w$), then $[\ ,\ ]$ is antisymmetric, hence alternating when $2$ is a nonzerodivisor.
\end{enumerate}
\end{proposition}

\begin{proof}
(i) For fixed $y$, $x\mapsto\langle Ax,y\rangle$ is in $L^\vee$, so equals $\langle-,z\rangle$ for a unique $z$ ($\beta$ is bijective); put $A^\dagger y=z$. In matrices
$x^{\mathsf T}A^{\mathsf T}Gy=x^{\mathsf T}GA^\dagger y$ for all $x,y$, so $GA^\dagger=A^{\mathsf T}G$. Then $\langle ABx,y\rangle=\langle Bx,A^\dagger y\rangle=\langle x,B^\dagger A^\dagger y\rangle$. If
$\langle x,y\rangle=\pm\langle y,x\rangle$, then $\langle A^\dagger x,y\rangle=\pm\langle y,A^\dagger x\rangle=\pm\langle Ay,x\rangle=\langle x,Ay\rangle$.
(ii) $y\in(\mf aL)^\perp$ iff $\langle ax,y\rangle=\langle x,ay\rangle=0$ for all $a\in\mf a$ and $x\in L$, iff $ay=0$ for all $a\in\mf a$ (perfectness).
(iii) Elements of $\T$ are self-adjoint by hypothesis, and self-adjointness is preserved by $S$-linear combinations, so every element of
$\T\otimes S$ is self-adjoint for the extended pairing (which is perfect, \Cref{prop:modules:gram}(iii)). For $x\in eL$ and $y\in(1-e)L$,
$\langle x,y\rangle=\langle ex,y\rangle=\langle x,ey\rangle=0$ since $ey=e(1-e)y=0$; likewise $\langle y,x\rangle=0$. The summands are free (direct summands of a lattice over a
principal ideal domain), and in a basis adapted to the decomposition the Gram matrix is block diagonal, with unit determinant equal to the
product of the two block determinants; so each block determinant is a unit.
(iv) $x\mapsto[x,-]$ is the composite of $\alpha$ with the isomorphism induced by $w$, so $[\ ,\ ]$ is perfect. For $A\in\T$,
$[Ax,y]=\langle x,A^\dagger wy\rangle=\langle x,wAw^{-1}wy\rangle=[x,Ay]$. If $\langle\ ,\ \rangle$ is antisymmetric and $w^\dagger=w$, then
$[y,x]=\langle y,wx\rangle=-\langle wx,y\rangle=-\langle x,w^\dagger y\rangle=-[x,y]$. An involution preserving the pairing satisfies $\langle wx,y\rangle=\langle w^2x,wy\rangle=\langle x,wy\rangle$.
\end{proof}

\begin{paperlink}\label{pl:modules:adjoint}
Section 4.1: the cup product on $H^1(X_0(M^\Sigma),\Zl)$ is perfect and alternating, the Atkin-Lehner involution $w=w_{M^\Sigma}$ preserves it, and
$wT_nw^{-1}$ is the adjoint of $T_n$ for the cup product (\Cref{ch:curves,ch:moddeg}). By \Cref{prop:modules:adjoint}(iv) the pairing
$\hat\delta_!$, which is $\frac12\langle x,wy\rangle$ on classes from $X_0(M^\Sigma)$, is perfect (the factor $\frac12$ being a unit as $\ell$ is odd),
alternating (antisymmetric over the domain $\Zl$ of characteristic $0$, \Cref{prop:modules:alternating}) and makes every Hecke operator, $U_p$ included, self-adjoint: ``this is what the
composition with $w$ is for''. In the paragraph after Lemma lem:hecke (and in Step 2 of the proof of Theorem thm:quant) the paper uses
\Cref{prop:modules:adjoint}(iii): ``the idempotent $e_{\mf m}$ is self-adjoint for $\hat\delta_!$, so $H^1(X_0(N),\Zl)$ splits orthogonally into its local
components and the pairing is perfect on each''. The idempotent $e_{\mf m}$ lies in $\T\otimes\Zl$, not in $\T$, which is why (iii) allows $\T\otimes S$.
\end{paperlink}

\begin{proposition}[an algebra paired with a module]\label{prop:modules:algebra-duality}
Let $\T$ be a commutative $R$-algebra that is free of finite rank as an $R$-module, $S$ a $\T$-module free of finite rank over $R$, and
$\langle\ ,\ \rangle\colon S\times\T\to R$ a perfect $R$-bilinear pairing with $\langle ts,u\rangle=\langle s,tu\rangle$ for $s\in S$ and $t,u\in\T$.
\begin{enumerate}[label=(\roman*)]
\item $s\mapsto\langle s,-\rangle$ is an isomorphism of $\T$-modules $S\simeq\T^\vee=\Hom_R(\T,R)$, where $(t\varphi)(u)=\varphi(tu)$.
\item For every ideal $\mf a\subseteq\T$, $S[\mf a]$ corresponds to $\Hom_R(\T/\mf a,R)$, the functionals vanishing on $\mf a$; it is saturated in $S$, of rank
$\rank_R(\T/\mf a)$.
\item Let $\mf a,\mf b\subseteq\T$ be ideals with $\mf a\cap\mf b=0$, $\T/\mf a$ and $\T/\mf b$ torsion-free, and $\rank\T/\mf a+\rank\T/\mf b=\rank\T$. Then
$\T\into\T/\mf a\times\T/\mf b$ has torsion cokernel $\T/(\mf a+\mf b)$ (\Cref{prop:modules:two-submodules}), $S[\mf a]\oplus S[\mf b]\subseteq S$, and
$S/(S[\mf a]\oplus S[\mf b])\simeq\Hom_R(\T/(\mf a+\mf b),K/R)$; for $R=\Z$ or $\Zl$ this is $(\T/(\mf a+\mf b))^*$, of the same order as $\T/(\mf a+\mf b)$.
\item $\T$ acts faithfully on $S$.
\end{enumerate}
\end{proposition}

\begin{proof}
(i) It is an $R$-isomorphism because the pairing is perfect, and $\T$-linear: $\langle ts,u\rangle=\langle s,tu\rangle=(t\langle s,-\rangle)(u)$.
(ii) $s\in S[\mf a]$ iff $\langle as,u\rangle=\langle s,au\rangle=0$ for all $a\in\mf a$, $u\in\T$, iff $\langle s,-\rangle$ vanishes on the ideal $\mf a$. Such functionals are
$\Hom_R(\T/\mf a,R)=\Hom_R((\T/\mf a)/\tors,R)$, of rank $\rank(\T/\mf a)$; saturation is \Cref{prop:modules:saturated}(iii).
(iii) The cokernel of $\T\to\T/\mf a\times\T/\mf b$ is $\T/(\mf a+\mf b)$, and the two lattices have the same rank by hypothesis, so it is torsion. Apply
\Cref{prop:modules:dual-inclusion} to this injection of lattices: $(\T/\mf a\times\T/\mf b)^\vee\into\T^\vee$, with cokernel $\Hom(\T/(\mf a+\mf b),K/R)$. Under (i)
and (ii) the image of $(\T/\mf a)^\vee\oplus(\T/\mf b)^\vee$ in $\T^\vee$ is $S[\mf a]+S[\mf b]$; this sum is direct because the source is a direct sum and the map is injective.
(iv) If $tS=0$ then $\langle s,tu\rangle=\langle ts,u\rangle=0$ for all $s,u$, so $tu=0$ for all $u$ (perfectness); with $u=1$, $t=0$.
\end{proof}

\begin{paperlink}\label{pl:modules:algebra-duality}
In Lemma lem:crit, $R=\Z$, $\T$ is the Hecke algebra, $S=S_2(\Gamma_0(N);\Z)$ and $\langle g,t\rangle=a_1(tg)$ is perfect (\Cref{ch:hecke}); it satisfies
$\langle tg,u\rangle=a_1(utg)=\langle g,tu\rangle$. Take $\mf a=\ker\phi_f$ and $\mf b=\Ann_\T(\mf a)$, the paper's $I$ and $J$ (the paper defines $J=\Ann_\T(f^\perp)$; it equals
$\Ann_\T(I)$ because $\T\otimes\Q=\Q\times\T'$ with $\T'$ acting faithfully on $f^\perp$). Then $S[\mf a]=\Z f$ ($f$ is the primitive form killed by $\mf a$), and
$S[\mf b]=L$ is the orthogonal complement lattice. By \Cref{prop:modules:congruence-ideal}(iii), $\mf a=\T\cap(0\times\T')$ and $\mf b=\T\cap(\Q\times0)$, so
$\T/\mf a$ and $\T/\mf b$ are torsion-free (\Cref{prop:modules:saturated}(iii)), and $\mf b$ is a lattice in $\Q\times0$, so $\rank\T/\mf b=\rank\T-1$ and the
ranks are complementary. So the paper's dualising step (\Cref{pl:modules:crit-second}) is \Cref{prop:modules:algebra-duality}(iii). The paper does not state these hypotheses, but its assertion that
$\T/(I+J)\simeq\Z/\phi_f(J)$ is cyclic \emph{of order} $r'$ presupposes them: as $I\cap J=0$, the quotient is finite iff $\phi_f(J)\ne0$, iff $J\ne0$, which is
the rank condition.
\end{paperlink}

\section[Finite free algebras and generalised eigenspaces]{Finite free algebras over a discrete valuation ring; idempotents; generalised eigenspaces}\label{sec:modules:algebras}

The Hecke algebra $\T$ is a commutative ring that is free of finite rank over $\Z$. The paper uses three of its structural features. After
tensoring with $\Q$ it is a product of local Artinian algebras, one of which is $\Q$, the factor of $f$ (``$\T\otimes\Q=\Q\times\T'$'', Lemma lem:crit).
After tensoring with $\Zl$ it is a product of local rings $\T_{\mf m}$ cut out by idempotents (Lemma lem:hecke, Remark rem:free, Theorem thm:quant).
And its action on forms decomposes into generalised eigenspaces, on which the operators need not be semisimple (Lemma lem:old, Theorem
thm:twist, Step 1 of Theorem thm:quant). In the language of linear algebra this section is the Jordan decomposition of a commuting family of
matrices, done over $\Z$ and $\Zl$ rather than over $\C$.

\begin{theorem}[finite-dimensional commutative algebras]\label{thm:modules:artinian-algebras}
Let $A$ be a commutative algebra of finite dimension over a field $F$.
\begin{enumerate}[label=(\roman*)]
\item Every prime ideal of $A$ is maximal; $A$ has finitely many maximal ideals $\mf m_1,\dots,\mf m_s$, with $s\le\dim_FA$; and the Jacobson radical
$\mf J=\bigcap_i\mf m_i$ is nilpotent.
\item $A\simeq\prod_{i=1}^sA_{\mf m_i}$, each factor a local ring with nilpotent maximal ideal; the factors are cut out by idempotents $e_1,\dots,e_s$ with
$\sum e_i=1$ and $e_ie_j=0$ for $i\ne j$, and $A_{\mf m_i}\simeq e_iA\simeq A/\mf m_i^N$ for $N$ large.
\item $A$ is reduced (has no nonzero nilpotents) iff $\mf J=0$ iff $A$ is a product of fields.
\item Every $A$-module $V$ is $\bigoplus_ie_iV$, and $e_iV=V_{\mf m_i}$.
\end{enumerate}
\end{theorem}

\begin{proof}
(i) If $P$ is prime, $A/P$ is a finite-dimensional domain, hence a field (multiplication by $a\ne0$ is an injective, hence surjective, linear map).
If $\mf m_1,\dots,\mf m_s$ are distinct maximal ideals, they are pairwise comaximal and $A\to\prod A/\mf m_i$ is surjective by the Chinese remainder theorem
(\Cref{ch:rings}), so $s\le\dim A$. The powers $\mf J^k$ form a decreasing chain of subspaces, so $\mf J^N=\mf J^{N+1}=\mf J\cdot\mf J^N$ for some $N$; by
\Cref{thm:modules:nakayama}(i) (for an ideal in the Jacobson radical), $\mf J^N=0$.
(ii) The ideals $\mf m_i^N$ are pairwise comaximal (if $\mf m_i^N+\mf m_j^N$ were proper it would lie in a maximal ideal containing both $\mf m_i$ and $\mf m_j$), and
$\prod_i\mf m_i^N\subseteq\mf J^N=0$. By the Chinese remainder theorem $A\simeq\prod_iA/\mf m_i^N$. The maximal ideals of $A/\mf m_i^N$ correspond to maximal ideals
containing $\mf m_i^N$, hence $\mf m_i$; so $A/\mf m_i^N$ is local with maximal ideal $\mf m_i/\mf m_i^N$, nilpotent. The idempotents are the unit vectors of the
product, and localisation at $\mf m_i$ picks out the $i$-th factor (\Cref{thm:modules:localisation}(vii)).
(iii) If $A$ is reduced, the nilpotent $\mf J$ is $0$, and $A\simeq\prod A/\mf m_i$ by (ii) with $N=1$. A product of fields is reduced.
(iv) $v=\sum e_iv$, and $e_iV\cap\sum_{j\ne i}e_jV=0$ by applying $e_i$. In $A_{\mf m_i}$, $e_i$ becomes a unit ($e_i\notin\mf m_i$) and $e_j=0$ for $j\ne i$ (as
$e_ie_j=0$); so $(e_jV)_{\mf m_i}=0$, while elements outside $\mf m_i$ act invertibly on $e_iV$, a module over the local ring $e_iA$; hence $V_{\mf m_i}=e_iV$.
\end{proof}

\begin{theorem}[finite free algebras over a discrete valuation ring]\label{thm:modules:finite-algebra-dvr}
Let $\T$ be a commutative $\mathcal O$-algebra that is finite and free as an $\mathcal O$-module (a \emph{finite free} $\mathcal O$-algebra).
\begin{enumerate}[label=(\roman*)]
\item $\T$ is Noetherian. Every maximal ideal $\mf M$ of $\T$ contains $\varpi$, so the maximal ideals of $\T$ correspond to those of the finite
$k$-algebra $\T/\varpi\T$; they are finite in number, and their residue fields are finite extensions of $k$.
\item $\T$ is torsion-free and $\T\subseteq\T_K:=\T\otimes_{\mathcal O}K$, a finite-dimensional $K$-algebra. $\T$ is reduced iff $\T_K$ is, iff $\T$ embeds in a product
of fields.
\item Let $\pi\colon\T\to\mathcal O$ be a surjective homomorphism of $\mathcal O$-algebras, with kernel $\wp$. Then $\wp$ is a prime ideal, $\T=\mathcal O\cdot1\oplus\wp$ as
$\mathcal O$-modules, and $\mf m_\pi=\pi^{-1}(\varpi\mathcal O)=\wp+\varpi\T$ is a maximal ideal with residue field $k$. The extension $\pi_K\colon\T_K\to K$ is a surjective
homomorphism; it factors through exactly one of the local factors of $\T_K$ (\Cref{thm:modules:artinian-algebras}).
\end{enumerate}
\end{theorem}

\begin{proof}
(i) $\T$ is a finitely generated module over the Noetherian $\mathcal O$ (\Cref{prop:modules:noetherian}). Let $\mf M$ be maximal; $\T/\mf M$ is a field which is a finitely
generated module over the image $D=\mathcal O/(\mf M\cap\mathcal O)$ of $\mathcal O$, a domain. Then $D$ is a field: for $0\ne x\in D$, the element $y=x^{-1}\in\T/\mf M$ satisfies a monic
equation $y^n+d_1y^{n-1}+\cdots+d_n=0$ over $D$ (\Cref{lem:modules:determinant-trick} with $\mf a=D$, applied to multiplication by $y$), and multiplying by $x^{n-1}$
gives $y=-(d_1+d_2x+\cdots+d_nx^{n-1})\in D$. So $\mf M\cap\mathcal O$ is maximal, that is, $\varpi\mathcal O$, and $\varpi\in\mf M$. The rest follows from
\Cref{thm:modules:isomorphism}(iii) and \Cref{thm:modules:artinian-algebras}(i), as $\T/\mf M$ is a quotient of the finite-dimensional $k$-algebra $\T/\varpi\T$.
(ii) Free modules are torsion-free and embed in their tensor product with $K$ (\Cref{prop:modules:torsion}(i)). If $x\in\T_K$ is nilpotent, so is
$\varpi^kx\in\T$ for $k$ large; so $\T$ is reduced iff $\T_K$ is, and then $\T\subseteq\T_K$, a product of fields (\Cref{thm:modules:artinian-algebras}(iii)).
Conversely a subring of a product of fields is reduced.
(iii) $\T/\wp\simeq\mathcal O$ is a domain. $t=\pi(t)\cdot1+(t-\pi(t)\cdot1)$ with $t-\pi(t)\cdot1\in\wp$, and $\mathcal O\cdot1\cap\wp=0$ because $\pi(a\cdot1)=a$. If $\pi(t)=\varpi a$ then
$t=\varpi a\cdot1+(t-\pi(t)\cdot1)\in\varpi\T+\wp$; the reverse inclusion is clear, and $\T/\mf m_\pi\simeq\mathcal O/\varpi=k$. The kernel of $\pi_K$ is a maximal ideal of
$\T_K=\prod_iA_i$, hence of the form $A_1\times\cdots\times\mf m_{A_i}\times\cdots\times A_s$ for one $i$ (\Cref{ch:rings}: ideals of a product); so $\pi_K(e_j)=0$ for $j\ne i$
and $\pi_K(e_i)=1$.
\end{proof}

\begin{theorem}[idempotent lifting and the decomposition $\T=\prod\T_{\mf m}$]\label{thm:modules:idempotent-lifting}
Let $\mathcal O$ be a \emph{complete} discrete valuation ring (for example $\Zl$, or $\mathcal O_\lambda$ of \Cref{ch:local}), and $\T$ a finite free $\mathcal O$-algebra.
\begin{enumerate}[label=(\roman*)]
\item Every idempotent $\bar e$ of $\T/\varpi\T$ lifts to a unique idempotent of $\T$.
\item $\T\simeq\prod_{\mf m}\T_{\mf m}$, the product over the maximal ideals of $\T$, where $\T_{\mf m}=e_{\mf m}\T$ for an idempotent $e_{\mf m}$, and $\T_{\mf m}$ is the
localisation of $\T$ at $\mf m$; each $\T_{\mf m}$ is a local ring and a finite free $\mathcal O$-algebra. Moreover $e_{\mf m}\notin\mf m$ and $e_{\mf m'}\in\mf m$ for $\mf m'\ne\mf m$.
\item For every $\T$-module $L$, $L=\bigoplus_{\mf m}e_{\mf m}L$ and $L_{\mf m}=e_{\mf m}L$; if $L$ is finitely generated and free over $\mathcal O$, so is each $L_{\mf m}$.
\item A ring homomorphism $\phi$ from $\T$ to a local ring factors through exactly one $\T_{\mf m}$. For $\pi\colon\T\to\mathcal O$ as in
\Cref{thm:modules:finite-algebra-dvr}(iii): $\pi(e_{\mf m_\pi})=1$ and $\pi(e_{\mf m'})=0$ for $\mf m'\ne\mf m_\pi$.
\item Completeness is necessary: $\Z_{(11)}[x]/(x^2+x-1)$ is a finite free $\Z_{(11)}$-algebra with two maximal ideals, $(11,x-3)$ and $(11,x-7)$, but no
idempotent other than $0$ and $1$.
\end{enumerate}
\end{theorem}

\begin{proof}
(i) \emph{Existence.} Let $e_0\in\T$ lift $\bar e$, and define $e_{j+1}=3e_j^2-2e_j^3$. Put $d_j=e_j^2-e_j$. A direct computation gives
\[
e_{j+1}-e_j=(1-2e_j)d_j,\qquad e_{j+1}^2-e_{j+1}=d_j^2\,(4d_j-3).
\]
(Expand both sides as polynomials in $e=e_j$, with $d=e^2-e$: the first is $3e^2-2e^3-e$ on both sides, and the second is
$4e^6-12e^5+9e^4+2e^3-3e^2$ on both sides.) Since
$d_0\in\varpi\T$, induction gives $d_j\in\varpi^{2^j}\T$, and $e_{j+1}-e_j\in\varpi^{2^j}\T$. As $\T\simeq\mathcal O^n$ is complete
(\Cref{thm:modules:zl}(vi) holds verbatim over any complete $\mathcal O$), the $e_j$ converge to some $e$, with $e\equiv e_0\pmod\varpi$ and $e^2-e=\lim d_j=0$.
\emph{Uniqueness.} If $e$, $e'$ are idempotents with $e\equiv e'\pmod\varpi$, then $x=e-e'$ satisfies $x^3=x$ (expand, using $e^2=e$, $e'^2=e'$ and
commutativity: $x^2=e-2ee'+e'$ and $x^3=e-e'$). So $x=x^3=x^9=\cdots\in\bigcap_k\varpi^k\T=0$.
(ii) By \Cref{thm:modules:artinian-algebras}, $\T/\varpi\T=\prod_{\bar{\mf m}}A_{\bar{\mf m}}$ with orthogonal idempotents $\bar e_{\mf m}$ summing to $1$; the $\bar{\mf m}$ correspond to the
maximal ideals $\mf m$ of $\T$ (\Cref{thm:modules:finite-algebra-dvr}(i)). Lift them to idempotents $e_{\mf m}$. For $\mf m\ne\mf m'$, $e_{\mf m}e_{\mf m'}$ is an
idempotent lifting $0$, hence $0$; and $\sum e_{\mf m}$ is an idempotent lifting $1$, hence $1$. So $\T=\prod e_{\mf m}\T$. Each $e_{\mf m}\T$ is a direct summand of $\T$, so
finite free over $\mathcal O$ (\Cref{thm:modules:submodule-free}), and a finite free $\mathcal O$-algebra with identity $e_{\mf m}$; its maximal ideals contain $\varpi$ and
correspond to those of $e_{\mf m}\T/\varpi e_{\mf m}\T=A_{\bar{\mf m}}$, which is local. So $e_{\mf m}\T$ is local, and the maximal ideal of $\T$ corresponding to it is the
product of its maximal ideal with the other factors, which is $\mf m$; it contains $e_{\mf m'}$ for $\mf m'\ne\mf m$ and not $e_{\mf m}$. Localisation at $\mf m$ picks out the
factor $e_{\mf m}\T$ (\Cref{thm:modules:localisation}(vii)).
(iii) As in \Cref{thm:modules:artinian-algebras}(iv). The summands of a free $\mathcal O$-module of finite rank are free.
(iv) The image of an idempotent in a local ring is $0$ or $1$ ($e(1-e)=0$ and one of $e$, $1-e$ is a unit). The images of the $e_{\mf m}$ are orthogonal and
sum to $1$, so exactly one is $1$ and $\phi$ factors through that factor. For $\pi$, reduce modulo $\varpi$: $\bar\pi\colon\T\to k$ has kernel $\mf m_\pi$, which contains
$e_{\mf m'}$ for $\mf m'\ne\mf m_\pi$ and not $e_{\mf m_\pi}$, by (ii).
(v) The ring is a subring of $\Q(\sqrt5)$, a domain (the polynomial is irreducible over $\Q$, having discriminant $5$), so $e(1-e)=0$ forces $e\in\{0,1\}$. It is free
with basis $1,x$ over $\Z_{(11)}$, so its maximal ideals contain $11$ (\Cref{thm:modules:finite-algebra-dvr}(i)) and correspond to those of
$\F_{11}[x]/((x-3)(x-7))\simeq\F_{11}\times\F_{11}$, since $x^2+x-1\equiv(x-3)(x-7)\pmod{11}$.
\end{proof}

\begin{corollary}[Hecke-type algebras]\label{cor:modules:hecke-type}
Let $\T$ be a commutative ring that is free of finite rank as a $\Z$-module. Then $\T\otimes\Zl$ is a finite free $\Zl$-algebra; its maximal ideals correspond
to the maximal ideals $\mf m$ of $\T$ containing $\ell$ (as $\mf m\otimes\Zl$); and $\T\otimes\Zl=\prod_{\mf m\ni\ell}\T_{\mf m}$ with $\T_{\mf m}=(\T\otimes\Zl)_{\mf m}$, the convention
of \Cref{conv:modules:tm}.
\end{corollary}

\begin{proof}
$\T\otimes\Zl\simeq\Zl^n$. Its maximal ideals correspond to those of $(\T\otimes\Zl)/\ell=\T/\ell\T$ (\Cref{thm:modules:finite-algebra-dvr}(i) and
\Cref{thm:modules:zl}(iii)), which are the maximal ideals of $\T$ containing $\ell$. Apply \Cref{thm:modules:idempotent-lifting}.
\end{proof}

\begin{remark}\label{rem:modules:completion}
One can show that $\T_{\mf m}$ is also the $\mf m$-adic completion $\varprojlim\T/\mf m^k$ of $\T$; we do not use this, and do not prove it (see \cite[Chapter 7]{Eisenbud}).
\end{remark}

\begin{example}[the Hecke algebra of level $23$]\label{ex:modules:level23}
Let $\T=\Z[x]/(x^2+x-1)=\Z[\frac{1+\sqrt5}2]$. This is the Hecke algebra of $S_2(\Gamma_0(23))$, with $x=T_2$ (\Cref{ch:ant}): the newforms of level $23$ have
$a_2=(-1\pm\sqrt5)/2$, the roots of $x^2+x-1$ (PARI: \texttt{charpoly(mfheckemat(mfinit([23,2],1),2))}); $\T_{23}\otimes\Q$ has dimension
$\dim S_2(\Gamma_0(23))=2$ (\Cref{prop:modules:algebra-duality}(i)), so it is $\Q(\sqrt5)$; and $\Z[T_2]\subseteq\T_{23}$ is an order of discriminant $5$, squarefree,
hence the maximal order, so $\T_{23}=\Z[T_2]$.
\begin{itemize}
\item $\ell=11$: $x^2+x-1\equiv(x-3)(x-7)\pmod{11}$, so $\T\otimes\Z_{11}\simeq\Z_{11}\times\Z_{11}$, by $x\mapsto(r_1,r_2)$ with $r_1\equiv7$, $r_2\equiv3$ the roots in $\Z_{11}$.
Modulo $11^6=1771561$ they are $r_1=728020$ and $r_2=1043540$ (PARI's \texttt{polrootspadic}; $r_1+r_2=11^6-1\equiv-1$). The idempotent of the first
factor is $e=(x-r_2)/(r_1-r_2)$, which is $\equiv784x+1058\pmod{11^3}$ and $\equiv1708457x+854229\pmod{11^6}$; it is also the limit of the iteration
$e\mapsto3e^2-2e^3$ started at $(x-3)/4$ (five steps give it modulo $11^6$).
\item $\ell=5$: $x^2+x-1\equiv(x-2)^2\pmod5$, so $\T\otimes\Z_5$ is local; it is the ring of integers of $\Q_5(\sqrt5)$, a ramified discrete valuation ring.
\item $\ell=2,3$: $x^2+x-1$ is irreducible modulo $2$ and modulo $3$, so $\T\otimes\Z_2$ and $\T\otimes\Z_3$ are local, unramified discrete valuation rings with residue
fields $\F_4$ and $\F_9$.
\end{itemize}
A second example: $\T=\Z[x]/(x^2-4)\simeq\{(a,b)\in\Z^2:a\equiv b\pmod4\}$ (by $x\mapsto(2,-2)$). Then $\T\otimes\Z_3\simeq\Z_3\times\Z_3$ with idempotent $(x+2)/4$, while
$\T\otimes\Z_2$ is local ($x^2-4\equiv x^2\pmod2$) and not a domain ($(x-2)(x+2)=0$).
\begin{lstlisting}[language=GP]
polrootspadic(x^2+x-1, 11, 6)      \\ roots = 3 + ... and 7 + ... in Z_11
e = Mod(Mod(1,11^6)*(x-3)/4, x^2+x-1);
for(i=1,5, e = 3*e^2 - 2*e^3); e  \\ Mod(1708457*x + 854229, x^2+x-1)
factormod(x^2+x-1, 5)              \\ (x+3)^2, that is (x-2)^2
\end{lstlisting}
\end{example}

\begin{paperlink}\label{pl:modules:tm}
``$\mf m$ the maximal ideal of $\T\otimes\Zl$ containing $\ell$ and $\ker\phi_f$'' (Lemma lem:hecke) is $\mf m_{\phi_f}$ of
\Cref{thm:modules:finite-algebra-dvr}(iii), for $\phi_f\otimes\Zl\colon\T\otimes\Zl\to\Zl$. By \Cref{thm:modules:idempotent-lifting}(iv), $\phi_f$ factors through
$\T_{\mf m}$ and kills the other local factors; in Step 1 of the proof of Theorem thm:quant, ``projecting $h$ to the $\mf m$-component (an idempotent of
$\T_\ell$\dots)'' uses the idempotent $e_{\mf m}$ of (ii). Remark rem:free is about $H^1(X_0(N),\Zl)_{\mf m}=e_{\mf m}H^1(X_0(N),\Zl)$, a free $\Zl$-module by (iii).
\end{paperlink}

\begin{theorem}[generalised eigenspaces]\label{thm:modules:generalised-eigenspaces}
\begin{enumerate}[label=(\roman*)]
\item Let $V$ be a nonzero finite-dimensional vector space over a field $F$ and $\T\subseteq\End_F(V)$ a commutative subalgebra such that the characteristic
polynomial of every $t\in\T$ splits into linear factors over $F$. For an $F$-algebra homomorphism $\phi\colon\T\to F$ (a \emph{system of eigenvalues}) put
\[
V(\phi)=\{v\in V:(t-\phi(t))^{\dim V}v=0\text{ for all }t\in\T\}.
\]
Then $V=\bigoplus_\phi V(\phi)$, the sum over all systems of eigenvalues, and each $V(\phi)\ne0$. The projection onto $V(\phi)$ is an idempotent $e_\phi\in\T$, the
idempotent of the local factor of $\T$ at $\ker\phi$ (\Cref{thm:modules:artinian-algebras}). Each $V(\phi)$ contains a common eigenvector $v\ne0$, with $tv=\phi(t)v$ for
all $t\in\T$. Every $\T$-stable subspace $W$ satisfies $W=\bigoplus_\phi(W\cap V(\phi))$.
\item (Integral version.) Let $\mathcal O$ be a complete discrete valuation ring, $L$ a lattice, and $\T\subseteq\End_{\mathcal O}(L)$ a commutative subalgebra, such that the
characteristic polynomials of the elements of $\T$ on $V=L\otimes K$ split over $K$. Then every system of eigenvalues $\phi\colon\T\to K$ takes values in $\mathcal O$;
$\mf m_\phi=\phi^{-1}(\varpi\mathcal O)$ is a maximal ideal of $\T$; $\mf m_\phi=\mf m_\psi$ iff $\phi\equiv\psi\pmod\varpi$ on $\T$; and for each maximal ideal $\mf m$,
\[
e_{\mf m}L\otimes K=\bigoplus_{\phi:\ \mf m_\phi=\mf m}V(\phi).
\]
(If the eigenvalues are not in $K$, replace $\mathcal O$ by the ring of integers of a finite extension, \Cref{ch:local}.)
\end{enumerate}
\end{theorem}

\begin{proof}
(i) $\T$ is a finite-dimensional commutative $F$-algebra, so $\T=\prod_i\T_i$ with local factors $\T_i=e_i\T$ (\Cref{thm:modules:artinian-algebras}). The residue
field $\T/\mf m_i$ is generated over $F$ by the images of the $t\in\T$. Each such image is a root of the minimal polynomial of $t$ as an element of $\T$, which is the
minimal polynomial of $t$ as an operator ($\T\subseteq\End V$), a divisor of the characteristic polynomial, hence a product of linear factors over $F$; so the
image lies in $F$, and $\T/\mf m_i=F$. This gives a system of eigenvalues $\phi_i\colon\T\to\T/\mf m_i=F$, and conversely every $\phi$ has a maximal kernel; so the
$\phi$ correspond to the $\mf m_i$. Put $V_i=e_iV$; then $V=\bigoplus V_i$, and $V_i\ne0$ because $e_i\ne0$ in $\T\subseteq\End V$. On $V_i$, $t-\phi_i(t)$ acts through
the nilpotent maximal ideal of $\T_i$, so it is a nilpotent operator and $(t-\phi_i(t))^{\dim V}=0$ there: $V_i\subseteq V(\phi_i)$. Conversely, if $v\in V(\phi_i)$ and
$j\ne i$, take $t=e_j$: $\phi_i(e_j)=0$, so $e_j^{\dim V}v=e_jv=0$; hence $v=e_iv\in V_i$. So $V(\phi_i)=V_i$ and $e_{\phi_i}=e_i$. A common eigenvector: let $\mf n$ be the
maximal ideal of $\T_i$ and $c\ge0$ the largest integer with $\mf n^cV_i\ne0$ (with $\mf n^0=\T_i$; it exists as $V_i\ne0$ and $\mf n$ is nilpotent); then
$\mf n^{c+1}V_i=0$, so a nonzero $v\in\mf n^cV_i$ is killed by $\mf n$, and
$tv=\phi_i(t)v$. Finally $e_iW\subseteq W\cap V_i$ for $W$ stable, and $w=\sum e_iw$.
(ii) $\T$ is a finitely generated $\mathcal O$-submodule of $\End_{\mathcal O}(L)\simeq\mathcal O^{n^2}$, hence finite free, and complete $\mathcal O$ makes
\Cref{thm:modules:idempotent-lifting} available. For $\phi$ a system of eigenvalues on $V$, $\phi(t)$ is an eigenvalue of $t$ (by (i) over $K$), hence a root of the
characteristic polynomial of $t$ on $L$, which is monic with coefficients in $\mathcal O$; a root $x\in K$ of such a polynomial has $v(x)\ge0$ (if $v(x)<0$ the
leading term has strictly smaller valuation than all others). So $\phi(\T)\subseteq\mathcal O$, and $\phi$ is an $\mathcal O$-algebra map, so $\T\to\mathcal O\to k$ is onto with kernel
$\mf m_\phi$, maximal. If $\mf m_\phi=\mf m_\psi=\mf m$, the reductions $\bar\phi,\bar\psi\colon\T\to k$ both induce $\mathcal O$-algebra isomorphisms $\T/\mf m\to k$, which coincide since
$\T/\mf m$ is generated by the image of $\mathcal O$ ($\T/\mf m\simeq k$ via either, and $k$ is generated by $\mathcal O$); conversely congruent systems have the same kernel modulo
$\varpi$. For the last statement: $e_{\mf m}$ acts on $V(\phi)$ as an idempotent whose only eigenvalue is $\phi(e_{\mf m})\in\{0,1\}$, so as the scalar $\phi(e_{\mf m})$; and
$\phi(e_{\mf m})\equiv1$ iff $e_{\mf m}\notin\mf m_\phi$ iff $\mf m_\phi=\mf m$ (\Cref{thm:modules:idempotent-lifting}(ii)). Sum over $\phi$ using (i).
\end{proof}

\begin{paperlink}\label{pl:modules:eigenspaces}
In Lemma lem:crit, ``$\T\otimes\Q=\Q\times\T'$ with $\T'$ acting faithfully on $f^\perp$'' is \Cref{thm:modules:artinian-algebras} for $\T\otimes\Q$ together with the fact
that the local factor through which $\phi_f$ factors is $\Q$ itself, which is multiplicity one for the newform $f$ (\Cref{ch:hecke}; compare
\Cref{prop:modules:congruence-ideal}(ii)). In Theorem thm:twist, ``$g^\perp$, the $\T$-stable complement of the $g$-line (the sum of the generalised eigenspaces
of the other systems of eigenvalues)'' is $\bigoplus_{\phi\ne\phi_g}V(\phi)$ of \Cref{thm:modules:generalised-eigenspaces}(i), applied to the algebra generated by the $T_p$,
$p\nmid N_S$, over $\overline\Q$. In Step 1 of the proof of Theorem thm:quant, ``the $\mf m$-component of $S_2(\Gamma_0(N);\Zl)$ is the sum of the line of $g^S$ and of
generalised eigenspaces of systems of eigenvalues with $U_q=0$'' is part (ii).
\end{paperlink}

\begin{proposition}[projecting a congruence to a local component]\label{prop:modules:idempotent-congruence}
In \Cref{thm:modules:generalised-eigenspaces}(ii) let $\mf m$ be a maximal ideal of $\T$ and $e=e_{\mf m}$.
\begin{enumerate}[label=(\roman*)]
\item $e$ maps $L$ into $L$ and every $\T$-stable subspace $W\subseteq V$ into itself, and fixes $V(\phi)$ pointwise when $\mf m_\phi=\mf m$.
\item If $x\in V(\phi)\cap L$ with $\mf m_\phi=\mf m$, $y\in L\cap W$ and $x\equiv y\pmod{\varpi^tL}$, then $x\equiv ey\pmod{\varpi^tL}$, and $ey\in L\cap W\cap eV$.
\end{enumerate}
\end{proposition}

\begin{proof}
(i) $e\in\T$ and $W$ is $\T$-stable; $e$ acts on $V(\phi)$ by $\phi(e)=1$ (proof of \Cref{thm:modules:generalised-eigenspaces}(ii)). (ii) $ex=x$, and $e(x-y)\in e\varpi^tL\subseteq\varpi^tL$.
\end{proof}

\begin{paperlink}\label{pl:modules:projection}
Step 1 of the proof of Theorem thm:quant: ``projecting $h$ to the $\mf m$-component (an idempotent of $\T_\ell$, so the projection stays in the complement and the
congruence survives) we may take $h$ in the sum of the generalised eigenspaces above''. This is \Cref{prop:modules:idempotent-congruence}(ii) with $x=g^S$, $W$ the
complement of the line of $g^S$, and $y=h$.
\end{paperlink}

\begin{proposition}[semisimple and non-semisimple operators]\label{prop:modules:semisimple-operators}
Let $A\in\End_F(V)$ with $V$ finite-dimensional over a field $F$, and let $\mu_A$ be its minimal polynomial.
\begin{enumerate}[label=(\roman*)]
\item $A$ is \emph{semisimple} if every $A$-stable subspace has an $A$-stable complement. This holds iff $\mu_A$ is squarefree, iff $F[A]$ is reduced; for $F$ perfect
(of characteristic $0$, or finite) it holds iff $A$ is diagonalisable over an algebraic closure $\overline F$.
\item If $A$ has a cyclic vector ($v$ with $v,Av,\dots,A^{n-1}v$ a basis), then $\mu_A$ equals the characteristic polynomial; so $A$ is not semisimple as soon as
the characteristic polynomial has a repeated root.
\item A commutative subalgebra $\T\subseteq\End_F(V)$ is reduced iff each of its elements is semisimple. If some $t\in\T$ is not semisimple, $\T$ contains a nonzero
nilpotent element.
\item (Jordan-Chevalley.) For $F$ perfect, $A=A_s+A_n$ with $A_s$ semisimple, $A_n$ nilpotent, $A_sA_n=A_nA_s$, and $A_s$, $A_n$ polynomials in $A$; the decomposition is
unique.
\end{enumerate}
\end{proposition}

\begin{proof}
(i) $F[A]\simeq F[X]/(\mu_A)$. If $\mu_A=g^2h$ with $g$ nonconstant then $gh(A)\ne0$ and $(gh)^2(A)=0$; and if $\mu_A=\prod g_i$ with distinct irreducible $g_i$ then
$F[X]/(\mu_A)\simeq\prod F[X]/(g_i)$ is a product of fields (Chinese remainder theorem). So $\mu_A$ squarefree iff $F[A]$ reduced. If $F[A]=\prod F_i$ is a product of fields
with idempotents $e_i$, an $A$-stable $W$ is an $F[A]$-submodule, $W=\bigoplus e_iW$ with $e_iW$ an $F_i$-subspace of the $F_i$-vector space $e_iV$; choose $F_i$-complements
$C_i$; then $\bigoplus C_i$ is an $A$-stable complement. If $\mu_A$ is not squarefree, let $N=gh(A)$ as above, $N\ne0=N^2$, and $W=\ker N\supseteq\im N\ne0$. If $V=W\oplus C$ with $C$
stable, then $N(C)\subseteq C\cap\im N\subseteq C\cap W=0$, so $C\subseteq\ker N=W$, $C=0$ and $N=0$: a contradiction. For perfect $F$, an irreducible polynomial has distinct
roots in $\overline F$; so $\mu_A$ is squarefree over $F$ iff it has distinct roots in $\overline F$, iff $A$ is diagonalisable over $\overline F$. (The minimal polynomial does
not change under field extension, being determined by the linear dependences among $1,A,A^2,\dots$.)
(ii) No nonzero polynomial of degree $<n$ kills $v$, so $\deg\mu_A\ge n$, and $\mu_A$ divides the characteristic polynomial (Cayley-Hamilton). A repeated root makes
$\mu_A$ non-squarefree.
(iii) If $\T$ is reduced so is each $F[t]\subseteq\T$; conversely a nonzero nilpotent $t$ has $\mu_t=X^k$ with $k\ge2$. If $t$ is not semisimple, $F[t]$ contains the nonzero
nilpotent $gh(t)$.
(iv) We do not prove this; see \cite[Section 12.3]{DummitFoote} (section number not verified) or any text on linear algebra.
\end{proof}

\begin{example}[the shift operator: $U_p$ on oldforms]\label{ex:modules:shift}
Let $V$ have basis $v_0,\dots,v_e$ ($e\ge1$) over a field $F$ and let $Uv_j=v_{j-1}$ for $j\ge1$.
\begin{enumerate}[label=(\alph*)]
\item If $Uv_0=av_0-pv_1$ with $p\ne0$, the characteristic polynomial of $U$ is $X^{e-1}(X^2-aX+p)$, and $v_e$ is a cyclic vector ($U^jv_e=v_{e-j}$). So for $e\ge3$ the
eigenvalue $0$ has a single Jordan block, of size $e-1\ge2$, and $U$ is not semisimple. Nevertheless every eigenvalue has an eigenvector: for $e\ge2$,
$U(v_0-av_1+pv_2)=(av_0-pv_1)-av_0+pv_1=0$; and if $\alpha,\beta$ are the roots of $X^2-aX+p$, $U(v_0-\beta v_1)=(a-\beta)v_0-pv_1=\alpha(v_0-\beta v_1)$.
\item If $Uv_0=av_0$ with $a\ne0$, the characteristic polynomial is $X^e(X-a)$, $v_e$ is again cyclic, $U$ is not semisimple for $e\ge2$, and $U(v_0-av_1)=0$.
\end{enumerate}
The characteristic polynomials (for $e\le4$, with $a$, $p$ as indeterminates) and the cyclicity of $v_e$ were checked with PARI. With $v_j=h(p^jz)$ for a newform $h$
of level $M$, (a) is $U_p$ on the span of the $h(p^jz)$ when $p\nmid M$, and (b) when $p\,\|\,M$ (\Cref{ch:hecke}); for instance at level $88$ with $h=f_{11}$, $p=2$,
$a=a_2(f_{11})=-2$ and $e=3$, the characteristic polynomial is $X^2(X^2+2X+2)$ and the algebra generated by $U_2$ on this span is not reduced
(\Cref{prop:modules:semisimple-operators}(iii)).
\end{example}

\begin{paperlink}\label{pl:modules:old}
Lemma lem:old: ``The operator $U_p$ is not semisimple on this span when $0$ is a repeated eigenvalue, but each listed eigenvalue has an eigenvector'' is
\Cref{ex:modules:shift}, with the eigenvectors the $p$-stabilisations $h-\beta_ph(pz)$, $h-\alpha_ph(pz)$ and the depleted forms. The same example shows
that the full Hecke algebra need not be reduced, which matters in Lemma lem:ddt417 (\Cref{prop:modules:congruence-ideal}(iv)).
\end{paperlink}

\begin{lemma}[common eigenvectors]\label{lem:modules:common-eigenvector}
Let $F$ be algebraically closed and $V\ne0$ finite-dimensional over $F$.
\begin{enumerate}[label=(\roman*)]
\item Let $A_1,\dots,A_k$ be commuting operators on $V$ and $\lambda=(\lambda_1,\dots,\lambda_k)\in F^k$, with joint generalised eigenspace
$V(\lambda)=\bigcap_i\ker(A_i-\lambda_i)^{\dim V}$. If $V(\lambda)\ne0$, it contains a common eigenvector of the $A_i$, which can be chosen to be an eigenvector of every
element of any commutative algebra $\T\subseteq\End_F(V)$ containing the $A_i$.
\item If $V=V_1\otimes\cdots\otimes V_k$ and $A_i$ acts on the $i$-th factor (as $1\otimes\cdots\otimes A_i\otimes\cdots\otimes1$), then $V(\lambda)=V_1(\lambda_1)\otimes\cdots\otimes V_k(\lambda_k)$, which is nonzero
iff each $\lambda_i$ is an eigenvalue of $A_i$ on $V_i$.
\end{enumerate}
\end{lemma}

\begin{proof}
(i) $V(\lambda)$ is stable under $\T$, since $\T$ commutes with the $A_i$. Apply \Cref{thm:modules:generalised-eigenspaces}(i) to the image of $\T$ in $\End V(\lambda)$
(characteristic polynomials split as $F$ is algebraically closed): $V(\lambda)$ contains a common eigenvector $v$ of $\T$, and $A_iv=\lambda_iv$ because $A_i-\lambda_i$ is
nilpotent on $V(\lambda)$.
(ii) Decompose each $V_i=\bigoplus_\mu V_i(\mu)$ into generalised eigenspaces of $A_i$; then $V$ is the direct sum of the $V_1(\mu_1)\otimes\cdots\otimes V_k(\mu_k)$. On such a summand
$A_i-\mu_i$ is nilpotent, and if $\mu_i\ne\lambda_i$ then $A_i-\lambda_i$ is invertible on it. So the summands in $V(\lambda)$ are exactly those with $\mu=\lambda$.
\end{proof}

\begin{paperlink}\label{pl:modules:common}
Lemma lem:old: ``Since $V$ is the tensor product of these local spaces and $\T$ is commutative, for every choice of one such eigenvalue $\mu_p$ at each $p\mid N/M$
there is a common eigenform $g\in V\otimes\overline\Q$ of all of $\T$''. Here $V=\bigotimes_{p\mid N/M}\operatorname{span}\{h(p^jz):0\le j\le e_p\}$, with $U_p$ acting on the $p$-th
factor, and \Cref{lem:modules:common-eigenvector}(ii) then (i) give $g$.
\end{paperlink}

\section[Congruence ideals and Gorenstein algebras]{Congruence ideals, Gorenstein algebras, and the ingredients of Lemma lem:ddt417}\label{sec:modules:congruence}

This section assembles the previous ones into the commutative algebra of congruences. A homomorphism $\pi\colon\T\to\mathcal O$ (the eigencharacter of a
form) has a \emph{congruence ideal} $\eta=\pi(\Ann_\T\ker\pi)$, which measures how far the factor of $\pi$ is congruent to the rest of $\T$; it is the paper's
$r$ in Lemma lem:crit and the $\eta$ of Lemma lem:ddt417. We prove every module-theoretic step of the proof of Lemma lem:ddt417 and stop short of the
lemma itself, which \Cref{ch:congruence} proves; we then define Gorenstein algebras and prove the freeness criterion behind Remark rem:free. Throughout,
$\mathcal O$ is a discrete valuation ring with fraction field $K$, and $\T$ is a finite free $\mathcal O$-algebra.

\begin{definition}[congruence ideal and congruence module]\label{def:modules:congruence-ideal}
Let $\pi\colon\T\to\mathcal O$ be a surjective $\mathcal O$-algebra homomorphism, $\wp=\ker\pi$ and $I=\Ann_\T(\wp)$. The \emph{congruence ideal} of $\pi$ is $\eta=\pi(I)\subseteq\mathcal O$,
and the \emph{congruence module} is $\T/(\wp+I)$, which is isomorphic to $\mathcal O/\eta$ by \Cref{thm:modules:isomorphism}(iv). Let $\pi_K\colon\T_K\to K$ be the extension
of $\pi$ and $\T_K=A_\pi\times\T'_K$ the decomposition in which $A_\pi$ is the local factor through which $\pi_K$ factors (\Cref{thm:modules:finite-algebra-dvr}(iii)),
with idempotent $e_\pi$.
\end{definition}

\begin{proposition}[congruence ideals]\label{prop:modules:congruence-ideal}
In the setting of \Cref{def:modules:congruence-ideal}:
\begin{enumerate}[label=(\roman*)]
\item $I\ne0$; indeed $I=\T\cap\Ann_{\T_K}(\wp_K)$, where $\wp_K=\wp\otimes K=\ker\pi_K$.
\item The following are equivalent: (a) $\eta\ne0$; (b) $\wp\cap I=0$; (c) $A_\pi=K$, that is, $\T_K=K\times\T'_K$ with $\pi_K$ the first projection.
\item If $\eta\ne0$: $\pi|_I\colon I\to\eta$ is an isomorphism, so $I$ is free of rank one over $\mathcal O$, generated by any $t\in I$ with $\pi(t)\mathcal O=\eta$; $\Ann_\T(I)=\wp$;
$I=\T\cap(K\times0)$ and $\wp=\T\cap(0\times\T'_K)$; and $\eta=\{a\in\mathcal O:ae_\pi\in\T\}$, the ``denominator'' of the idempotent $e_\pi$.
\item If $\T_K$ is reduced, (c) holds automatically. $\T_K$ need not be reduced (\Cref{ex:modules:shift}); then (c) must be checked directly. For the
eigencharacter of a newform it is multiplicity one (\Cref{ch:hecke}).
\item (Compatibility.) If $\mathcal O$ is complete and $\mf m=\mf m_\pi$, write $\T=\T_{\mf m}\times\T''$ with $\T''=\prod_{\mf m'\ne\mf m}\T_{\mf m'}$ (\Cref{thm:modules:idempotent-lifting}). Then
$\wp=\wp_{\mf m}\times\T''$ with $\wp_{\mf m}=\ker(\pi|_{\T_{\mf m}})$, and $\Ann_\T(\wp)=\Ann_{\T_{\mf m}}(\wp_{\mf m})\times0$; so $\eta(\T,\pi)=\eta(\T_{\mf m},\pi|_{\T_{\mf m}})$: \emph{the congruence ideal is
computed in $\T_{\mf m}$}. The other local factors lie in $\wp$, and their component of the annihilator is $0$. For a flat extension $\mathcal O\to\mathcal O'$ of discrete valuation
rings (for example $\Z_{(\ell)}\to\Zl$), $\eta(\T\otimes\mathcal O',\pi\otimes\mathcal O')=\eta\mathcal O'$.
\end{enumerate}
\end{proposition}

\begin{proof}
(i) For $t\in\T$, $t\wp=0$ iff $t\wp_K=0$, as $\wp_K=K\wp$. Write $\T_K=A_\pi\times\T'_K$; then $\wp_K=\mf m_{A_\pi}\times\T'_K$ and
$\Ann_{\T_K}(\wp_K)=\Ann_{A_\pi}(\mf m_{A_\pi})\times0$. In the local Artinian ring $A_\pi$ the maximal ideal is nilpotent (\Cref{thm:modules:artinian-algebras}); if
$\mf m_{A_\pi}^{c+1}=0\ne\mf m_{A_\pi}^c$ (with $\mf m^0=A_\pi$), then $0\ne\mf m_{A_\pi}^c\subseteq\Ann(\mf m_{A_\pi})$. Every nonzero $y\in\T_K$ has a nonzero multiple $\varpi^sy\in\T$; so
$I\ne0$.
(ii) (a)$\Rightarrow$(b): choose $t\in I$ with $\pi(t)$ a generator of $\eta$. For $x\in\wp\cap I$ write $\pi(t)x=tx+(\pi(t)-t)x$. Here $tx\in I\wp=0$, and $(\pi(t)-t)x\in\wp I=0$
because $\pi(t)-t\in\wp$. So $\pi(t)x=0$, and $x=0$ because $\T$ is torsion-free and $\pi(t)\ne0$.
(b)$\Rightarrow$(a): if $\eta=0$ then $I\subseteq\wp$, so $\wp\cap I=I\ne0$ by (i).
(a)$\Leftrightarrow$(c): if $A_\pi$ is not a field, then $\mf m_{A_\pi}\ne0$ and $\Ann_{A_\pi}(\mf m_{A_\pi})\subseteq\mf m_{A_\pi}$ (a unit cannot kill a nonzero ideal), so
$I\subseteq\T\cap(\mf m_{A_\pi}\times0)\subseteq\wp_K$ and $\eta=\pi(I)=0$. If $A_\pi$ is a field, it is $K$ ($\pi_K$ maps it onto $K$ as a $K$-algebra). Then $e_\pi=(1,0)$ has a
nonzero multiple $ae_\pi\in\T$ with $a\in\mathcal O$; it kills $\wp_K=0\times\T'_K$, so lies in $I$, and $\pi(ae_\pi)=a\ne0$.
(iii) The kernel of $\pi|_I$ is $\wp\cap I=0$, and its image is $\eta$ by definition; $\eta\simeq\mathcal O$. If $xI=0$ then $\pi(x)\pi(t)=\pi(xt)=0$, so $\pi(x)=0$; conversely
$\wp I=0$. By (c), $\wp_K=0\times\T'_K$ and $\Ann_{\T_K}(\wp_K)=K\times0$, so (i) gives $I=\T\cap(K\times0)$, and $\wp=\T\cap\ker\pi_K=\T\cap(0\times\T'_K)$. An element of
$\T\cap(K\times0)$ is $ce_\pi$ with $c=\pi_K(ce_\pi)$; so $\eta=\{c:ce_\pi\in\T\}$ (such $c$ lie in $\mathcal O$ automatically).
(iv) If $\T_K$ is reduced its local factors are fields (\Cref{thm:modules:artinian-algebras}(iii)), and $A_\pi=K$ as in (ii).
(v) $\pi$ kills $\T''$ (\Cref{thm:modules:idempotent-lifting}(iv)), so $\wp=\wp_{\mf m}\times\T''$. An element $(a,b)$ kills $\wp$ iff $a\wp_{\mf m}=0$ and $b\T''=0$, that is, $b=0$; so
$\Ann_\T(\wp)=\Ann_{\T_{\mf m}}(\wp_{\mf m})\times0$, and $\pi(\Ann_\T\wp)=\pi(\Ann_{\T_{\mf m}}\wp_{\mf m})$. For a flat $\mathcal O\to\mathcal O'$, $\ker(\pi\otimes\mathcal O')=\wp\otimes\mathcal O'$ and
$\Ann(\wp\otimes\mathcal O')=I\otimes\mathcal O'$ (\Cref{prop:modules:flat}(iii); $\wp$ is finitely generated), so $\eta(\T\otimes\mathcal O')=\pi(I)\mathcal O'$.
\end{proof}

\begin{paperlink}\label{pl:modules:congruence-ideal}
Lemma lem:crit: ``$I=\T\cap(0\times\T')$ and $J=\T\cap(\Q\times0)$, so $I\cap J=0$ and $\T/(I+J)\simeq\Z/\phi_f(J)$'' is \Cref{prop:modules:congruence-ideal}(iii) over $\mathcal O=\Z_{(\ell)}$
(or over $\Z$ verbatim, the proofs using only that $\Z$ is a principal ideal domain), with $\pi=\phi_f$, $\wp=\ker\phi_f$ (the paper's $I$ there) and $\Ann(\wp)$ (the paper's
$J$); Appendix B, convention B-C13(c), explains the clash of letters between Lemmas lem:crit and lem:ddt417. In Lemma lem:ddt417 the paper now proves $\wp\cap I=0$ by
the argument of (ii), and remarks that reducedness would not do. The paragraph after Lemma lem:hecke, and Step 1 of the proof of Theorem thm:quant, say that
``the congruence ideal is computed in $\T_{\mf m}$: the other local factors of $\T\otimes\Zl$ lie in the kernel $\wp$, $\phi_f$ factors through $\T_{\mf m}$, and $\Ann(\wp)$
decomposes along the idempotents'': this is (v).
\end{paperlink}

\begin{example}[toy congruences]\label{ex:modules:toy-congruence}
All entries were computed by hand from the definitions (and, for $\varpi=3$, $r=1,2$, checked by computer).
\begin{center}\small
\begin{tabular}{lllll}
\toprule
$\T$, $\pi$ & $\wp$ & $I=\Ann_\T(\wp)$ & $\eta$ & Gorenstein?\\
\midrule
(a) $\T_r$, $\pi(a,b)=a$ & $0\times\varpi^r\mathcal O$ & $\varpi^r\mathcal O\times0$ & $\varpi^r\mathcal O$ & yes\\
(b) $\Z[x]/(x^2-4)$, $\pi(x)=2$ & $\Z(x-2)$ & $\Z(x+2)$ & $4\Z$ & yes\\
(c) $\T^\flat_3$, $\pi(a,b,c)=a$ & $0\times\varpi\mathcal O\times\varpi\mathcal O$ & $\varpi\mathcal O\times0\times0$ & $\varpi\mathcal O$ & no\\
(d) $\mathcal O[x]/(x^2)$, $\pi(x)=0$ & $\mathcal Ox$ & $\mathcal Ox$ & $0$ & yes\\
(e) $\Z[x]/(x^2(x-1))$, $\pi(x)=0$ & $(x)$ & $\Z(x-x^2)$ & $0$ & yes\\
(e$'$) $\Z[x]/(x^2(x-1))$, $\pi(x)=1$ & $(x-1)$ & $\Z x^2$ & $\Z$ & yes\\
\bottomrule
\end{tabular}
\end{center}
Here $\T_r=\{(a,b)\in\mathcal O^2:a\equiv b\bmod\varpi^r\}$ and $\T^\flat_3=\{(a,b,c)\in\mathcal O^3:a\equiv b\equiv c\bmod\varpi\}$.
(a) $\T_r\simeq\mathcal O[x]/(x(x-\varpi^r))$ with $x=(\varpi^r,0)$; $e_\pi=(1,0)=x/\varpi^r$, with denominator $\varpi^r$; the congruence module is $\mathcal O/\varpi^r$: ``two eigenforms
congruent modulo $\varpi^r$''. (b) $(a+bx)(x-2)=(a-2b)x+(4b-2a)$ vanishes iff $a=2b$, so $I=\Z(x+2)$, $\pi(x+2)=4$, $e_\pi=(x+2)/4$, congruence module $\Z/4$; under
$\T\simeq\{(a,b):a\equiv b\bmod4\}$ this is (a) with $\varpi^r=4$. (c) An element $(a,b,c)$ kills $\wp$ iff $b=c=0$; three pairwise congruent eigenforms, congruence ideal $\varpi\mathcal O$.
(d) Here $\wp\cap I=\mathcal Ox\ne0$ and $\eta=0$: the factor of $\pi$ in $\T_K=K[x]/(x^2)$ is not a field. (e) For $\pi(x)=0$ the factor of $\T_\Q$ is $\Q[x]/(x^2)$; the
annihilator of $x$ consists of the $c(x^2-x)$ (use $x^3=x^2$), and $\pi(x-x^2)=0$. For $\pi(x)=1$ the factor is $\Q$; $I=\Z x^2$ with $x^2$ an idempotent of $\T$, so
$\eta=\Z$ and the congruence module is $0$: no congruence. The Gorenstein column is \Cref{prop:modules:gorenstein}(iii).
\end{example}

\begin{proposition}[the module ingredients of Lemma lem:ddt417]\label{prop:modules:ddt-ingredients}
In the setting of \Cref{def:modules:congruence-ideal} with $\eta\ne0$, choose $t\in I$ with $\pi(t)\mathcal O=\eta$. Let $L$ be a $\T$-module that is free of finite rank over
$\mathcal O$, with a perfect $\mathcal O$-bilinear pairing for which every element of $\T$ is self-adjoint. Then:
\begin{enumerate}[label=(\roman*)]
\item $L[\wp]$ and $L[I]$ are saturated in $L$.
\item $L[\wp]\cap L[I]=0$, so the sum $L[\wp]+L[I]$ is direct.
\item $L[\wp]^\perp=L[I]={}^\perp L[\wp]$ and $L[I]^\perp=L[\wp]={}^\perp L[I]$.
\item More generally $(\mf aL)^\perp=L[\mf a]$ for every ideal $\mf a\subseteq\T$.
\item For every $x\in L$, $\pi(t)x=tx+(\pi(t)-t)x$ with $tx\in L[\wp]$ and $(\pi(t)-t)x\in L[I]$; so $Q=L/(L[\wp]\oplus L[I])$ is killed by $\eta$.
\item If $L\simeq\T^r$ as $\T$-modules, then $L[\wp]=I^r$ (free of rank $r$ over $\mathcal O$), $L[I]=\wp^r$, and $Q\simeq(\T/(\wp+I))^r\simeq(\mathcal O/\eta)^r$.
\end{enumerate}
\end{proposition}

\begin{proof}
(i) \Cref{prop:modules:saturated}(iii). (ii) $\pi(t)=t+(\pi(t)-t)\in I+\wp$, so an element killed by $\wp$ and by $I$ is killed by $\pi(t)\ne0$, hence is $0$ ($L$ is torsion-free
over $\mathcal O$). Directness is \Cref{prop:modules:two-submodules}.
(iii) Let $x\in L[\wp]$ and $y\in L[I]$. Since $(\pi(t)-t)x\in\wp x=0$, $\pi(t)\langle x,y\rangle=\langle tx,y\rangle=\langle x,ty\rangle=0$, and $\langle x,y\rangle=0$; likewise $\langle y,x\rangle=0$. So
$L[I]\subseteq L[\wp]^\perp\cap{}^\perp L[\wp]$. Conversely let $y\in L[\wp]^\perp$. For $s\in I$ and $z\in L$, $sz\in L[\wp]$ (as $\wp s=0$), so $\langle z,sy\rangle=\langle sz,y\rangle=0$;
by perfectness $sy=0$, and $y\in L[I]$. The same argument with the arguments exchanged gives ${}^\perp L[\wp]=L[I]$. Then $L[I]^\perp=({}^\perp L[\wp])^\perp=L[\wp]$
and ${}^\perp L[I]={}^\perp(L[\wp]^\perp)=L[\wp]$ by \Cref{prop:modules:orthogonal}(ii), $L[\wp]$ being saturated.
(iv) is \Cref{prop:modules:adjoint}(ii).
(v) $\wp tx\subseteq\wp Ix=0$; and for $s\in I$, $s(\pi(t)-t)\in I\wp=0$.
(vi) Componentwise, $\T^r[\mf a]=\Ann_\T(\mf a)^r$ (\Cref{prop:modules:annihilator}(iii)); $\Ann_\T(\wp)=I$ and $\Ann_\T(I)=\wp$ (\Cref{prop:modules:congruence-ideal}(iii)). So
$Q=\T^r/(I^r\oplus\wp^r)=(\T/(I+\wp))^r$, and $\T/(\wp+I)\simeq\mathcal O/\eta$.
\end{proof}

\begin{definition}[Gorenstein algebras]\label{def:modules:gorenstein}
Let $\T^\vee=\Hom_{\mathcal O}(\T,\mathcal O)$, a $\T$-module by $(t\varphi)(u)=\varphi(tu)$. The finite free $\mathcal O$-algebra $\T$ is \emph{Gorenstein} if $\T^\vee$ is free of rank one
as a $\T$-module.
\end{definition}

\begin{proposition}[Gorenstein criteria]\label{prop:modules:gorenstein}
\begin{enumerate}[label=(\roman*)]
\item $\T$ is Gorenstein iff there is a perfect $\mathcal O$-bilinear pairing $\T\times\T\to\mathcal O$ with $\langle st,u\rangle=\langle s,tu\rangle$; one may take $\langle s,t\rangle=\lambda(st)$ for a
generator $\lambda$ of $\T^\vee$.
\item If $\mathcal O$ is complete, $\T$ is Gorenstein iff each local factor $\T_{\mf m}$ is. If $\T$ is local, with maximal ideal $\mf m$ and residue field $k_\T$, the following
are equivalent: $\T$ is Gorenstein; $\dim_{k_\T}(\T/\varpi\T)[\mf m]=1$ (the \emph{socle} of $\T/\varpi\T$ is a line); $\dim_{k_\T}\T^\vee/\mf m\T^\vee=1$.
\item Examples: $\mathcal O[x]/(f)$ for monic $f$ is Gorenstein; $\T_r$ of \Cref{ex:modules:toy-congruence}(a) is Gorenstein; $\T^\flat_3$ of (c) is not, although
$\T^\flat_3\otimes K=K^3$ is reduced; $\mathcal O[x]/(x^2)$ is Gorenstein and not reduced.
\item Complete intersections (quotients $\mathcal O[[x_1,\dots,x_n]]/(f_1,\dots,f_n)$ that are finite free over $\mathcal O$) are Gorenstein; and for $\T$ local Gorenstein
and $\pi$ as in \Cref{def:modules:congruence-ideal}, Wiles's numerical criterion compares $\eta$ with $\wp/\wp^2$ (stated; see \Cref{ch:congruence}).
\end{enumerate}
\end{proposition}

\begin{proof}
(i) If $\lambda$ generates $\T^\vee$ freely, $s\mapsto s\lambda=\lambda(s\,\cdot)$ is an isomorphism $\T\to\T^\vee$, which says that $\lambda(st)$ is perfect. Conversely, given the
pairing, $s\mapsto\langle s,-\rangle$ is a $\T$-linear isomorphism $\T\to\T^\vee$ (\Cref{prop:modules:algebra-duality}(i) with $S=\T$).
(ii) $\T^\vee=\prod_{\mf m}\T_{\mf m}^\vee$ compatibly with $\T=\prod\T_{\mf m}$, and a module over a product is free of rank one iff each component is. Let $\T$ be local and
$A=\T/\varpi\T$, a local $k$-algebra with maximal ideal $\bar{\mf m}=\mf m/\varpi\T$ and $A/\bar{\mf m}=k_\T$. As $\T$ is free, $\T^\vee/\varpi\T^\vee=\Hom_k(A,k)=:A'$. Under the
perfect pairing $A\times A'\to k$, $(x,\varphi)\mapsto\varphi(x)$, the orthogonal of the subspace $\bar{\mf m}A'$ is $\{x:\varphi(ax)=0\ \forall a\in\bar{\mf m},\varphi\in A'\}=A[\bar{\mf m}]$. So
$\dim_k\T^\vee/\mf m\T^\vee=\dim_kA'/\bar{\mf m}A'=\dim_kA[\bar{\mf m}]$, and dividing by $[k_\T:k]$ gives the equality of $k_\T$-dimensions. By Nakayama
(\Cref{thm:modules:nakayama}(iii)), $\dim_{k_\T}\T^\vee/\mf m\T^\vee=1$ iff $\T^\vee$ is cyclic, $\T\onto\T^\vee$. Both are free $\mathcal O$-modules of the same rank, so this surjection is an
isomorphism (compose with an $\mathcal O$-isomorphism $\T^\vee\simeq\T$ and use \Cref{lem:modules:determinant-trick}(b)), and $\T^\vee$ is free of rank one.
(iii) For $f$ monic of degree $n$, with basis $1,x,\dots,x^{n-1}$, let $\lambda$ be the coefficient of $x^{n-1}$. Then $\lambda(x^ix^j)=0$ for $i+j<n-1$ and $=1$ for $i+j=n-1$,
so the Gram matrix of $\lambda(st)$ is triangular with respect to the antidiagonal, with $1$'s on it, and has determinant $\pm1$; apply (i). This covers $\mathcal O[x]/(x^2)$
and $\T_r\simeq\mathcal O[x]/(x(x-\varpi^r))$. Directly for $\T_r$: $\lambda(a,b)=(a-b)/\varpi^r$ is $\mathcal O$-valued on $\T_r$, and in the basis $(1,1)$, $(\varpi^r,0)$ the Gram matrix of
$\lambda(st)$ is $\Mtwo011{\varpi^r}$, of determinant $-1$. For $\T^\flat_3$, local with basis $(1,1,1)$, $u=(\varpi,0,0)$, $w=(0,\varpi,0)$: modulo $\varpi$, $u^2=\varpi u\equiv0$,
$uw=0$, $w^2\equiv0$, so $\T^\flat_3/\varpi=k[u,w]/(u^2,uw,w^2)$, whose socle $ku+kw$ has dimension $2$.
(iv) We do not prove this; see \cite[Chapter 21]{Eisenbud} and \cite[Section 4]{DDT}.
\end{proof}

\begin{theorem}[a freeness criterion in the style of Mazur]\label{thm:modules:mazur-freeness}
Let $\T$ be a local finite free $\mathcal O$-algebra, with maximal ideal $\mf m$ and residue field $k_\T$, and let $L\ne0$ be a $\T$-module, free of finite rank over $\mathcal O$, with a
perfect $\mathcal O$-bilinear pairing for which every element of $\T$ is self-adjoint. Suppose
\begin{enumerate}[label=(\alph*)]
\item $\dim_{k_\T}L/\mf mL=d$, and
\item $\rank_{\mathcal O}L=d\cdot\rank_{\mathcal O}\T$.
\end{enumerate}
Then $L\simeq\T^d$ as $\T$-modules, and $\T$ is Gorenstein.
\end{theorem}

\begin{proof}
By Nakayama (\Cref{thm:modules:nakayama}(iii)) there is a surjection $\T^d\onto L$. Both are free over $\mathcal O$ of rank $d\rank\T$, so it is an isomorphism (as in the proof of
\Cref{prop:modules:gorenstein}(ii)). The pairing gives a $\T$-linear isomorphism $L\simeq L^\vee$, $y\mapsto\langle-,y\rangle$ (linear because
$\langle-,ty\rangle=\langle t\,\cdot,y\rangle$), and $L^\vee\simeq(\T^d)^\vee=(\T^\vee)^d$. So $(\T^\vee)^d\simeq\T^d$, and reducing modulo $\mf m$, $(\T^\vee/\mf m\T^\vee)^d\simeq k_\T^d$, so
$\dim_{k_\T}\T^\vee/\mf m\T^\vee=1$ and $\T$ is Gorenstein by \Cref{prop:modules:gorenstein}(ii).
\end{proof}

\begin{paperlink}\label{pl:modules:free}
Remark rem:free: ``Freeness of $H^1(X_0(N),\Zl)_{\mf m}$ over $\T_{\mf m}$ holds when $J_0(N)[\mf m]$ is two-dimensional, since the localised Hecke algebra is then Gorenstein and the
Tate module is free of rank two''. In \Cref{thm:modules:mazur-freeness} take $\mathcal O=\Zl$, $\T=\T_{\mf m}$, $L=H^1(X_0(N),\Zl)_{\mf m}$ with the pairing $\hat\delta_!$ (perfect on the
$\mf m$-component by \Cref{prop:modules:adjoint}(iii)), and $d=2$. Hypothesis (a) is the two-dimensionality of $J_0(N)[\mf m]$, transported to $L/\mf mL$ by duality, and (b) is a
rank count from the Eichler-Shimura isomorphism; both are \Cref{ch:moddeg}. The conclusion $L\simeq\T_{\mf m}^2$ is the hypothesis of the equality clause of Lemma
lem:ddt417. After Proposition prop:pullback the paper no longer needs this for Theorems thm:main and thm:quant.
\end{paperlink}

\begin{remark}[the statement of Lemma lem:ddt417, and where each step is proved]\label{rem:modules:ddt417-preview}
In the book's notation the paper's Lemma lem:ddt417 reads: \emph{let $\T$ be a finite free $\mathcal O$-algebra acting on a module $L$ that is free of finite rank over $\mathcal O$,
with a perfect pairing for which $\T$ is self-adjoint; let $\pi\colon\T\onto\mathcal O$, $\wp=\ker\pi$, $I=\Ann_\T(\wp)$ and $\eta=\pi(I)$; suppose $L[\wp]$ has rank $d$ and Gram
determinant $\Delta\ne0$. Then $v(\Delta)\le d\,v(\eta)$, that is, $\eta^d\subseteq\Delta\mathcal O$, with equality if $L$ is free over $\T$.} It is proved in \Cref{ch:congruence}.
Every step except the final count is proved in this chapter:
\begin{center}\small
\begin{tabular}{ll}
\toprule
step of the paper's proof & proved in\\
\midrule
$\wp\cap I=0$ when $\eta\ne0$ & \Cref{prop:modules:congruence-ideal}(ii)\\
$L[\wp]+L[I]$ is direct & \Cref{prop:modules:ddt-ingredients}(ii)\\
$Q=L/(L[\wp]\oplus L[I])$ is killed by $\eta$ & \Cref{prop:modules:ddt-ingredients}(v)\\
$L[\wp]$ is saturated & \Cref{prop:modules:saturated}(iii), \Cref{prop:modules:ddt-ingredients}(i)\\
$L/L[\wp]^\perp\simeq\Hom(L[\wp],\mathcal O)$ & \Cref{thm:modules:gram-cokernel}(i)\\
$L[\wp]^\perp=L[I]$ & \Cref{prop:modules:ddt-ingredients}(iii)\\
$Q\simeq\coker(L[\wp]\to L[\wp]^\vee)$, of length $v(\Delta)$ & \Cref{thm:modules:gram-cokernel}(ii)\\
$Q$ is generated by $d$ elements & \Cref{thm:modules:gram-cokernel}(ii) ($L[\wp]^\vee\simeq\mathcal O^d$ surjects onto $Q$)\\
if $L\simeq\T^d$, $Q\simeq(\mathcal O/\eta)^d$ & \Cref{prop:modules:ddt-ingredients}(vi)\\
the Gram determinant in rank two is $\eta^2$ & \Cref{prop:modules:alternating}(ii)\\
\bottomrule
\end{tabular}
\end{center}
The remaining step is that a module generated by $d$ elements and killed by $\eta$ is a quotient of $(\mathcal O/\eta)^d$, of length at most $d\,v(\eta)$.
\emph{Checks on the toy algebra $\T_r$} (all computed): with $L=\T_r$ and $\langle s,u\rangle=\lambda(su)$, $\lambda(a,b)=(a-b)/\varpi^r$, one has $L[\wp]=I=\mathcal O(\varpi^r,0)$, $d=1$ and
$\Delta=\lambda((\varpi^r,0)^2)=\varpi^r$, so $v(\Delta)=r=d\,v(\eta)$ (free case, equality). With $L=\mathcal O^2$, $\T_r$ acting diagonally and the standard pairing, $L[\wp]=\mathcal Oe_1$
and $\Delta=1$, so $0=v(\Delta)<r$ (strict inequality; $L$ is not free over $\T_r$). With $L=\T_r^2$ and the alternating pairing $\langle x,y\rangle=\lambda(x_1y_2-x_2y_1)$ (perfect, with
$\T_r$ self-adjoint), $L[\wp]=I^2$ has basis $(u,0)$, $(0,u)$ with $u=(\varpi^r,0)$, $d=2$, and $\Delta=\lambda(u^2)^2=\varpi^{2r}=\eta^2$: the shape of Lemma lem:hecke and of Step 2 of the
proof of Theorem thm:quant. Finally, when $\eta=0$ the inequality is empty, and if moreover $L\simeq\T^r$, then $I\subseteq\wp$ gives $I\cdot I=0$, so the pairing vanishes on
$L[\wp]=I^r$ ($\langle\sum a_ie_i,\sum b_je_j\rangle=\sum\langle b_ja_ie_i,e_j\rangle=0$) and $\Delta=0$: the hypothesis $\Delta\ne0$ excludes this case.

\emph{The paper's proof.} The Commentary on this lemma required by decision D1 (the paper counts orders over a finite residue field, the
book states lengths over any discrete valuation ring) is given with its proof in \Cref{ch:congruence}. Two remarks concern steps proved here.
The converse inclusion ``if $y\perp L[\wp]$ then \dots\ $ty=0$'' shows only that the single element $t$ kills $y$; this suffices because $I=\mathcal Ot$
(\Cref{prop:modules:congruence-ideal}(iii)), which the paper does not say, and \Cref{prop:modules:ddt-ingredients}(iii) uses every $s\in I$ instead. In
``since $L$ is torsion-free and $\wp\cap I=0$, the sum $L[\wp]+L[I]$ is direct'' the operative reason is $\pi(t)\in\wp+I$
(\Cref{prop:modules:ddt-ingredients}(ii)); $\wp\cap I=0$ gives it only after an unstated rank argument ($I\ne0$ embeds in $\T/\wp\simeq\mathcal O$, so
$\wp+I$ has finite index in $\T$ and contains a nonzero scalar).
\end{remark}

\section{Exercises}\label{sec:modules:exercises}

The exercises are graded from routine to harder; those marked ``check the paper'' verify a sentence of the paper with the tools of this chapter.

\begin{exercise}\label{ex:modules:ex-tensor}
Compute $\Z/m\otimes\Z/n$, $\Q\otimes\Z/n$, $\Zl\otimes\Z/n$, $\Qlmod\otimes\Z/n$, $\Hom(\Z/m,\Z/n)$ and $\Hom_{\Zl}(\Qlmod,\Z/\ell^k)$. Show that $\Hom(\Z/2,-)$ is not right exact
and $-\otimes\Z/2$ is not left exact on $0\to\Z\xrightarrow2\Z\to\Z/2\to0$. (Answers: $\Z/\gcd(m,n)$, $0$, $\Z/\ell^{v_\ell(n)}$, $0$, $\Z/\gcd(m,n)$, $0$.)
\end{exercise}

\begin{exercise}[check the paper: Lemma lem:IV]\label{ex:modules:ex-IV}
Let $T$ be free of rank $2$ over $\Z_3$ with an automorphism $\sigma$ such that $\sigma^2+\sigma+1=0$. Using that $\Z_3[\zeta_3]$ is a discrete valuation ring with uniformiser
$1-\zeta_3$ (\Cref{ch:ant}), show that $T$ is free of rank one over $\Z_3[X]/(X^2+X+1)\simeq\Z_3[\zeta_3]$ (\Cref{prop:modules:torsion}(iii)), that
$T/3T\simeq\Z_3[\zeta_3]/(1-\zeta_3)^2$ (note $(1-\zeta_3)^2=-3\zeta_3$), and that the vectors of $T/3T$ fixed by $\sigma$ form a line. Deduce that in a suitable basis
$\sigma=\Mtwo0{-1}1{-1}$.
\end{exercise}

\begin{exercise}\label{ex:modules:ex-fitting}
Show that Fitting ideals are multiplicative in short exact sequences of finite $\Zl$-modules. Show that they are not over $\Z[x]$, using
$0\to R/(2,x)\to R/(4,x)\to R/(2,x)\to0$ (\Cref{prop:modules:fitting-properties}(vi)). Over $R=k[x]/(x^2)$, show that $\Fitt(k)=(x)$ and that
$0\to k\to R\to k\to0$ and $0\to k\to k\oplus k\to k\to0$ both satisfy $\Fitt(M)=\Fitt(M')\Fitt(M'')$ (both sides are $0$). Compute $\Fitt$ and $\Ann$ of $\Z[x]/(2,x)$ and of its
square.
\end{exercise}

\begin{exercise}[higher Fitting ideals]\label{ex:modules:ex-fitting-i}
For $M$ with $n$ generators and a relation matrix $A$, let $\Fitt_i(M)$ be the ideal generated by the $(n-i)\times(n-i)$ minors of $A$, with $\Fitt_i(M)=R$
for $i\ge n$ (so $\Fitt_0$ is the Fitting ideal of \Cref{def:modules:fitting}).
(a) Show that $\Fitt_i(M)$ does not depend on the presentation, and that $\Fitt_0(M)\subseteq\Fitt_1(M)\subseteq\cdots$. (Hint: the proof of
\Cref{lem:modules:fitting-well-defined} works for minors of size $n-i$; for an added generator, column operations turn the new relation matrix into
$\operatorname{diag}(A,1)$, and Laplace expansion puts its minors of size $n+1-i$ in the ideal of minors of size $n-i$ of $A$.)
(b) Over a discrete valuation ring let $M=\bigoplus_{j=1}^s\mathcal O/\varpi^{a_j}$ with $1\le a_1\le\cdots\le a_s$. Show that $\Fitt_i(M)=\varpi^{a_1+\cdots+a_{s-i}}\mathcal O$
for $i<s$ and $\Fitt_i(M)=\mathcal O$ for $i\ge s$. Deduce that the $\Fitt_i$ determine a torsion $M$ up to isomorphism, that $\dim_kM[\varpi]$ is the least
$i$ with $\Fitt_i(M)=\mathcal O$, and that $M$ is cyclic iff $\Fitt_1(M)=\mathcal O$.
(c) Over $\Zl$, compute $\Fitt_0$ and $\Fitt_1$ of $\Z/\ell^2$ and of $(\Z/\ell)^2$. (Answers: $\Fitt_0=\ell^2\Zl$ for both; $\Fitt_1=\Zl$ and $\ell\Zl$. So $\Fitt_1$
separates the two modules that \Cref{thm:modules:fitting-dvr}(iv) shows $\Fitt_0$ cannot.)
\end{exercise}

\begin{exercise}[cofinitely generated modules]\label{ex:modules:ex-cofinite}
(a) Show that a quotient of a divisible module is divisible. Let $W$ be cofinitely generated of corank $r$ and $F\subseteq W$ a finite submodule;
show that $W/F$ is cofinitely generated of corank $r$, and that $W/F\simeq(\Qlmod)^r$ if $W$ is divisible. (Hint: dualising, $(W/F)^*$ is the kernel of
$W^*\to F^*$, a submodule of finite index of $W^*$; use \Cref{thm:modules:pontryagin-zl}.)
(b) Show that multiplication by $\ell$ identifies $\Qlmod/(\Qlmod)[\ell]$ with $\Qlmod$, and that under this identification the dual of the quotient map
$\Qlmod\to\Qlmod/(\Qlmod)[\ell]$ is the inclusion $\ell\Zl\subseteq\Zl$.
(c) For $W=\Qlmod\oplus\Z/\ell$ compute $W[\ell]$, the divisible part $\bigcap_k\ell^kW$, the corank and $W^*$. (Answers: $(\Z/\ell)^2$; $\Qlmod\oplus0$; $1$;
$\Zl\oplus\Z/\ell$.)
\end{exercise}

\begin{exercise}[the Gram determinant of a sublattice of finite index; the prime $2$]\label{ex:modules:ex-gram-index}
(a) Let $L'\subseteq L$ be lattices of the same rank over $\mathcal O$, $B$ the matrix whose rows are the coordinates of a basis of $L'$ in a basis of $L$, and $G$
the Gram matrix of a pairing on $L$. Show that the Gram matrix of $L'$ is $BGB^{\mathsf T}$, so that $v(\Delta_{L'})=v(\Delta_L)+2\len(L/L')$
(\Cref{thm:modules:stacked-bases}).
(b) On $L=\Z_2^2$ with the perfect alternating pairing $\langle x,y\rangle=x_1y_2-x_2y_1$, let $L'=\Z_2e_1\oplus\Z_2\cdot2e_2$, of index $2$. Compute the Gram
determinant and the ideal of values of $L'$, and check (a) and \Cref{prop:modules:alternating}(ii). (Answer: $\Delta=4$, values $2\Z_2$.)
(c) Show that $\frac12\langle\ ,\ \rangle$ is perfect on $L'$ but not $\Z_2$-valued on $L$, while for $\ell$ odd $\frac12\langle\ ,\ \rangle$ is perfect on $\Zl^2$. Deduce
that passing between $\hat\delta_!$ and the cup product, which differ by the factor $\frac12$ (\Cref{pl:modules:adjoint}), preserves perfectness and Gram
determinants up to units only for $\ell$ odd: over $\Z_2$ the factor $\frac12$ changes the ideal of values by $2$ and
the Gram determinant of a rank-two lattice by $4$.
\end{exercise}

\begin{exercise}[the saturation of $\wp L$; reduction modulo $\varpi$]\label{ex:modules:ex-saturation}
In the setting of \Cref{prop:modules:ddt-ingredients}:
(a) Show that $\operatorname{sat}(\wp L)=L[I]$ and $\operatorname{sat}(IL)=L[\wp]$. (Hint: \Cref{prop:modules:adjoint}(ii), \Cref{prop:modules:orthogonal}(ii)
and \Cref{prop:modules:ddt-ingredients}(iii).)
(b) Let $\T_r$ ($r\ge1$, \Cref{ex:modules:toy-congruence}) act on $L=\mathcal O^2$ coordinatewise, with the standard pairing and $\pi(a,b)=a$. Show that
$\wp L=0\oplus\varpi^r\mathcal O$ is not saturated, that $L[I]=0\oplus\mathcal O$, and that the map $L[\wp]/\varpi L[\wp]\to(L/\varpi L)[\wp]$, injective for
every $L$ because $L[\wp]$ is saturated (\Cref{prop:modules:saturated}(i)(d)), is not surjective here: the reduction of $L[\wp]$ is a line, while $\wp$
kills all of $L/\varpi L=k^2$.
\end{exercise}

\begin{exercise}[check the paper: Lemma lem:crit]\label{ex:modules:ex-crit}
Let $\T$ act on $S$ with a perfect pairing $S\times\T\to\Z$ as in \Cref{prop:modules:algebra-duality}, and let $\mf a$, $\mf b$ be ideals as in its part (iii). Show that
$S/(S[\mf a]\oplus S[\mf b])$ has the same order as $\T/(\mf a+\mf b)$. Work it out for $\T=\Z[x]/(x^2-4)$ acting on $S=\T^\vee$, $\mf a=(x-2)$, $\mf b=(x+2)$ (answer:
$4$), and interpret the result as a congruence modulo $4$ between two ``eigenforms'', the elements of $S[\mf a]$ and $S[\mf b]$. Show that the rank hypothesis
in (iii) cannot be dropped. (Hint: $\mf a=\mf b=0$.)
\end{exercise}

\begin{exercise}[three eigenvalues]\label{ex:modules:ex-three}
Let $\T=\Z[x]/(x(x-2)(x-6))$, which embeds in $\Z^3$ by $x\mapsto(0,2,6)$, and $\pi(x)=0$.
(a) Show that $\wp=(x)$, $I=\Ann_\T(\wp)=\Z\,(x-2)(x-6)$ and $\eta=12\Z$, and that $12$ is the denominator of the idempotent $e_\pi=(x-2)(x-6)/12$
(\Cref{prop:modules:congruence-ideal}(iii); its proof works verbatim over $\Z$, \Cref{pl:modules:congruence-ideal}).
(b) Show that $\T\otimes\Z_2$ is local and that $\T\otimes\Z_3=\T_{\mf m}\times\T_{\mf m'}$ with $\mf m=(3,x)$ and $\mf m'=(3,x-2)$, with idempotent
$e_{\mf m'}=x(x-6)/(-8)$. Compute the congruence ideal in $\T_{\mf m}$ and in $\T\otimes\Z_2$, and compare with \Cref{prop:modules:congruence-ideal}(v).
(Answers: $3\Z_3$ and $4\Z_2$.)
(c) $2$-adically, the eigenvalue $0$ is congruent to $2$ and to $6$ only modulo $2$, yet $v_2(\eta)=2$. Explain this: $\eta$ is generated by $\pi((x-2)(x-6))=(0-2)(0-6)$, and
it measures the congruence of $\pi$ with the whole complementary factor $\Q\times\Q$ of $\T\otimes\Q$, not with each other eigenvalue separately
(\Cref{prop:modules:congruence-decomposition}).
\end{exercise}

\begin{exercise}[check the paper: the congruence ideal is computed in $\T_{\mf m}$]\label{ex:modules:ex-tm}
For $\T\otimes\Zl=\prod_{\mf m'}\T_{\mf m'}$ and $\pi$ factoring through $\T_{\mf m}$, show that $\Ann(\ker\pi)=\Ann_{\T_{\mf m}}(\wp_{\mf m})\times0$. Do the other local factors lie in the
annihilator or in the kernel? Compute both for $\T=\Z[x]/(x^2-4)$, $\ell=3$ and $\pi(x)=2$.
\end{exercise}

\begin{exercise}[check the paper: non-semisimple $U_p$ in Lemma lem:old]\label{ex:modules:ex-shift}
For the operator $U$ of \Cref{ex:modules:shift}(a) compute the characteristic polynomial, show that $v_e$ is a cyclic vector, and find an eigenvector for each eigenvalue;
deduce that for $e\ge3$ the algebra generated by $U$ is not reduced. With \Cref{ch:hecke}: carry this out for the full Hecke algebra of level $88$ acting on the span of
$f_{11}(z)$, $f_{11}(2z)$, $f_{11}(4z)$, $f_{11}(8z)$, where $a_2(f_{11})=-2$, and check with PARI that $U_2$ has characteristic polynomial $X^7(X^2+2X+2)$ on all of $S_2(\Gamma_0(88))$.
\end{exercise}

\begin{exercise}[Gorenstein algebras and faithful modules]\label{ex:modules:ex-gorenstein}
(a) Show that $\T_r$ ($r\ge1$) and $\Zl[x]/(x^2)$ are local and Gorenstein, and that $\T^\flat_3=\{(a,b,c)\in\Zl^3:a\equiv b\equiv c\bmod\ell\}$ is local and not, by
computing the socles of their reductions modulo $\ell$. For which monic $f$ is $\Zl[x]/(f)$ local? (Answer: when $f\bmod\ell$ is a power of one irreducible
polynomial.)
(b) Let $\T$ be a local finite free $\Zl$-algebra. Show that there is a \emph{faithful} $\T$-module $L$, free of finite rank over $\Zl$, with a perfect pairing
for which $\T$ is self-adjoint and with $\dim_{k_\T}L/\mf mL=1$, if and only if $\T$ is Gorenstein. (Hint: by Nakayama $L$ is cyclic, so
$L\simeq\T/\Ann_\T(L)=\T$, and \Cref{thm:modules:mazur-freeness} applies with $d=1$; conversely take $L=\T$ with the pairing of
\Cref{prop:modules:gorenstein}(i).)
(c) Show that faithfulness cannot be dropped in (b): $\T^\flat_3$ acting on $L=\Zl$ through $\pi(a,b,c)=a$, with the pairing $xy$, has $\dim L/\mf mL=1$. Show that
$\T^\flat_3$ acts faithfully on $L=\Zl^3$ coordinatewise, self-adjointly for the standard pairing, with $\dim_{\Fl}L/\mf mL=3$, so that hypothesis (b) of
\Cref{thm:modules:mazur-freeness} fails ($\rank L=3\ne3\rank\T^\flat_3$).
\end{exercise}

\section{Notes and references}\label{sec:modules:notes}

The material of Sections 1 to 4 is standard and follows Atiyah and Macdonald \cite[Chapters 1 to 3, 6, 7]{AtiyahMacdonald} and Dummit and Foote
\cite[Chapters 10 and 12]{DummitFoote}; Lang \cite[Chapters III and X]{Lang} and Eisenbud \cite[Chapters 2 to 4]{Eisenbud} give the same results in
greater generality. Our proof of right exactness of the tensor product through the adjunction is the one in Atiyah and Macdonald; the snake and five lemmas
are in Weibel \cite[Section 1.3]{Weibel}, whose numbering \Cref{ch:cohomology} uses. Hilbert's basis theorem, the equivalence of finite length with
``Noetherian and Artinian'', and the characterisation of discrete valuation rings as integrally closed one-dimensional Noetherian local domains are
cited, not proved. Every chapter, section or theorem number quoted from a book in this chapter (Atiyah-Macdonald, Dummit-Foote, Lang,
Eisenbud, Weibel, Jacobson, Matsumura, Serre, Neukirch-Schmidt-Wingberg, Rubin, Mazur, Darmon-Diamond-Taylor) is a pointer, not a verified citation.

The Smith normal form over a principal ideal domain (\Cref{thm:modules:smith}) is proved here by induction on the number of prime factors, which avoids
the Euclidean hypothesis; Jacobson \cite[Section 3.7]{Jacobson85} gives the same argument. The structure theorem and stacked bases are in Dummit and Foote
\cite[Section 12.1]{DummitFoote}. Saturated submodules are the pure submodules of Bourbaki. The congruence module of a decomposition
(\Cref{prop:modules:congruence-decomposition}) and the duality of \Cref{prop:modules:algebra-duality} are the commutative algebra behind Hida's congruence modules
\cite{Hida81} and Ribet's treatment of the congruence number \cite{Ribet83}; Agashe, Ribet and Stein \cite{ARS12} use them for the modular degree.

Discrete valuation rings and the ring $\Zl$ follow Serre \cite[Chapters I and II]{SerreLocalFields}; the analytic theory is \Cref{ch:local}. Baer's criterion and
the injectivity of divisible modules are in Weibel \cite[Section 2.3]{Weibel}. Pontryagin duality for finite groups is in Serre \cite[Chapter VI, Section 1]{SerreCourse};
for compact and discrete $\Zl$-modules see Neukirch, Schmidt and Wingberg \cite[Chapter I, Section 1]{NSW} and Rubin \cite[Appendix B]{Rubin}. Our proof of the
structure of cofinitely generated modules (\Cref{thm:modules:pontryagin-zl}(iii)) is algebraic and avoids compactness.

Fitting ideals were introduced by Fitting in 1936. We follow Northcott \cite{Northcott76} and Eisenbud \cite[Section 20.2]{Eisenbud}; the proof of
\Cref{lem:modules:fitting-well-defined} is the standard one through the ideal of determinants of all tuples of
relations. Mazur and Wiles \cite[Appendix]{MazurWiles84} use Fitting ideals of $\Zl$-modules as here, and the Tamagawa factors of the paper are Fitting ideals because
of Fontaine and Perrin-Riou \cite[I.4.2.2]{FPR94}, as recalled by Diamond, Flach and Guo \cite{DFG04}.

Sections 10 and 11 follow Darmon, Diamond and Taylor \cite[Section 4]{DDT}, especially the discussion of congruence modules leading to their Lemma 4.17, whose
inequality half is the paper's Lemma lem:ddt417; the toy algebras $\T_r$ and $\T^\flat_3$ are the standard examples of a congruence and of a non-Gorenstein algebra.
The structure of Artinian rings is \cite[Chapter 8]{AtiyahMacdonald}, and idempotent lifting over complete rings is in Eisenbud \cite[Section 7.5]{Eisenbud}; our Newton
iteration $e\mapsto3e^2-2e^3$ needs no Hensel lemma. Gorenstein rings are treated in Matsumura \cite[Chapter 18]{Matsumura86} and Eisenbud \cite[Chapter 21]{Eisenbud};
the socle criterion (\Cref{prop:modules:gorenstein}(ii)) is the form used for Hecke algebras by Mazur \cite[Chapter II, Section 15]{Mazur77} and Tilouine \cite{Tilouine97}.
\Cref{thm:modules:mazur-freeness} abstracts Mazur's argument that multiplicity one gives freeness and the Gorenstein property; the geometric input for the Hecke
algebra is \Cref{ch:moddeg}, with Ribet and Stein \cite{RibetStein}.

\emph{Rigour.} Everything in this chapter is proved except: Hilbert's basis theorem; ``finite length iff Noetherian and Artinian''; criterion (e) for discrete valuation
rings; completions in general (Artin-Rees); localisation of $\Hom$; topological Nakayama (sketched); Pontryagin duality for general compact modules; the Pfaffian
identity in general rank and the symplectic normal form (sketched); Jordan-Chevalley; complete intersections are Gorenstein, and Wiles's numerical criterion;
and the identification of $\T_{\mf m}$ with an $\mf m$-adic completion. The computations quoted (Smith forms, $11$-adic roots and idempotents, characteristic
polynomials of the shift operator and of $U_2$ at level $88$, the Pfaffian in rank $4$) were checked with PARI/GP \cite{PARI}.
